% !TEX program = pdflatex
\documentclass[12pt]{article}

\usepackage[utf8]{inputenc}
\usepackage{CJKutf8}
\usepackage{amsmath,amssymb}
\usepackage{mathtools}
\usepackage{amsthm}
\usepackage{geometry}
\geometry{
  a4paper,
  top=25mm,
  bottom=25mm,
  left=20mm,
  right=20mm
}
\usepackage{xcolor}
\usepackage{url}
\usepackage[unicode,colorlinks=true,linkcolor=blue,urlcolor=blue]{hyperref}
\usepackage{xurl}
\Urlmuskip=0mu plus 4mu
\sloppy
\usepackage{tocloft}

\renewcommand{\cftsecleader}{\cftdotfill{\cftdotsep}}
\renewcommand{\refname}{参考文献}

\newtheorem{definition}{定义}[section]
\newtheorem{axiom}{公理}[section]
\newtheorem{assumption}{假设}[section]
\newtheorem{theorem}{定理}[section]
\newtheorem{proposition}{命题}[section]
\newtheorem{corollary}{推论}[section]
\newtheorem{lemma}{引理}[section]
\newtheorem{remark}{备注}[section]
\newtheorem{Certificate}{证书}[section]

\newcommand{\NN}{\mathbb{N}}
\newcommand{\RR}{\mathbb{R}}
\newcommand{\GFramework}{\textsf{G-Framework}}
\newcommand{\DefEq}{\coloneqq}
\newcommand{\Pcal}{\mathcal P}
\newcommand{\Bcal}{\mathcal B}
\newcommand{\Ucal}{\mathcal U}
\newcommand{\Vcal}{\mathcal V}
\newcommand{\Hcal}{\mathcal H}
\newcommand{\Rcal}{\mathcal R}
\newcommand{\Gcal}{\mathcal G}
\newcommand{\TPH}{\ensuremath{\mathsf{TPH}}}
\newcommand{\LHTS}{\ensuremath{\mathsf{LHTS}}}
\newcommand{\TPHLHTS}{\ensuremath{\mathsf{TPH\mbox{-}LHTS}}}
\newcommand{\PDE}{\ensuremath{\mathsf{PDE}}}
\newcommand{\norm}[1]{\left\lVert #1 \right\rVert}
\newcommand{\ip}[2]{\left\langle #1,#2 \right\rangle}
\newcommand{\code}[1]{\path{#1}}
\DeclareMathOperator{\Sol}{Sol}
\DeclareMathOperator{\Res}{Res}
\DeclareMathOperator{\Gain}{Gain}
\DeclareMathOperator{\Audit}{Audit}
\DeclareMathOperator{\Cert}{Cert}
\DeclareMathOperator{\Corr}{Corr}
\DeclareMathOperator{\Var}{Var}
\DeclareMathOperator{\Cov}{Cov}
\DeclareMathOperator{\Eff}{Eff}
\DeclareMathOperator{\Hit}{Hit}
\DeclareMathOperator{\JacGap}{JacGap}
\DeclareMathOperator{\Eq}{Eq}
\DeclareMathOperator{\Bal}{Bal}
\DeclareMathOperator{\DD}{DD}
\DeclareMathOperator{\Exposure}{Exposure}
\DeclareMathOperator{\UnclearRisk}{UnclearRisk}
\DeclareMathOperator{\Risk}{Risk}
\DeclareMathOperator{\Cost}{Cost}
\DeclareMathOperator{\PositiveImpact}{PositiveImpact}
\DeclareMathOperator{\NegativeImpact}{NegativeImpact}
\DeclareMathOperator{\Tunnel}{Tunnel}
\DeclareMathOperator{\Subst}{Subst}
\DeclareMathOperator{\ExprSubst}{ExprSubst}
\DeclareMathOperator{\SolveSubst}{SolveSubst}
\DeclareMathOperator{\InterpEquiv}{InterpEquiv}
\DeclareMathOperator{\Solve}{Solve}
\DeclareMathOperator{\Fid}{Fid}
\DeclareMathOperator{\Obs}{Obs}
\DeclareMathOperator{\Supp}{Supp}
\DeclareMathOperator{\Sep}{Sep}
\DeclareMathOperator{\Faith}{Faith}
\DeclareMathOperator{\Cover}{Cover}
\DeclareMathOperator{\Comm}{Comm}
\DeclareMathOperator{\Maturity}{Maturity}
\DeclareMathOperator{\Meaningful}{Meaningful}
\DeclareMathOperator{\Directional}{Directional}
\DeclareMathOperator{\Composable}{Composable}
\DeclareMathOperator{\Auditable}{Auditable}
\DeclareMathOperator{\Integrable}{Integrable}
\DeclareMathOperator{\Usable}{Usable}
\DeclareMathOperator{\Marg}{Marg}
\DeclareMathOperator{\Ctrl}{Ctrl}
\DeclareMathOperator{\Desc}{Desc}
\DeclareMathOperator{\Hom}{Hom}
\DeclareMathOperator{\Interp}{Interp}
\DeclareMathOperator{\Proj}{Proj}
\DeclareMathOperator{\Dom}{Dom}
\DeclareMathOperator{\Opt}{Opt}
\newcommand{\Lmeas}{L\text{-}\mathrm{meas}}
\newcommand{\Meaning}{\mathrm{Meaning}}
\DeclareMathOperator*{\argmin}{arg\,min}
\DeclareMathOperator*{\argmax}{arg\,max}
\DeclareMathOperator*{\Argmin}{Argmin}

\setlength{\emergencystretch}{8em}
\setlength{\parindent}{0pt}
\setlength{\parskip}{0.6em}

\begin{document}
\begin{CJK*}{UTF8}{gkai}

\title{从 PDE 残差下降到 L-测度微观隧穿：\\表达替代、求解替代与解释等价\\的形式逻辑证书\\（工作笔记：形式系统版）}
\author{Gao Zheng（高政）}
\date{2026-05-18}
\maketitle

\section*{前言}
\addcontentsline{toc}{section}{前言}

本文的目标是把 PDE 求解、偏微残差下降、有限回放工程兑现、微观隧穿/同伦命中，
以及 \(L\)-积分/测度化替代统一写成一个可定义、可证明、可审计的形式系统。
纵观全文，核心结论已经不再只是“微观隧穿可以帮助 PDE 求解”，而是如下三层替代结构：
\[
  \boxed{
  \begin{gathered}
  \text{PDE 的局部微分表达}
  \Rightarrow
  L\text{-积分/测度表达替代}\\
  \Rightarrow
  L\text{-路径测度/微观隧穿求解替代}
  \Rightarrow
  \text{Law-Space 解释严格等价}.
  \end{gathered}
  }
\]
因此，本文中的“替代”不是 \(L_0\) 裸差分对 PDE 的粗糙替代，而是把 PDE 的局部偏微分表达提升为
Law-Space 中的 \(L\)-积分/测度表达；把 PDE 的函数求解对象提升为 \(L\)-路径测度空间中的低作用量命中；
并要求二者在同一解释函子 \(\mathcal C_L\) 下保持严格解释等价。

全文首先从最小 PDE 残差口径出发。设偏微系统为
\[
  \Pcal[u]=0,
  \qquad
  \Bcal[u]=0,
\]
并定义残差能量
\[
  E(u)=\frac12\norm{\Pcal[u]}^2+\frac{\lambda}{2}\norm{\Bcal[u]}^2.
\]
第一章证明求解 PDE 可写成残差能量下降与可审计增益兑现之间的正相关传导。这里的“求解”不是泛指闭式公式，
而是把 PDE 与边界/初值条件诱导的残差压低到零或工程阈值；“偏微下降”是沿函数空间残差能量泛函下降；
“增益兑现”是目标泛函中可验证、可使用、可审计的收益增加。

第二章把该残差下降口径放入有限回放证据域 \(\mathcal D_{\mathrm{replayid}}\)。通过
\[
  \Eff_{\partial\JacGap}
  \wedge
  \Hit_{\TPH}^{\partial_L\to\Phi}
  \wedge
  \Eff_{\mathrm{equity}},
\]
证明偏微正影响累计量严格大于偏微负影响累计量：
\[
  \sum_t I^+_{\partial_L}(t)
  >
  \sum_t I^-_{\partial_L}(t).
\]
该结论是有限区间上的累计工程证书，不要求每个 tick 点态为正。

第三章进一步说明，微观隧穿/同伦命中不是普通有限差分法，也不是用 \(L_0\) 裸差分替代 PDE；
它是在路径空间中寻找使残差下降、JacGap 下降、风险受控且目标兑现的有效变分路径。
第四章把该口径提升为 \(L\)-积分/测度替代：若 \(D_Lu\) 是 \(u\) 的 \(L\)-变化测度，且
\[
  \partial_Lu=\frac{dD_Lu}{d\mu_L},
\]
则偏微分是 \(L\)-测度变化的局部密度。于是 PDE 的局部微分表达可替代为 \(L\)-积分弱表达，
PDE 求解可替代为 \(L\)-路径测度低作用量命中，并在解释函子 \(\mathcal C_L\) 下保持语义等价。
第五章把“PDE 显式表达下的边际函数偏微控制”形式化迁移为
“PDE \(L\)-测度隐式表达、同伦扭曲、\TPHLHTS{} 微观隧穿命中与偏微下降工程兑现”的全局变分范式，
并证明二者在 PDE 解集合与 Law-Space 工程解释上等价，在表达域和方法搜索能力上形成优化。
第六章把全文增益登记为成熟度证书：有限证据域内强支持命题已经闭合，一般理论框架已经成型；
全称版本仍需测试函数分离性、投影忠实性、解释函子交换性与路径族覆盖性。
第七章进一步把传统 PDE 求解与 \(L\)-测度同伦隧穿求解放在同一形式逻辑中比较，
证明后者不是“更快算一个 PDE”的单点改良，而是把函数解搜索升级为
可审计、可兑现、可控风险的全局命中路径搜索，从而形成战略工程降维优势。

本文继承仓内 tex37 的“定义--证书--谓词因果链--备注”体例
\cite{RepoTex37SemanticGauge}，并与仓内 PDE 低残差命中范式、PDE 策略簇、\LHTS{} 与 \TPHLHTS{} 证书链相衔接
\cite{RepoTex55PDEParadigm,RepoTex53StrategyCluster,RepoTex56LHTS,RepoTex57TPHLHTS}。
关于 \GFramework{} 中 JacGap、障碍度量与工程目标的背景，参考
\cite{RepoAppMathVol1,RepoEconomicsVol}。外部数学背景方面，PDE、弱形式、有限元残差、泛函分析、
测度论、梯度流、数值优化与路径积分分别参考
\cite{Evans2010PDE,BrennerScott2008FEM,Brezis2011FunctionalAnalysis,Bogachev2007MeasureTheory,
AmbrosioGigliSavare2008GradientFlows,NocedalWright2006NumericalOptimization,FeynmanHibbs1965PathIntegrals,
FreidlinWentzell2012RandomPerturbations}。
本文只引用 \(NoteBook\) 下的 \(\_pub\) 稿件和 \(docs\) 发布稿，不引用 \(NoteBook/\) 内非 \(\_pub\) 目录稿件。

全文的最终形式化结论可登记为
\[
  \boxed{
  \begin{gathered}
  \mathcal D_{\mathrm{replayid}}
  \models
  \left(
  \Eff_{\partial\JacGap}
  \wedge
  \Hit_{\TPH}^{\partial_L\to\Phi}
  \wedge
  \Eff_{\mathrm{equity}}
  \right)
  \Rightarrow
  \Subst_{\PDE}^{\Lmeas}(\Tunnel_G)=\top,\\
  \Subst_{\PDE}^{\Lmeas}(\Tunnel_G)
  \Longleftrightarrow
  \ExprSubst_{\PDE}^{\Lmeas}
  \wedge
  \SolveSubst_{\PDE}^{\Lmeas}
  \wedge
  \InterpEquiv_{\PDE}^{\Lmeas}.
  \end{gathered}
  }
\]
这正是本文标题中“表达替代、求解替代与解释等价”的含义。

\setcounter{tocdepth}{3}
\setcounter{secnumdepth}{3}
\tableofcontents
\clearpage

%%%%%%%%%%%%%%%%%%%%%%%%%%%%%%%%%%%%%%%%%%%%%%%%%%%%%%%%%%%%
\section{形式逻辑：偏微残差下降与可审计增益兑现}
\label{sec:tex58-pde-residual-descent-gain}
%%%%%%%%%%%%%%%%%%%%%%%%%%%%%%%%%%%%%%%%%%%%%%%%%%%%%%%%%%%%

\begin{remark}[核心陈述（严格化后）]
设偏微系统由
\[
  \Pcal[u]=0,\qquad x\in\Omega,
\]
与边界或初值条件
\[
  \Bcal[u]=0,\qquad x\in\partial\Omega
\]
共同给出。若残差能量
\[
  E(u)
  =
  \frac12\norm{\Pcal[u]}_{\Hcal_P}^{2}
  +
  \frac{\lambda}{2}\norm{\Bcal[u]}_{\Hcal_B}^{2},
  \qquad \lambda>0,
\]
忠实度量偏微系统是否被满足，并且求解过程沿 \(E\) 的下降方向演化，则求解 PDE 的核心结构链为
\[
  \text{偏微残差障碍}
  \Longrightarrow
  \text{残差能量}
  \Longrightarrow
  \text{偏微下降}
  \Longrightarrow
  \text{残差消除或压低}
  \Longrightarrow
  \text{可审计增益兑现}.
\]
在线性增益传导
\[
  G(u)=G_{\max}-cE(u),\qquad c>0
\]
下，残差下降量与增益兑现量满足精确正比例关系。
\end{remark}

\begin{remark}[阅读导航与引用口径]
本章的写法对齐仓内 tex37 的“定义--证书--谓词因果链--备注”体例
\cite{RepoTex37SemanticGauge}。
关于 PDE 从显式求解对象转化为低残差集合、命中泛函与保真性证书的口径，参考
\cite{RepoTex55PDEParadigm}；关于 PDE 策略簇、残差、障碍下降与证书命中的软件层范式，参考
\cite{RepoTex53StrategyCluster}；关于 \GFramework{} 中 JacGap 与障碍度量的应用数学背景，参考
\cite{RepoAppMathVol1}。
本文只引用 \(NoteBook\) 下的 \(\_pub\) 稿件或 \(docs\) 下的发布稿，不引用工作区内非发布目录稿件。
\end{remark}

\subsection{对象域与严格化词典}
\label{subsec:tex58-objects}

\begin{definition}[偏微系统与函数空间]
\label{def:tex58-pde-system}
令 \(\Omega\subset\RR^d\) 为区域，令 \(\Vcal\) 为待求函数所在的函数空间。
设
\[
  \Pcal:\Vcal\to\Hcal_P,
  \qquad
  \Bcal:\Vcal\to\Hcal_B
\]
分别为偏微分算子与边界/初值算子，其中 \(\Hcal_P,\Hcal_B\) 是带范数的线性空间。
偏微系统记为
\[
  \mathsf{PDE}(\Pcal,\Bcal)
  \DefEq
  \{u\in\Vcal:\Pcal[u]=0,\ \Bcal[u]=0\}.
\]
\end{definition}

\begin{definition}[残差算子、加权残差范数与残差能量]
\label{def:tex58-residual-energy}
定义残差算子
\[
  R(u)
  \DefEq
  \bigl(\Pcal[u],\Bcal[u]\bigr).
\]
给定权重 \(\lambda>0\)，定义加权残差范数
\[
  \norm{R(u)}_{\lambda}^{2}
  \DefEq
  \norm{\Pcal[u]}_{\Hcal_P}^{2}
  +
  \lambda\norm{\Bcal[u]}_{\Hcal_B}^{2},
\]
并定义残差能量泛函
\[
  E(u)
  \DefEq
  \frac12\norm{R(u)}_{\lambda}^{2}
  =
  \frac12\norm{\Pcal[u]}_{\Hcal_P}^{2}
  +
  \frac{\lambda}{2}\norm{\Bcal[u]}_{\Hcal_B}^{2}.
\]
\end{definition}

\begin{definition}[严格解、阈值解与“求解偏微方程”]
\label{def:tex58-solve-pde}
定义严格解集
\[
  \Sol_0
  \DefEq
  \{u\in\Vcal:E(u)=0\}.
\]
给定工程阈值 \(\varepsilon\ge0\)，定义阈值解集
\[
  \Sol_{\varepsilon}
  \DefEq
  \{u\in\Vcal:\norm{R(u)}_{\lambda}\le\varepsilon\}
  =
  \left\{u\in\Vcal:E(u)\le\frac{\varepsilon^2}{2}\right\}.
\]
本文中，“求解偏微方程”指构造 \(u\in\Sol_0\)，或在工程口径下构造
\[
  u\in\Sol_{\varepsilon}
\]
并给出残差证书。
\end{definition}

\begin{definition}[偏微残差障碍与偏微下降]
\label{def:tex58-pde-descent}
称
\[
  \Rcal_{\mathrm{obs}}(u)\DefEq E(u)
\]
为偏微残差障碍。若 \(E\) 在函数空间 \(\Vcal\) 的 Hilbert 结构下 Fréchet 可微，
并且人工演化参数 \(\tau\) 下的求解轨道 \(u(\tau)\) 满足
\[
  \frac{\partial u}{\partial\tau}
  =
  -\nabla E(u)
  =
  -\frac{\delta E}{\delta u},
\]
则称该演化为连续偏微下降流。

离散情形下，若序列 \(u_0,u_1,\ldots\) 满足
\[
  E(u_{k+1})\le E(u_k),
  \qquad k=0,1,\ldots,
\]
则称其为有效离散偏微下降序列。
\end{definition}

\begin{definition}[可审计增益兑现]
\label{def:tex58-gain-realization}
令目标增益泛函为
\[
  G:\Vcal\to\RR.
\]
本文采用最小线性传导模型
\[
  G(u)=G_{\max}-cE(u),
  \qquad c>0.
\]
给定两个状态 \(u,v\in\Vcal\)，若存在证书包
\[
  \Pi_G(u,v)
  =
  \bigl(E(u),E(v),G(u),G(v),c\bigr)
\]
使审计器验证
\[
  \Audit\bigl(\Pi_G(u,v)\bigr)=\top
\]
并且
\[
  G(v)-G(u)=c\bigl(E(u)-E(v)\bigr),
\]
则称从 \(u\) 到 \(v\) 的残差下降已经兑现为可审计增益。
\end{definition}

\begin{remark}[词典压缩]
上述定义把三个自然语言词压缩为三个可判定对象：
\[
  \text{求解偏微方程}
  \equiv
  E(u)=0\ \text{或}\ E(u)\le\varepsilon^2/2,
\]
\[
  \text{偏微下降}
  \equiv
  \frac{\partial u}{\partial\tau}=-\nabla E(u)
  \ \text{或}\
  E(u_{k+1})\le E(u_k),
\]
\[
  \text{增益兑现}
  \equiv
  \Audit(\Pi_G)=\top
  \ \text{且}\
  \Delta G=c(-\Delta E).
\]
\end{remark}

\subsection{公理（降级为定理）：离散残差下降到增益兑现的归纳闭合}
\label{subsec:tex58-axiom-demoted-induction}

\begin{theorem}[公理降级：离散残差下降--增益兑现传导闭合]
\label{thm:tex58-discrete-gain-induction}
给定离散序列
\[
  u_0,u_1,\ldots,u_n
\]
以及线性增益
\[
  G_k\DefEq G(u_k)=G_{\max}-cE(u_k),
  \qquad c>0.
\]
若对所有 \(0\le k<n\) 有
\[
  E(u_{k+1})\le E(u_k),
\]
则对所有 \(0\le k<n\) 有
\[
  G_{k+1}\ge G_k,
\]
并且每一步精确满足
\[
  G_{k+1}-G_k
  =
  c\bigl(E(u_k)-E(u_{k+1})\bigr).
\]
若某一步满足严格下降
\[
  E(u_{k+1})<E(u_k),
\]
则对应增益严格增加
\[
  G_{k+1}>G_k.
\]
\end{theorem}

\begin{Certificate}[数学归纳法证书：残差下降步推出增益兑现步]
\label{cert:tex58-discrete-gain-induction}
定义谓词
\[
\begin{aligned}
  P(m):\quad
  &\forall k\in\{0,\ldots,m-1\},\
  E(u_{k+1})\le E(u_k)\\
  &\Rightarrow
  G_{k+1}\ge G_k
  \ \text{且}\
  G_{k+1}-G_k=c(E(u_k)-E(u_{k+1})).
\end{aligned}
\]
归纳证书登记为
\[
  \pi_{\mathrm{ind}}
  =
  \bigl(P(1),\ \forall m\in\{1,\ldots,n-1\},\ P(m)\Rightarrow P(m+1)\bigr).
\]
其中归纳推进只使用增益定义
\[
  G_k=G_{\max}-cE(u_k),
  \qquad c>0.
\]
\end{Certificate}

\begin{proof}[证明（数学归纳法证书；基于数学符号推演逻辑的谓词因果链）]
\textbf{基步 \(P(1)\)：}
由定义
\[
  G_1-G_0
  =
  \bigl(G_{\max}-cE(u_1)\bigr)
  -
  \bigl(G_{\max}-cE(u_0)\bigr).
\]
化简得
\[
  G_1-G_0
  =
  c\bigl(E(u_0)-E(u_1)\bigr).
\]
若 \(E(u_1)\le E(u_0)\)，则
\[
  E(u_0)-E(u_1)\ge0,
\]
所以
\[
  G_1-G_0\ge0,
  \qquad
  G_1\ge G_0.
\]
故 \(P(1)\) 成立。

\textbf{归纳步：}
假设 \(P(m)\) 成立。若还给定第 \(m\) 步残差下降
\[
  E(u_{m+1})\le E(u_m),
\]
则同样由增益定义可得
\[
  G_{m+1}-G_m
  =
  \bigl(G_{\max}-cE(u_{m+1})\bigr)
  -
  \bigl(G_{\max}-cE(u_m)\bigr)
  =
  c\bigl(E(u_m)-E(u_{m+1})\bigr).
\]
因为 \(c>0\) 且 \(E(u_m)-E(u_{m+1})\ge0\)，得到
\[
  G_{m+1}\ge G_m.
\]
这把 \(P(m)\) 扩展为 \(P(m+1)\)。

由数学归纳法，\(P(m)\) 对所有 \(m=1,\ldots,n\) 成立。
严格下降情形下，
\[
  E(u_k)-E(u_{k+1})>0
\]
进一步推出
\[
  G_{k+1}-G_k
  =
  c\bigl(E(u_k)-E(u_{k+1})\bigr)>0.
\]
证毕。
\end{proof}

\subsection{定理：求解偏微方程等价于残差能量归零}
\label{subsec:tex58-solution-energy-theorem}

\begin{theorem}[残差能量归零等价定理]
\label{thm:tex58-solution-energy-equivalence}
假设 \(\norm{\cdot}_{\Hcal_P}\) 与 \(\norm{\cdot}_{\Hcal_B}\) 均忠实，即
\[
  \norm{y}=0\Longleftrightarrow y=0.
\]
则对任意 \(u^\ast\in\Vcal\)，有
\[
  u^\ast\in\mathsf{PDE}(\Pcal,\Bcal)
  \Longleftrightarrow
  E(u^\ast)=0.
\]
若严格解存在，则
\[
  u^\ast\in\mathsf{PDE}(\Pcal,\Bcal)
  \Longleftrightarrow
  u^\ast\in\argmin_{u\in\Vcal}E(u)
  \ \text{且}\
  E(u^\ast)=0.
\]
\end{theorem}

\begin{Certificate}[证书：范数非负 + 忠实性 + 零能量最小性]
\label{cert:tex58-solution-energy-equivalence}
证书登记谓词链：
\[
\begin{aligned}
  S_1 &: \norm{\Pcal[u]}_{\Hcal_P}^{2}\ge0,\quad
         \norm{\Bcal[u]}_{\Hcal_B}^{2}\ge0,\quad \lambda>0,\\
  S_2 &: S_1\Rightarrow E(u)\ge0,\\
  S_3 &: \Pcal[u^\ast]=0\wedge\Bcal[u^\ast]=0
         \Rightarrow E(u^\ast)=0,\\
  S_4 &: E(u^\ast)=0
         \Rightarrow \norm{\Pcal[u^\ast]}_{\Hcal_P}^{2}=0
         \wedge\norm{\Bcal[u^\ast]}_{\Hcal_B}^{2}=0,\\
  S_5 &: S_4\wedge\text{范数忠实}
         \Rightarrow \Pcal[u^\ast]=0\wedge\Bcal[u^\ast]=0,\\
  S_6 &: S_2\wedge E(u^\ast)=0
         \Rightarrow u^\ast\in\argmin E.
\end{aligned}
\]
\end{Certificate}

\begin{proof}[证明（基于数学符号推演逻辑的谓词因果链）]
由范数非负性与 \(\lambda>0\)，对任意 \(u\in\Vcal\) 有
\[
  E(u)
  =
  \frac12\norm{\Pcal[u]}_{\Hcal_P}^{2}
  +
  \frac{\lambda}{2}\norm{\Bcal[u]}_{\Hcal_B}^{2}
  \ge0.
\]
这给出 \(S_1\Rightarrow S_2\)。

若 \(u^\ast\) 是偏微系统严格解，则
\[
  \Pcal[u^\ast]=0,
  \qquad
  \Bcal[u^\ast]=0.
\]
代入残差能量定义得
\[
  E(u^\ast)
  =
  \frac12\cdot0
  +
  \frac{\lambda}{2}\cdot0
  =
  0.
\]
这给出 \(S_3\)。

反向设 \(E(u^\ast)=0\)。由于
\[
  \frac12\norm{\Pcal[u^\ast]}_{\Hcal_P}^{2}\ge0,
  \qquad
  \frac{\lambda}{2}\norm{\Bcal[u^\ast]}_{\Hcal_B}^{2}\ge0,
\]
且二者之和为零，只能有
\[
  \norm{\Pcal[u^\ast]}_{\Hcal_P}^{2}=0,
  \qquad
  \norm{\Bcal[u^\ast]}_{\Hcal_B}^{2}=0.
\]
由范数忠实性得到
\[
  \Pcal[u^\ast]=0,
  \qquad
  \Bcal[u^\ast]=0.
\]
故 \(u^\ast\) 是严格解。

若严格解存在，则上面已经给出某个 \(u^\ast\) 满足 \(E(u^\ast)=0\)。
又由 \(E(u)\ge0\)，得到
\[
  E(u^\ast)=0=\inf_{u\in\Vcal}E(u),
\]
所以
\[
  u^\ast\in\argmin_{u\in\Vcal}E(u).
\]
反过来，若 \(u^\ast\in\argmin E\) 且 \(E(u^\ast)=0\)，由前述零能量推出严格解。
证毕。
\end{proof}

\begin{corollary}[阈值求解的残差证书]
\label{cor:tex58-threshold-solution}
对任意 \(\varepsilon\ge0\)，有
\[
  u\in\Sol_{\varepsilon}
  \Longleftrightarrow
  E(u)\le \frac{\varepsilon^2}{2}.
\]
因此工程意义下的 PDE 求解可以登记为残差证书
\[
  \Cert_{\varepsilon}(u)
  =
  \left(E(u),\varepsilon,\ E(u)\le\frac{\varepsilon^2}{2}\right).
\]
\end{corollary}

\begin{Certificate}[证书：阈值范数与能量阈值互换]
\label{cert:tex58-threshold-solution}
证书登记为
\[
  \norm{R(u)}_{\lambda}\le\varepsilon
  \Longleftrightarrow
  \norm{R(u)}_{\lambda}^{2}\le\varepsilon^2
  \Longleftrightarrow
  \frac12\norm{R(u)}_{\lambda}^{2}\le\frac{\varepsilon^2}{2}
  \Longleftrightarrow
  E(u)\le\frac{\varepsilon^2}{2}.
\]
\end{Certificate}

\begin{proof}[证明（基于数学符号推演逻辑的谓词因果链）]
由定义~\ref{def:tex58-residual-energy}，
\[
  E(u)=\frac12\norm{R(u)}_{\lambda}^{2}.
\]
当 \(\varepsilon\ge0\) 时，
\[
  \norm{R(u)}_{\lambda}\le\varepsilon
\]
等价于
\[
  \norm{R(u)}_{\lambda}^{2}\le\varepsilon^2.
\]
两端同除以 \(2\)，得到
\[
  E(u)\le\varepsilon^2/2.
\]
反向推理完全相同。证毕。
\end{proof}

\subsection{引理：偏微下降必然使残差能量不增加}
\label{subsec:tex58-gradient-descent-lemma}

\begin{lemma}[连续偏微下降的残差单调性]
\label{lem:tex58-continuous-descent}
设 \(E:\Vcal\to\RR_{\ge0}\) 在 Hilbert 空间结构下 Fréchet 可微。
若 \(u(\tau)\) 满足偏微下降流
\[
  \frac{\partial u}{\partial\tau}
  =
  -\nabla E(u),
\]
则
\[
  \frac{d}{d\tau}E(u(\tau))
  =
  -\norm{\nabla E(u(\tau))}^{2}
  \le0.
\]
若 \(\nabla E(u(\tau))\neq0\)，则
\[
  \frac{d}{d\tau}E(u(\tau))<0.
\]
\end{lemma}

\begin{Certificate}[证书：链式法则 + 下降流代入 + 内积正定]
\label{cert:tex58-continuous-descent}
证书登记谓词链：
\[
\begin{aligned}
  D_1 &: u_\tau=-\nabla E(u),\\
  D_2 &: \frac{d}{d\tau}E(u(\tau))=\ip{\nabla E(u)}{u_\tau},\\
  D_3 &: D_1\wedge D_2
         \Rightarrow
         \frac{dE}{d\tau}
         =
         \ip{\nabla E(u)}{-\nabla E(u)},\\
  D_4 &: D_3
         \Rightarrow
         \frac{dE}{d\tau}
         =
         -\norm{\nabla E(u)}^2\le0.
\end{aligned}
\]
\end{Certificate}

\begin{proof}[证明（基于数学符号推演逻辑的谓词因果链）]
由函数空间链式法则，
\[
  \frac{d}{d\tau}E(u(\tau))
  =
  \ip{\nabla E(u)}{\frac{\partial u}{\partial\tau}}.
\]
代入偏微下降流
\[
  \frac{\partial u}{\partial\tau}
  =
  -\nabla E(u),
\]
得到
\[
  \frac{d}{d\tau}E(u(\tau))
  =
  \ip{\nabla E(u)}{-\nabla E(u)}
  =
  -\norm{\nabla E(u)}^2.
\]
由于
\[
  \norm{\nabla E(u)}^2\ge0,
\]
可得
\[
  \frac{dE}{d\tau}\le0.
\]
若 \(\nabla E(u)\neq0\)，则 \(\norm{\nabla E(u)}^2>0\)，从而
\[
  \frac{dE}{d\tau}<0.
\]
证毕。
\end{proof}

\begin{lemma}[离散有效下降的残差单调性]
\label{lem:tex58-discrete-descent}
若离散序列满足
\[
  E(u_{k+1})\le E(u_k),
  \qquad k=0,1,\ldots,
\]
则残差障碍
\[
  \Rcal_{\mathrm{obs}}(u_k)=E(u_k)
\]
沿序列单调不增。
\end{lemma}

\begin{Certificate}[证书：残差障碍定义回代]
\label{cert:tex58-discrete-descent}
证书登记为
\[
  E(u_{k+1})\le E(u_k)
  \quad\wedge\quad
  \Rcal_{\mathrm{obs}}(u)=E(u)
  \Longrightarrow
  \Rcal_{\mathrm{obs}}(u_{k+1})\le \Rcal_{\mathrm{obs}}(u_k).
\]
\end{Certificate}

\begin{proof}[证明（基于数学符号推演逻辑的谓词因果链）]
由定义~\ref{def:tex58-pde-descent}，
\[
  \Rcal_{\mathrm{obs}}(u_k)=E(u_k),
  \qquad
  \Rcal_{\mathrm{obs}}(u_{k+1})=E(u_{k+1}).
\]
把离散下降条件
\[
  E(u_{k+1})\le E(u_k)
\]
直接回代，即得
\[
  \Rcal_{\mathrm{obs}}(u_{k+1})\le\Rcal_{\mathrm{obs}}(u_k).
\]
证毕。
\end{proof}

\subsection{命题：残差下降与增益兑现正相关}
\label{subsec:tex58-gain-correlation}

\begin{proposition}[连续时间正比例传导命题]
\label{prop:tex58-continuous-positive-correlation}
设
\[
  G(u)=G_{\max}-cE(u),
  \qquad c>0,
\]
且 \(u(\tau)\) 满足连续偏微下降流。定义残差下降强度
\[
  X(\tau)\DefEq -\frac{dE(u(\tau))}{d\tau}
\]
与增益兑现速率
\[
  Y(\tau)\DefEq \frac{dG(u(\tau))}{d\tau}.
\]
则
\[
  Y(\tau)=cX(\tau),
  \qquad
  X(\tau)\ge0,
  \qquad
  Y(\tau)\ge0.
\]
若在某个观测窗口上 \(X\) 非退化，即 \(\Var(X)>0\)，则
\[
  \Corr(X,Y)=1.
\]
\end{proposition}

\begin{Certificate}[证书：线性增益传导 + 梯度流能量恒等式]
\label{cert:tex58-continuous-positive-correlation}
证书登记谓词链：
\[
\begin{aligned}
  G_1 &: G=G_{\max}-cE,\ c>0,\\
  G_2 &: G_1\Rightarrow \frac{dG}{d\tau}=-c\frac{dE}{d\tau},\\
  G_3 &: \text{引理~\ref{lem:tex58-continuous-descent}}
         \Rightarrow -\frac{dE}{d\tau}=\norm{\nabla E(u)}^2\ge0,\\
  G_4 &: G_2\wedge G_3
         \Rightarrow
         \frac{dG}{d\tau}=c\norm{\nabla E(u)}^2
         =
         c\left(-\frac{dE}{d\tau}\right),\\
  G_5 &: Y=cX,\ c>0,\ \Var(X)>0
         \Rightarrow \Corr(X,Y)=1.
\end{aligned}
\]
\end{Certificate}

\begin{proof}[证明（基于数学符号推演逻辑的谓词因果链）]
由增益定义
\[
  G(u(\tau))=G_{\max}-cE(u(\tau))
\]
对 \(\tau\) 求导，得到
\[
  \frac{dG}{d\tau}
  =
  -c\frac{dE}{d\tau}.
\]
由引理~\ref{lem:tex58-continuous-descent}，
\[
  \frac{dE}{d\tau}
  =
  -\norm{\nabla E(u)}^2.
\]
因此
\[
  \frac{dG}{d\tau}
  =
  c\norm{\nabla E(u)}^2
  =
  c\left(-\frac{dE}{d\tau}\right).
\]
也就是说
\[
  Y(\tau)=cX(\tau).
\]
由于 \(c>0\) 且 \(X(\tau)=\norm{\nabla E(u)}^2\ge0\)，得
\[
  Y(\tau)\ge0.
\]
若在观测窗口上 \(\Var(X)>0\)，则
\[
  \Cov(X,Y)=\Cov(X,cX)=c\Var(X),
\]
\[
  \Var(Y)=\Var(cX)=c^2\Var(X).
\]
故
\[
  \Corr(X,Y)
  =
  \frac{\Cov(X,Y)}{\sqrt{\Var(X)\Var(Y)}}
  =
  \frac{c\Var(X)}{\sqrt{\Var(X)c^2\Var(X)}}
  =
  1.
\]
证毕。
\end{proof}

\begin{proposition}[离散时间正比例传导命题]
\label{prop:tex58-discrete-positive-correlation}
令
\[
  \Delta E_k^{-}\DefEq E(u_k)-E(u_{k+1}),
  \qquad
  \Delta G_k\DefEq G_{k+1}-G_k.
\]
在线性增益传导下，
\[
  \Delta G_k=c\Delta E_k^{-}.
\]
因此只要 \(\Delta E_k^{-}\ge0\)，就有 \(\Delta G_k\ge0\)；
若有限样本窗口中 \(\Delta E_k^{-}\) 非退化，则
\[
  \Corr(\Delta E^{-},\Delta G)=1.
\]
\end{proposition}

\begin{Certificate}[证书：一步代数恒等式 + 非退化相关系数]
\label{cert:tex58-discrete-positive-correlation}
证书登记为
\[
  \Delta G_k
  =
  \bigl(G_{\max}-cE(u_{k+1})\bigr)
  -
  \bigl(G_{\max}-cE(u_k)\bigr)
  =
  c\bigl(E(u_k)-E(u_{k+1})\bigr)
  =
  c\Delta E_k^{-}.
\]
\end{Certificate}

\begin{proof}[证明（基于数学符号推演逻辑的谓词因果链）]
由证书中的一步代数恒等式，直接得到
\[
  \Delta G_k=c\Delta E_k^{-}.
\]
当 \(\Delta E_k^{-}\ge0\) 且 \(c>0\) 时，
\[
  \Delta G_k\ge0.
\]
若样本窗口中 \(\Var(\Delta E^{-})>0\)，则与命题~\ref{prop:tex58-continuous-positive-correlation}
的相关系数计算相同：
\[
  \Corr(\Delta E^{-},\Delta G)=\Corr(\Delta E^{-},c\Delta E^{-})=1.
\]
证毕。
\end{proof}

\subsection{定理：PDE 求解是残差障碍到增益兑现的偏微下降过程}
\label{subsec:tex58-main-theorem}

\begin{theorem}[偏微残差障碍转化为可审计增益兑现定理]
\label{thm:tex58-main}
设以下条件成立：
\[
\begin{aligned}
  A_1 &: E(u)=0\Longleftrightarrow \Pcal[u]=0\wedge\Bcal[u]=0,\\
  A_2 &: u(\tau)\ \text{满足连续偏微下降流，或 } \{u_k\}\text{ 满足有效离散偏微下降},\\
  A_3 &: G(u)=G_{\max}-cE(u),\quad c>0,\\
  A_4 &: \Audit(\Pi_G)=\top\ \text{可验证}.
\end{aligned}
\]
则 PDE 求解过程可形式化为
\[
  \boxed{
  \text{残差障碍 }E
  \xrightarrow{\text{偏微下降}}
  \text{残差压低或消除}
  \xrightarrow{\text{线性传导}}
  \text{可审计增益兑现}.
  }
\]
在连续非退化情形下，
\[
  \frac{dG}{d\tau}
  =
  c\left(-\frac{dE}{d\tau}\right),
  \qquad c>0,
\]
并且残差下降强度与增益兑现速率正相关。
\end{theorem}

\begin{Certificate}[总证书：残差忠实 + 下降有效 + 传导审计]
\label{cert:tex58-main}
证书登记谓词链：
\[
\begin{aligned}
  M_1 &: A_1\Rightarrow \text{求解 PDE 等价于 }E(u)=0
         \quad\text{（定理~\ref{thm:tex58-solution-energy-equivalence}）},\\
  M_2 &: A_2\Rightarrow E\text{ 沿求解过程不增加}
         \quad\text{（引理~\ref{lem:tex58-continuous-descent} 或 \ref{lem:tex58-discrete-descent}）},\\
  M_3 &: A_3\Rightarrow \Delta G=c(-\Delta E)\\
       &\quad\text{（命题~\ref{prop:tex58-continuous-positive-correlation}}\\
       &\quad\text{或命题~\ref{prop:tex58-discrete-positive-correlation}）},\\
  M_4 &: A_4\Rightarrow \Delta G\text{ 可被证书包审计},\\
  M_5 &: M_1\wedge M_2\wedge M_3\wedge M_4
         \Rightarrow
         \text{PDE 求解是残差障碍到增益兑现的偏微下降过程}.
\end{aligned}
\]
\end{Certificate}

\begin{proof}[证明（基于数学符号推演逻辑的谓词因果链）]
由 \(A_1\) 与定理~\ref{thm:tex58-solution-energy-equivalence}，严格求解
\[
  \Pcal[u]=0,\qquad \Bcal[u]=0
\]
等价于
\[
  E(u)=0.
\]
工程阈值求解则由推论~\ref{cor:tex58-threshold-solution} 表示为
\[
  E(u)\le\varepsilon^2/2.
\]
因此 \(E\) 是求解目标的障碍度量，即 \(M_1\) 成立。

由 \(A_2\)，若采用连续偏微下降流，则引理~\ref{lem:tex58-continuous-descent} 给出
\[
  \frac{dE}{d\tau}\le0.
\]
若采用离散下降序列，则引理~\ref{lem:tex58-discrete-descent} 给出
\[
  E(u_{k+1})\le E(u_k).
\]
因此偏微下降使残差障碍不增加，即 \(M_2\) 成立。

由 \(A_3\)，增益被定义为
\[
  G=G_{\max}-cE.
\]
命题~\ref{prop:tex58-continuous-positive-correlation} 与
命题~\ref{prop:tex58-discrete-positive-correlation} 分别给出连续与离散情形下的传导关系：
\[
  \frac{dG}{d\tau}
  =
  c\left(-\frac{dE}{d\tau}\right),
  \qquad
  \Delta G_k=c\Delta E_k^{-}.
\]
这说明残差下降被正比例传导为增益兑现，即 \(M_3\) 成立。

由 \(A_4\)，证书包
\[
  \Pi_G
  =
  (E_{\mathrm{before}},E_{\mathrm{after}},
    G_{\mathrm{before}},G_{\mathrm{after}},c)
\]
可被审计器验证，因此该增益不是不可检查的口头收益，而是可审计增益，即 \(M_4\) 成立。

合并 \(M_1,M_2,M_3,M_4\)，得到结构链
\[
  \text{偏微残差}
  \Longrightarrow
  \text{残差能量障碍}
  \Longrightarrow
  \text{偏微下降}
  \Longrightarrow
  \text{残差消除或压低}
  \Longrightarrow
  \text{可审计增益兑现}.
\]
在连续非退化观测窗口中，命题~\ref{prop:tex58-continuous-positive-correlation} 进一步给出
\[
  \Corr\left(
  -\frac{dE}{d\tau},
  \frac{dG}{d\tau}
  \right)=1.
\]
定理成立。证毕。
\end{proof}

\subsection{推论：G 框架中的障碍--修复--收益链}
\label{subsec:tex58-gframework-corollary}

\begin{corollary}[G 框架语义下的 PDE 求解结构链]
\label{cor:tex58-gframework-chain}
在 \GFramework{} 语义下，可将本文的残差下降结构解释为
\[
  \text{局部微分法则不闭合}
  \Longleftrightarrow
  \text{偏微残差障碍},
\]
\[
  E(u)
  \Longleftrightarrow
  \text{障碍能量或 JacGap 型残差度量},
\]
\[
  -\nabla E(u)
  \Longleftrightarrow
  \text{障碍修复方向},
\]
\[
  G_{k+1}-G_k
  =
  c\bigl(E(u_k)-E(u_{k+1})\bigr)
  \Longleftrightarrow
  \text{障碍下降向收益兑现的传导证书}.
\]
因此，PDE 求解在该语义下不是单纯“找到一个函数”，而是完成
\[
  \boxed{
  \text{局部微分障碍}
  \to
  \text{残差度量}
  \to
  \text{偏微下降}
  \to
  \text{结构修复}
  \to
  \text{目标增益兑现}.
  }
\]
\end{corollary}

\begin{Certificate}[证书：偏微残差到 G 框架障碍语言的保守解释]
\label{cert:tex58-gframework-chain}
证书登记谓词链：
\[
\begin{aligned}
  C_1 &: \Pcal[u]\ne0\ \text{或}\ \Bcal[u]\ne0
         \Rightarrow E(u)>0
         \quad\text{（残差忠实）},\\
  C_2 &: E(u)>0\Rightarrow \text{存在待压低的障碍量},\\
  C_3 &: -\nabla E(u)\Rightarrow \text{给出一阶障碍下降方向},\\
  C_4 &: E(u_{k+1})\le E(u_k)\Rightarrow \text{障碍不增加},\\
  C_5 &: G_{k+1}-G_k=c(E(u_k)-E(u_{k+1}))
         \Rightarrow \text{障碍下降被登记为收益兑现},\\
  C_6 &: C_1\wedge C_2\wedge C_3\wedge C_4\wedge C_5
         \Rightarrow \text{G 框架障碍--修复--收益链成立}.
\end{aligned}
\]
\end{Certificate}

\begin{proof}[证明（基于数学符号推演逻辑的谓词因果链）]
由残差忠实性，若偏微方程或边界条件未被满足，则
\[
  R(u)\ne0
\]
从而
\[
  E(u)>0.
\]
这把局部微分法则不闭合登记为正障碍量。
由偏微下降定义，\(-\nabla E(u)\) 是降低该障碍量的一阶方向。
若连续或离散下降有效，则
\[
  E_{\mathrm{after}}\le E_{\mathrm{before}}.
\]
再由线性增益传导，
\[
  G_{\mathrm{after}}-G_{\mathrm{before}}
  =
  c(E_{\mathrm{before}}-E_{\mathrm{after}}).
\]
因此每一段可审计的残差下降都被转写为可审计收益增加。
这正是 \GFramework{} 中障碍--修复--收益链的保守解释。证毕。
\end{proof}

\subsection{备注：边界条件、非退化性与最终口径}
\label{subsec:tex58-remarks}

\begin{remark}[成立条件]
本文主命题成立需要四类条件。

第一，残差能量必须忠实：
\[
  E(u)=0
  \Longleftrightarrow
  \Pcal[u]=0\wedge\Bcal[u]=0.
\]
这要求残差范数忠实，且边界/初值条件与方程系统相容。

第二，下降过程必须有效：
\[
  \frac{dE}{d\tau}\le0
  \quad\text{或}\quad
  E(u_{k+1})\le E(u_k).
\]
若数值格式不稳定、步长过大、投影不相容或离散残差不能代表连续残差，则只能登记为尝试求解，
不能登记为已经兑现的残差下降。

第三，增益函数必须与残差下降单调对应。本文采用最小线性模型
\[
  G=G_{\max}-cE,\qquad c>0.
\]
更一般地，只要
\[
  \frac{dG}{dE}<0
\]
且审计器能验证对应传导，也可得到单调正相关；但精确相关系数等于 \(1\) 的结论依赖线性传导模型。

第四，正相关结论需要非退化观测窗口：
\[
  \Var\left(-\frac{dE}{d\tau}\right)>0
  \quad\text{或}\quad
  \Var(\Delta E^{-})>0.
\]
若残差已经完全停止变化，则比例恒等式仍成立，但相关系数作为统计量不再有非退化意义。
\end{remark}

\begin{remark}[最终严格命题]
本文最终命题可压缩为：
\[
  \boxed{
  \begin{gathered}
  \text{在残差忠实、边界相容、下降有效、增益单调传导与证书可审计的条件下，}\\
  \text{求解偏微方程等价于把偏微残差障碍经偏微下降压低或消除，}\\
  \text{并把该下降量按正传导系数兑现为可验证目标增益。}
  \end{gathered}
  }
\]
在线性传导模型下，其最短数学表达为
\[
  \boxed{
  \frac{dG}{d\tau}
  =
  c\left(-\frac{dE}{d\tau}\right),
  \qquad c>0.
  }
\]
离散情形下对应为
\[
  \boxed{
  G_{k+1}-G_k
  =
  c\bigl(E(u_k)-E(u_{k+1})\bigr),
  \qquad c>0.
  }
\]
\end{remark}

\begin{remark}[与仓内 PDE 命中范式的接口]
仓内 tex55 已经把高阶 PDE 求解从“直接显式求解”转化为“命中低残差集合”的问题
\cite{RepoTex55PDEParadigm}。本文提供的是该命中范式的函数空间底层解释：
\[
  \text{低残差命中}
  \Longleftrightarrow
  E(u)\le\varepsilon^2/2.
\]
仓内 tex53 进一步把 PDE 方法空间写成策略簇、残差、障碍下降与证书命中的动态系统
\cite{RepoTex53StrategyCluster}。本文提供的是该策略簇的最小收益传导接口：
\[
  \text{策略有效}
  \Longrightarrow
  E_{\mathrm{after}}\le E_{\mathrm{before}}
  \Longrightarrow
  \Delta G=c(E_{\mathrm{before}}-E_{\mathrm{after}})\ge0.
\]
因此，本文不是替换 tex53 或 tex55，而是给出它们共享的残差--下降--增益底层证书。
\end{remark}

\clearpage

%%%%%%%%%%%%%%%%%%%%%%%%%%%%%%%%%%%%%%%%%%%%%%%%%%%%%%%%%%%%
\section{形式逻辑：有限回放证据域中的 TPH-LHTS 传导与偏微正影响证书}
\label{sec:tex58-tph-lhts-finite-replay-impact}
%%%%%%%%%%%%%%%%%%%%%%%%%%%%%%%%%%%%%%%%%%%%%%%%%%%%%%%%%%%%

\begin{remark}[核心陈述（严格化后）]
在有限回放证据域 \(\mathcal D_{\mathrm{replayid}}\) 内，整体工程兑现
\[
  \Eff_{\mathrm{equity}}=\top
\]
只登记整体工程目标泛函已经相对初始基准正向兑现；偏微/微分下降
\[
  \Eff_{\partial\JacGap}=\top
\]
只登记 \(\partial_L\) 作用下的 JacGap 累计下降有效。要推出
\[
  \sum_{t\in T}I^+_{\partial_L}(t)
  >
  \sum_{t\in T}I^-_{\partial_L}(t),
\]
还必须引入桥接谓词
\[
  \Hit_{\TPH}^{\partial_L\to\Phi}=\top.
\]
本文中该谓词不只表示“TPH 发生命中”，而是表示 \TPHLHTS{} 命中把
\(\partial_L\) 的偏微/微分下降效果传导到整体工程目标泛函 \(\Phi\)。
因此严格命题写成
\[
  \boxed{
  \mathcal D_{\mathrm{replayid}}
  \models
  \left(
  \Eff_{\partial\JacGap}
  \wedge
  \Hit_{\TPH}^{\partial_L\to\Phi}
  \wedge
  \Eff_{\mathrm{equity}}
  \right)
  \Rightarrow
  \sum_{t\in T}I^+_{\partial_L}(t)
  >
  \sum_{t\in T}I^-_{\partial_L}(t).
  }
\]
\end{remark}

\begin{remark}[阅读导航与引用口径]
本章继承第~\ref{sec:tex58-pde-residual-descent-gain} 章的残差--下降--增益传导口径。
\TPHLHTS{} 的“\LHTS{} 作为嵌入运动内核、\TPH{} 作为轨迹函数参数同伦控制器”的分工，
参考仓内 \(\_pub\) 稿件 \cite{RepoTex56LHTS,RepoTex57TPHLHTS}；
PDE 低残差命中、命中泛函与保真性层级参考 \cite{RepoTex55PDEParadigm}；
策略簇、残差、障碍下降与命中证书参考 \cite{RepoTex53StrategyCluster}；
JacGap 作为结构权重和风险/收益相关结构量的仓内经济学背景参考
\cite{RepoEconomicsVol}。外部风险调整目标的背景可参考
\cite{Markowitz1952PortfolioSelection,RockafellarUryasev2000CVaR}。
本章仍只引用 \(NoteBook\) 下 \(\_pub\) 稿件与 \(docs\) 下发布稿，不引用 \(NoteBook/\) 内非 \(\_pub\) 目录稿件。
\end{remark}

\subsection{证据域、工程目标与兑现谓词}
\label{subsec:tex58-replay-domain-equity}

\begin{definition}[有限回放证据域]
\label{def:tex58-replay-domain}
设有限回放证据域为
\[
  \mathcal D_{\mathrm{replayid}}
  \DefEq
  \bigl(T,S_t,A_t,\JacGap_t,\TPH_t,\Eq_t,\Bal_t,\Risk_t\bigr)_{t=0}^{N},
\]
其中
\[
  T=\{0,1,\ldots,N\}.
\]
这里 \(S_t\) 是第 \(t\) 个状态，\(A_t\) 是第 \(t\) 个动作或算子执行记录，
\(\JacGap_t\) 是第 \(t\) 个状态下的雅可比间隙，\(\TPH_t\) 是第 \(t\) 个状态下的
\TPH{} 命中记录，\(\Eq_t\) 是净值，\(\Bal_t\) 是余额，\(\Risk_t\) 汇总回撤、
敞口、未出清风险、锁仓风险、保证金风险与交易成本等风险变量。
\end{definition}

\begin{definition}[整体工程目标泛函]
\label{def:tex58-engineering-objective}
定义整体工程目标泛函
\[
  \Phi_t
  \DefEq
  \alpha \Eq_t
  +
  \beta \Bal_t
  -
  \gamma_1 \DD_t
  -
  \gamma_2 \Exposure_t
  -
  \gamma_3 \UnclearRisk_t
  -
  \gamma_4 \Cost_t,
\]
其中
\[
  \alpha,\beta,\gamma_1,\gamma_2,\gamma_3,\gamma_4>0.
\]
因此 \(\Phi_t\) 不是单纯净值，而是综合净值、余额、回撤、敞口、未出清风险与成本之后的
风险调整工程目标。
\end{definition}

\begin{definition}[整体工程兑现谓词]
\label{def:tex58-equity-effectiveness}
定义
\[
  \Eff_{\mathrm{equity}}=\top
\]
当且仅当
\[
  \Phi_N-\Phi_0>0.
\]
若只使用净值与余额的弱化证据，则可登记
\[
  \Eq_N>\Eq_0,
  \qquad
  \Bal_N>\Bal_0.
\]
但本章采用风险调整泛函 \(\Phi\) 作为严格兑现口径。
\end{definition}

\begin{Certificate}[证书：工程目标正向兑现]
\label{cert:tex58-equity-effectiveness}
证书登记谓词链：
\[
\begin{aligned}
  E_1 &: \Phi_t
  =
  \alpha \Eq_t+\beta \Bal_t
  -\gamma_1\DD_t-\gamma_2\Exposure_t
  -\gamma_3\UnclearRisk_t-\gamma_4\Cost_t,\\
  E_2 &: \Phi_N-\Phi_0>0,\\
  E_3 &: E_1\wedge E_2
  \Rightarrow
  \mathcal D_{\mathrm{replayid}}\models \Eff_{\mathrm{equity}}=\top.
\end{aligned}
\]
\end{Certificate}

\begin{proof}[证明（基于数学符号推演逻辑的谓词因果链）]
由定义~\ref{def:tex58-engineering-objective}，\(\Phi_t\) 汇总净值、余额与风险扣减项。
由定义~\ref{def:tex58-equity-effectiveness}，
\[
  \Eff_{\mathrm{equity}}=\top
  \Longleftrightarrow
  \Phi_N-\Phi_0>0.
\]
因此 \(E_1\) 与 \(E_2\) 合取即推出
\[
  \mathcal D_{\mathrm{replayid}}\models\Eff_{\mathrm{equity}}=\top.
\]
证毕。
\end{proof}

\subsection{偏微 JacGap 下降与累计有效性}
\label{subsec:tex58-partial-jacgap-effectiveness}

\begin{definition}[偏微/微分下降算子与 JacGap 变化]
\label{def:tex58-partial-jacgap-delta}
令
\[
  \partial_L
\]
表示作用在 \LHTS{}/\TPH{} 结构上的偏微/微分下降算子。它作用的对象不是价格本身，
而是结构性残差或障碍量，例如 \(\JacGap_t\)。定义
\[
  \Delta_{\partial_L}\JacGap_t
  \DefEq
  \JacGap_{t+1}^{(\partial_L)}-\JacGap_t.
\]
若
\[
  \Delta_{\partial_L}\JacGap_t<0,
\]
则表示第 \(t\) 个 tick 上 \(\partial_L\) 使结构障碍下降。
\end{definition}

\begin{definition}[正下降量、反向增障量与累计有效性]
\label{def:tex58-positive-negative-jacgap}
定义正下降量
\[
  D^+_{\partial_L}(t)
  \DefEq
  [-\Delta_{\partial_L}\JacGap_t]_+
  =
  \max\{-\Delta_{\partial_L}\JacGap_t,0\},
\]
定义反向增障量
\[
  D^-_{\partial_L}(t)
  \DefEq
  [\Delta_{\partial_L}\JacGap_t]_+
  =
  \max\{\Delta_{\partial_L}\JacGap_t,0\}.
\]
定义
\[
  \Eff_{\partial\JacGap}=\top
\]
当且仅当
\[
  \sum_{t=0}^{N-1}D^+_{\partial_L}(t)
  >
  \sum_{t=0}^{N-1}D^-_{\partial_L}(t).
\]
\end{definition}

\begin{lemma}[累计偏微下降等价于 JacGap 净变化为负]
\label{lem:tex58-partial-jacgap-equivalence}
在定义~\ref{def:tex58-positive-negative-jacgap} 下，
\[
  \Eff_{\partial\JacGap}=\top
  \Longleftrightarrow
  \sum_{t=0}^{N-1}\Delta_{\partial_L}\JacGap_t<0.
\]
\end{lemma}

\begin{Certificate}[证书：正负部分分解]
\label{cert:tex58-partial-jacgap-equivalence}
对任意实数 \(x\)，有
\[
  x=[x]_+-[-x]_+.
\]
令
\[
  x_t=\Delta_{\partial_L}\JacGap_t.
\]
则
\[
  \Delta_{\partial_L}\JacGap_t
  =
  D^-_{\partial_L}(t)-D^+_{\partial_L}(t).
\]
\end{Certificate}

\begin{proof}[证明（基于数学符号推演逻辑的谓词因果链）]
由正负部分分解，
\[
  \Delta_{\partial_L}\JacGap_t
  =
  [\Delta_{\partial_L}\JacGap_t]_+
  -
  [-\Delta_{\partial_L}\JacGap_t]_+
  =
  D^-_{\partial_L}(t)-D^+_{\partial_L}(t).
\]
对 \(t=0,\ldots,N-1\) 求和，得到
\[
  \sum_{t=0}^{N-1}\Delta_{\partial_L}\JacGap_t
  =
  \sum_{t=0}^{N-1}D^-_{\partial_L}(t)
  -
  \sum_{t=0}^{N-1}D^+_{\partial_L}(t).
\]
因此
\[
  \sum_{t=0}^{N-1}D^+_{\partial_L}(t)
  >
  \sum_{t=0}^{N-1}D^-_{\partial_L}(t)
\]
等价于
\[
  \sum_{t=0}^{N-1}\Delta_{\partial_L}\JacGap_t<0.
\]
证毕。
\end{proof}

\begin{remark}[累计有效而非逐 tick 有效]
\(\Eff_{\partial\JacGap}=\top\) 的含义是有限回放区间上的累计降障量大于累计增障量。
它不要求
\[
  \forall t,\quad \Delta_{\partial_L}\JacGap_t<0.
\]
因此某些 tick 可以短期增障，只要区间累计净变化为负，偏微/微分下降机制仍在累计意义上成立。
\end{remark}

\subsection{偏微归因贡献与正负影响}
\label{subsec:tex58-partial-impact-decomposition}

\begin{definition}[偏微下降对整体工程目标的归因贡献]
\label{def:tex58-partial-contribution}
定义 \(\partial_L\) 在第 \(t\) 个 tick 上通过 \TPHLHTS{} 传导到目标泛函 \(\Phi\) 的归因贡献为
\[
  C_{\partial_L}(t)
  \DefEq
  \Delta_{\partial_L}\Phi_t.
\]
定义偏微正影响
\[
  I^+_{\partial_L}(t)
  \DefEq
  [C_{\partial_L}(t)]_+
  =
  \max\{C_{\partial_L}(t),0\},
\]
定义偏微负影响
\[
  I^-_{\partial_L}(t)
  \DefEq
  [-C_{\partial_L}(t)]_+
  =
  \max\{-C_{\partial_L}(t),0\}.
\]
\end{definition}

\begin{lemma}[偏微归因贡献的正负分解]
\label{lem:tex58-impact-decomposition}
对每个 \(t\)，有
\[
  C_{\partial_L}(t)
  =
  I^+_{\partial_L}(t)-I^-_{\partial_L}(t).
\]
因此
\[
  \sum_{t=0}^{N-1}C_{\partial_L}(t)
  =
  \sum_{t=0}^{N-1}I^+_{\partial_L}(t)
  -
  \sum_{t=0}^{N-1}I^-_{\partial_L}(t).
\]
特别地，
\[
  \sum_{t=0}^{N-1}I^+_{\partial_L}(t)
  >
  \sum_{t=0}^{N-1}I^-_{\partial_L}(t)
  \Longleftrightarrow
  \sum_{t=0}^{N-1}C_{\partial_L}(t)>0.
\]
\end{lemma}

\begin{Certificate}[证书：贡献正负部分代数恒等式]
\label{cert:tex58-impact-decomposition}
证书登记为
\[
  C=[C]_+-[-C]_+
  \quad\Longrightarrow\quad
  C_{\partial_L}(t)
  =
  I^+_{\partial_L}(t)-I^-_{\partial_L}(t).
\]
\end{Certificate}

\begin{proof}[证明（基于数学符号推演逻辑的谓词因果链）]
对任意实数 \(C\)，正负部分满足
\[
  C=[C]_+-[-C]_+.
\]
取 \(C=C_{\partial_L}(t)\)，并使用定义~\ref{def:tex58-partial-contribution}，得
\[
  C_{\partial_L}(t)
  =
  I^+_{\partial_L}(t)-I^-_{\partial_L}(t).
\]
对所有 \(t\) 求和，得到
\[
  \sum_t C_{\partial_L}(t)
  =
  \sum_t I^+_{\partial_L}(t)
  -
  \sum_t I^-_{\partial_L}(t).
\]
于是
\[
  \sum_t I^+_{\partial_L}(t)>\sum_t I^-_{\partial_L}(t)
\]
当且仅当
\[
  \sum_t C_{\partial_L}(t)>0.
\]
证毕。
\end{proof}

\subsection{公理（降级为定理）：累计净贡献的归纳闭合}
\label{subsec:tex58-impact-induction}

\begin{theorem}[公理降级：有限回放累计净贡献归纳公式]
\label{thm:tex58-impact-induction}
定义单步偏微归因净贡献
\[
  c_t
  \DefEq
  I^+_{\partial_L}(t)-I^-_{\partial_L}(t),
\]
并定义累计净贡献
\[
  K_n
  \DefEq
  \sum_{t=0}^{n}c_t,
  \qquad 0\le n\le N-1.
\]
则对所有 \(0\le n\le N-1\)，有
\[
  K_n
  =
  \sum_{t=0}^{n}
  \left(
  I^+_{\partial_L}(t)-I^-_{\partial_L}(t)
  \right).
\]
并且
\[
  K_{N-1}>0
  \Longleftrightarrow
  \sum_{t=0}^{N-1}I^+_{\partial_L}(t)
  >
  \sum_{t=0}^{N-1}I^-_{\partial_L}(t).
\]
\end{theorem}

\begin{Certificate}[数学归纳法证书：累计净贡献递推]
\label{cert:tex58-impact-induction}
定义谓词
\[
  P(n):
  K_n=
  \sum_{t=0}^{n}c_t.
\]
归纳证书登记为
\[
  \pi_{\mathrm{impact}}
  =
  \bigl(P(0),\ \forall n\in\{0,\ldots,N-2\},\ P(n)\Rightarrow P(n+1)\bigr),
\]
其中递推关系为
\[
  K_{n+1}=K_n+c_{n+1}.
\]
\end{Certificate}

\begin{proof}[证明（数学归纳法证书；基于数学符号推演逻辑的谓词因果链）]
\textbf{基步：}
由定义，
\[
  K_0=c_0.
\]
因此
\[
  K_0=\sum_{t=0}^{0}c_t,
\]
故 \(P(0)\) 成立。

\textbf{归纳步：}
假设 \(P(n)\) 成立，即
\[
  K_n=\sum_{t=0}^{n}c_t.
\]
由递推关系，
\[
  K_{n+1}=K_n+c_{n+1}.
\]
代入归纳假设得
\[
  K_{n+1}
  =
  \sum_{t=0}^{n}c_t+c_{n+1}
  =
  \sum_{t=0}^{n+1}c_t.
\]
故 \(P(n)\Rightarrow P(n+1)\)。

由数学归纳法，\(P(n)\) 对所有 \(0\le n\le N-1\) 成立。
再由
\[
  c_t=I^+_{\partial_L}(t)-I^-_{\partial_L}(t),
\]
得到
\[
  K_{N-1}
  =
  \sum_{t=0}^{N-1}I^+_{\partial_L}(t)
  -
  \sum_{t=0}^{N-1}I^-_{\partial_L}(t).
\]
因此 \(K_{N-1}>0\) 等价于
\[
  \sum_{t=0}^{N-1}I^+_{\partial_L}(t)
  >
  \sum_{t=0}^{N-1}I^-_{\partial_L}(t).
\]
证毕。
\end{proof}

\begin{remark}[归纳证书的含义]
该归纳法只证明有限回放区间上的累计净贡献如何由单步贡献相加得到。
它不要求
\[
  \forall t,\quad I^+_{\partial_L}(t)>I^-_{\partial_L}(t).
\]
所以短期浮亏、回撤扩大、敞口增加、未出清风险增加，都不破坏累计正贡献命题；
它们只进入 \(I^-_{\partial_L}(t)\) 并接受总和审计。
\end{remark}

\subsection{TPH-LHTS 传导证书与非退化桥接}
\label{subsec:tex58-tph-lhts-bridge}

\begin{definition}[TPH-LHTS 传导命中谓词]
\label{def:tex58-tph-bridge-hit}
定义
\[
  \Hit_{\TPH}^{\partial_L\to\Phi}=\top
\]
当且仅当存在传导系数
\[
  \kappa>0,
  \qquad
  \varepsilon\ge0,
\]
使得
\[
  \sum_{t=0}^{N-1}\Delta_{\partial_L}\Phi_t
  \ge
  \kappa
  \sum_{t=0}^{N-1}[-\Delta_{\partial_L}\JacGap_t]_+
  -
  \varepsilon
  \sum_{t=0}^{N-1}[\Delta_{\partial_L}\JacGap_t]_+.
\]
这里 \(\kappa\) 表示 JacGap 下降向工程目标泛函的正向传导率，
\(\varepsilon\) 表示偏微扭曲、执行噪声与风险暴露引入的负向损耗率。
\end{definition}

\begin{definition}[非退化传导条件]
\label{def:tex58-nondegenerate-bridge}
称 \(\TPHLHTS{}\) 传导在 \(\mathcal D_{\mathrm{replayid}}\) 内非退化，若
\[
  \kappa
  \sum_{t=0}^{N-1}[-\Delta_{\partial_L}\JacGap_t]_+
  >
  \varepsilon
  \sum_{t=0}^{N-1}[\Delta_{\partial_L}\JacGap_t]_+.
\]
\end{definition}

\begin{proposition}[TPH-LHTS 非退化传导推出目标归因正和]
\label{prop:tex58-bridge-positive-attribution}
若
\[
  \Hit_{\TPH}^{\partial_L\to\Phi}=\top
\]
且非退化传导条件成立，则
\[
  \sum_{t=0}^{N-1}\Delta_{\partial_L}\Phi_t>0.
\]
\end{proposition}

\begin{Certificate}[证书：传导下界 + 非退化严格正性]
\label{cert:tex58-bridge-positive-attribution}
证书登记谓词链：
\[
\begin{aligned}
  H_1 &: \Hit_{\TPH}^{\partial_L\to\Phi}=\top,\\
  H_2 &: H_1\Rightarrow
  \sum_t\Delta_{\partial_L}\Phi_t
  \ge
  \kappa\sum_t[-\Delta_{\partial_L}\JacGap_t]_+
  -
  \varepsilon\sum_t[\Delta_{\partial_L}\JacGap_t]_+,\\
  H_3 &: \kappa\sum_t[-\Delta_{\partial_L}\JacGap_t]_+
  >
  \varepsilon\sum_t[\Delta_{\partial_L}\JacGap_t]_+,\\
  H_4 &: H_2\wedge H_3
  \Rightarrow
  \sum_t\Delta_{\partial_L}\Phi_t>0.
\end{aligned}
\]
\end{Certificate}

\begin{proof}[证明（基于数学符号推演逻辑的谓词因果链）]
由 \(\Hit_{\TPH}^{\partial_L\to\Phi}=\top\)，存在 \(\kappa>0\) 与 \(\varepsilon\ge0\)，使
\[
  \sum_t\Delta_{\partial_L}\Phi_t
  \ge
  \kappa\sum_t[-\Delta_{\partial_L}\JacGap_t]_+
  -
  \varepsilon\sum_t[\Delta_{\partial_L}\JacGap_t]_+.
\]
由非退化条件，
\[
  \kappa\sum_t[-\Delta_{\partial_L}\JacGap_t]_+
  -
  \varepsilon\sum_t[\Delta_{\partial_L}\JacGap_t]_+
  >
  0.
\]
因此
\[
  \sum_t\Delta_{\partial_L}\Phi_t>0.
\]
证毕。
\end{proof}

\subsection{定理：偏微正影响累计大于偏微负影响}
\label{subsec:tex58-main-impact-theorem}

\begin{theorem}[有限回放偏微正影响严格支配定理]
\label{thm:tex58-positive-impact-dominates}
在有限回放证据域 \(\mathcal D_{\mathrm{replayid}}\) 中，若
\[
  \Eff_{\partial\JacGap}=\top,
  \qquad
  \Hit_{\TPH}^{\partial_L\to\Phi}=\top,
  \qquad
  \Eff_{\mathrm{equity}}=\top,
\]
且 \(\TPHLHTS{}\) 传导满足非退化条件
\[
  \kappa
  \sum_{t=0}^{N-1}[-\Delta_{\partial_L}\JacGap_t]_+
  >
  \varepsilon
  \sum_{t=0}^{N-1}[\Delta_{\partial_L}\JacGap_t]_+,
\]
则
\[
  \sum_{t=0}^{N-1}I^+_{\partial_L}(t)
  >
  \sum_{t=0}^{N-1}I^-_{\partial_L}(t).
\]
等价地，
\[
  \PositiveImpact(\partial_L)
  >
  \NegativeImpact(\partial_L).
\]
\end{theorem}

\begin{Certificate}[总证书：偏微下降 + TPH-LHTS 桥接 + 工程兑现]
\label{cert:tex58-positive-impact-dominates}
证书登记谓词链：
\[
\begin{aligned}
  T_1 &: \Eff_{\partial\JacGap}=\top
  \Rightarrow
  \sum_t[-\Delta_{\partial_L}\JacGap_t]_+
  >
  \sum_t[\Delta_{\partial_L}\JacGap_t]_+,\\
  T_2 &: \Hit_{\TPH}^{\partial_L\to\Phi}=\top
  \wedge\text{非退化传导}
  \Rightarrow
  \sum_t\Delta_{\partial_L}\Phi_t>0
  \quad\text{（命题~\ref{prop:tex58-bridge-positive-attribution}）},\\
  T_3 &: C_{\partial_L}(t)=\Delta_{\partial_L}\Phi_t
  \Rightarrow
  \sum_t C_{\partial_L}(t)>0,\\
  T_4 &: \sum_t C_{\partial_L}(t)>0
  \Rightarrow
  \sum_t I^+_{\partial_L}(t)>\sum_t I^-_{\partial_L}(t)
  \quad\text{（引理~\ref{lem:tex58-impact-decomposition}）},\\
  T_5 &: \Eff_{\mathrm{equity}}=\top
  \Rightarrow
  \Phi_N-\Phi_0>0
  \quad\text{（定义~\ref{def:tex58-equity-effectiveness}）}.
\end{aligned}
\]
\end{Certificate}

\begin{proof}[证明（基于数学符号推演逻辑的谓词因果链）]
由 \(\Eff_{\partial\JacGap}=\top\)，得到
\[
  \sum_t[-\Delta_{\partial_L}\JacGap_t]_+
  >
  \sum_t[\Delta_{\partial_L}\JacGap_t]_+.
\]
这登记偏微/微分下降在有限回放区间内累计有效。

由 \(\Hit_{\TPH}^{\partial_L\to\Phi}=\top\) 与非退化传导条件，命题
~\ref{prop:tex58-bridge-positive-attribution} 给出
\[
  \sum_t\Delta_{\partial_L}\Phi_t>0.
\]
又由定义~\ref{def:tex58-partial-contribution}，
\[
  C_{\partial_L}(t)=\Delta_{\partial_L}\Phi_t,
\]
所以
\[
  \sum_t C_{\partial_L}(t)>0.
\]
由引理~\ref{lem:tex58-impact-decomposition}，
\[
  \sum_t C_{\partial_L}(t)
  =
  \sum_t I^+_{\partial_L}(t)
  -
  \sum_t I^-_{\partial_L}(t).
\]
因此
\[
  \sum_t I^+_{\partial_L}(t)
  -
  \sum_t I^-_{\partial_L}(t)>0,
\]
移项得到
\[
  \sum_t I^+_{\partial_L}(t)
  >
  \sum_t I^-_{\partial_L}(t).
\]
同时，\(\Eff_{\mathrm{equity}}=\top\) 由定义登记
\[
  \Phi_N-\Phi_0>0,
\]
说明该偏微归因结论发生在整体工程目标已经正向兑现的证据域内，而不是孤立的局部下降记录。
证毕。
\end{proof}

\subsection{推论：三类可登记结论}
\label{subsec:tex58-impact-corollaries}

\begin{corollary}[整体工程兑现成立]
\label{cor:tex58-equity-realized}
若
\[
  \Eff_{\mathrm{equity}}=\top,
\]
则
\[
  \mathcal D_{\mathrm{replayid}}\models \Phi_N-\Phi_0>0.
\]
特别地，当 \(\Phi\) 的主要分量取净值与余额时，可登记
\[
  \Eq_N>\Eq_0,
  \qquad
  \Bal_N>\Bal_0,
\]
作为弱化工程兑现证据。
\end{corollary}

\begin{Certificate}[证书：兑现谓词展开]
\label{cert:tex58-equity-realized}
证书登记为
\[
  \Eff_{\mathrm{equity}}=\top
  \Longleftrightarrow
  \Phi_N-\Phi_0>0.
\]
\end{Certificate}

\begin{proof}[证明（基于数学符号推演逻辑的谓词因果链）]
该结论是定义~\ref{def:tex58-equity-effectiveness} 的直接展开。证毕。
\end{proof}

\begin{corollary}[偏微/微分下降成立]
\label{cor:tex58-partial-descent-realized}
若
\[
  \Eff_{\partial\JacGap}=\top,
\]
则
\[
  \sum_{t=0}^{N-1}\Delta_{\partial_L}\JacGap_t<0.
\]
因此 \(\JacGap\) 在 \(\partial_L\) 作用下的有限区间累计趋势为下降。
\end{corollary}

\begin{Certificate}[证书：累计降障大于累计增障]
\label{cert:tex58-partial-descent-realized}
证书登记为
\[
  \Eff_{\partial\JacGap}=\top
  \Longleftrightarrow
  \sum_tD^+_{\partial_L}(t)>\sum_tD^-_{\partial_L}(t)
  \Longleftrightarrow
  \sum_t\Delta_{\partial_L}\JacGap_t<0.
\]
\end{Certificate}

\begin{proof}[证明（基于数学符号推演逻辑的谓词因果链）]
由定义~\ref{def:tex58-positive-negative-jacgap} 与引理~\ref{lem:tex58-partial-jacgap-equivalence}，
命题立即成立。证毕。
\end{proof}

\begin{corollary}[偏微正影响累计大于负影响]
\label{cor:tex58-positive-impact-realized}
在定理~\ref{thm:tex58-positive-impact-dominates} 的条件下，
\[
  \mathcal D_{\mathrm{replayid}}
  \models
  \left(
  \Eff_{\partial\JacGap}
  \wedge
  \Hit_{\TPH}^{\partial_L\to\Phi}
  \wedge
  \Eff_{\mathrm{equity}}
  \right)
  \Rightarrow
  \PositiveImpact(\partial_L)>\NegativeImpact(\partial_L).
\]
\end{corollary}

\begin{Certificate}[证书：总定理实例化]
\label{cert:tex58-positive-impact-realized}
证书登记为
\[
  \PositiveImpact(\partial_L)
  \DefEq
  \sum_t I^+_{\partial_L}(t),
  \qquad
  \NegativeImpact(\partial_L)
  \DefEq
  \sum_t I^-_{\partial_L}(t),
\]
并调用定理~\ref{thm:tex58-positive-impact-dominates}。
\end{Certificate}

\begin{proof}[证明（基于数学符号推演逻辑的谓词因果链）]
由定理~\ref{thm:tex58-positive-impact-dominates}，
\[
  \sum_t I^+_{\partial_L}(t)
  >
  \sum_t I^-_{\partial_L}(t).
\]
将左右两端分别替换为 \(\PositiveImpact(\partial_L)\) 与
\(\NegativeImpact(\partial_L)\)，即得推论。证毕。
\end{proof}

\subsection{备注：有限证据域边界与最终口径}
\label{subsec:tex58-impact-remarks}

\begin{remark}[有限证据域边界]
本章证明的严格范围是有限回放证据域
\[
  \mathcal D_{\mathrm{replayid}}.
\]
它证明的是
\[
  \text{有限区间累计工程效应为正}.
\]
它不证明逐 tick 点态正贡献：
\[
  \forall t,\quad I^+_{\partial_L}(t)>I^-_{\partial_L}(t).
\]
它也不证明所有未来运行都成立：
\[
  \forall\mathcal D,\quad
  \sum_t I^+_{\partial_L}(t)>\sum_t I^-_{\partial_L}(t).
\]
它更不证明 \(\partial_L\) 在无限候选集上达到全局最优。
\end{remark}

\begin{remark}[最终严格结论]
本章最终结论为：
\[
  \boxed{
  \begin{gathered}
  \text{在有限回放证据域 }\mathcal D_{\mathrm{replayid}}\text{ 内，}\\
  \text{若整体工程兑现成立、\TPHLHTS{} 命中传导成立、偏微/微分下降成立，}\\
  \text{且传导非退化，则偏微正影响累计量严格大于偏微负影响累计量。}
  \end{gathered}
  }
\]
形式化写作即
\[
  \boxed{
  \mathcal D_{\mathrm{replayid}}
  \models
  \left(
  \Eff_{\partial\JacGap}
  \wedge
  \Hit_{\TPH}^{\partial_L\to\Phi}
  \wedge
  \Eff_{\mathrm{equity}}
  \right)
  \Rightarrow
  \PositiveImpact(\partial_L)
  >
  \NegativeImpact(\partial_L).
  }
\]
\end{remark}

%%%%%%%%%%%%%%%%%%%%%%%%%%%%%%%%%%%%%%%%%%%%%%%%%%%%%%%%%%%%
\clearpage

\section{形式逻辑：微观隧穿作为 PDE 有效下降路径搜索的工程替代}
\label{sec:tex58-microtunnel-pde-descent-substitution}
%%%%%%%%%%%%%%%%%%%%%%%%%%%%%%%%%%%%%%%%%%%%%%%%%%%%%%%%%%%%

\begin{remark}[核心陈述（严格化后）]
本章证明的命题不是
\[
  \text{微观隧穿/同伦命中用粗糙离散差分替代 PDE 理论整体}.
\]
严格结论是：
\[
  \boxed{
  \text{微观隧穿/同伦命中用统计力学路径族、同伦修复与 JacGap 证书，}
  }
\]
\[
  \boxed{
  \text{替代 PDE 工程求解中的有效下降路径搜索任务。}
  }
\]
也就是说，它替代的不是偏微分方程的存在性、唯一性、正则性或弱解理论，而是工程求解中最关键的搜索环节：
\[
  \text{寻找使残差下降、目标兑现、风险受控的变分路径}.
\]
形式化目标写成
\[
  \boxed{
  \mathcal D_{\mathrm{replayid}}
  \models
  \left(
  \Eff_{\partial\JacGap}
  \wedge
  \Hit_{\TPH}^{\partial_L\to\Phi}
  \wedge
  \Eff_{\mathrm{equity}}
  \wedge
  \Fid_{\JacGap\to\PDE}
  \right)
  \Rightarrow
  \Subst_{\PDE}^{\mathrm{eng}}(\Tunnel_G)=\top.
  }
\]
\end{remark}

\begin{remark}[阅读导航与引用口径]
本章承接第~\ref{sec:tex58-pde-residual-descent-gain} 章关于 PDE 残差能量与增益兑现的基础证明，
以及第~\ref{sec:tex58-tph-lhts-finite-replay-impact} 章关于 \(\mathcal D_{\mathrm{replayid}}\) 内
\TPHLHTS{} 传导和偏微正影响支配的证据链。仓内背景方面，统计力学同伦命中对高阶 PDE
求解范式的替代、微观隧穿可控化、路径空间命中与保真性层级参考 \cite{RepoTex55PDEParadigm}；
\LHTS{} 的格点同伦扭曲、延迟信用和命中证书参考 \cite{RepoTex56LHTS}；
\TPH{} 的轨迹函数参数同伦、概率路径同伦类、JacGap 证书化与 D 结构闭合参考
\cite{RepoTex57TPHLHTS}；Law-Space 路径积分与 JacGap 背景参考
\cite{RepoAppMathVol1,RepoEconomicsVol}。

外部数学背景方面，普通有限差分法的标准证明链是“一致性、稳定性、收敛性”，可参考
\cite{LeVeque2007FiniteDifference,Strikwerda2004FiniteDifference}；路径积分与统计力学路径权重
可参考 \cite{FeynmanHibbs1965PathIntegrals,FreidlinWentzell2012RandomPerturbations}；
梯度流与变分下降背景仍参考 \cite{AmbrosioGigliSavare2008GradientFlows,NocedalWright2006NumericalOptimization}。
\end{remark}

\subsection{层级定义：L0、L1、partial-L 与工程替代谓词}
\label{subsec:tex58-layers-and-substitution}

\begin{definition}[Law-Space 结构与裸离散差异层]
\label{def:tex58-lawspace-l0}
设
\[
  \mathfrak L
\]
为给定的 Law-Space、语义结构或工程约束结构。定义 \(L_0\) 为裸离散差异层。
其对象是相邻记录之间的差异：
\[
  \Delta_0 x_i
  \DefEq
  x_{i+1}-x_i.
\]
在没有 \(\mathfrak L\) 解释之前，\(\Delta_0 x_i\) 只登记两个离散记录之间的数值差异，
不自动登记方向、因果、约束、连续性、可积性、闭合性、可修复性或可审计性。
\end{definition}

\begin{definition}[\(L_1\) 层、Law-Space 偏位与语义解释映射]
\label{def:tex58-lawspace-l1}
定义 \(L_1\) 为在 \(\mathfrak L\) 结构下具有可解释连续性的最小偏位层。
设
\[
  \partial_L:\text{状态空间}\longrightarrow\text{有意义的最小变化方向}.
\]
形式上，
\[
  \partial_Lu
  \DefEq
  \lim_{h\to0}^{\mathfrak L}
  \frac{u(x+h)-u(x)}{h}.
\]
这里的 \(\lim^{\mathfrak L}\) 不是裸极限，而是在给定拓扑、度量、联络、坐标、约束与语义解释之后的极限。
定义语义解释映射
\[
  \sigma_{\mathfrak L}:L_0\to L_1,
\]
并记可用差分为
\[
  \Delta_{\mathrm{usable}}
  \DefEq
  \sigma_{\mathfrak L}(\Delta_0).
\]
\end{definition}

\begin{definition}[PDE 下降路径搜索任务]
\label{def:tex58-pde-descent-search}
令
\[
  E_{\PDE}(u)
  \DefEq
  \frac12\norm{\Pcal[u]}^2
  +
  \frac{\lambda}{2}\norm{\Bcal[u]}^2,
  \qquad \lambda>0.
\]
定义 PDE 工程求解中的下降路径搜索任务为
\[
  \Solve_{\PDE}^{\mathrm{descent}}
  \DefEq
  \left\{
  \text{寻找路径 }u(\tau)\text{ 使 }
  E_{\PDE}\downarrow,\ \Phi\uparrow,\ \Risk\text{ 受控}
  \right\}.
\]
该对象不表示 PDE 理论整体，而只表示工程求解中的有效下降路径搜索模块。
\end{definition}

\begin{definition}[G-Framework 微观隧穿算子与工程替代谓词]
\label{def:tex58-tunnel-substitution}
令
\[
  \Tunnel_G
\]
表示 \GFramework{} 中的微观隧穿/同伦命中算子。定义
\[
  \Subst_{\PDE}^{\mathrm{eng}}(\Tunnel_G)=\top
\]
当且仅当在给定有限证据域内，\(\Tunnel_G\) 能以可审计方式完成
\(\Solve_{\PDE}^{\mathrm{descent}}\) 所要求的有效下降路径搜索，即
\[
  \Tunnel_G
  \succcurlyeq_{\mathrm{eng}}
  \Solve_{\PDE}^{\mathrm{descent}}.
\]
\end{definition}

\subsection{公理（降级为定理）：有效替代证据的归纳闭合}
\label{subsec:tex58-substitution-induction}

\begin{theorem}[公理降级：有限回放替代证据累计归纳公式]
\label{thm:tex58-substitution-induction}
给定有限回放区间 \(0,\ldots,N-1\)。定义单步替代证据得分
\[
  q_t
  \DefEq
  a\,[-\Delta_{\partial_L}E_{\PDE,t}]_+
  +
  b\,[-\Delta_{\partial_L}\JacGap_t]_+
  +
  c\,[-\Delta_{\partial_L}\Risk_t]_+
  +
  d\,[\Delta_{\partial_L}\Phi_t]_+,
\]
其中 \(a,b,c,d>0\)。定义累计替代证据
\[
  Q_n
  \DefEq
  \sum_{t=0}^{n}q_t,
  \qquad 0\le n\le N-1.
\]
则
\[
  Q_n=\sum_{t=0}^{n}q_t
\]
对所有 \(0\le n\le N-1\) 成立。若存在阈值 \(\theta_{\mathrm{eng}}>0\) 使
\[
  Q_{N-1}\ge\theta_{\mathrm{eng}},
\]
则有限回放证据支持
\[
  \Tunnel_G\succcurlyeq_{\mathrm{eng}}\Solve_{\PDE}^{\mathrm{descent}}.
\]
\end{theorem}

\begin{Certificate}[数学归纳法证书：单步证据到累计证据]
\label{cert:tex58-substitution-induction}
定义谓词
\[
  P(n):\ Q_n=\sum_{t=0}^{n}q_t.
\]
归纳证书登记为
\[
  \pi_{\mathrm{subst}}
  =
  \bigl(P(0),\ \forall n\in\{0,\ldots,N-2\},\ P(n)\Rightarrow P(n+1)\bigr).
\]
唯一使用的递推关系是
\[
  Q_{n+1}=Q_n+q_{n+1}.
\]
\end{Certificate}

\begin{proof}[证明（数学归纳法证书；基于数学符号推演逻辑的谓词因果链）]
\textbf{基步：}
由定义，
\[
  Q_0=q_0=\sum_{t=0}^{0}q_t.
\]
故 \(P(0)\) 成立。

\textbf{归纳步：}
假设 \(P(n)\) 成立，即
\[
  Q_n=\sum_{t=0}^{n}q_t.
\]
由累计证据递推，
\[
  Q_{n+1}=Q_n+q_{n+1}.
\]
代入归纳假设得
\[
  Q_{n+1}
  =
  \sum_{t=0}^{n}q_t+q_{n+1}
  =
  \sum_{t=0}^{n+1}q_t.
\]
故 \(P(n)\Rightarrow P(n+1)\)。

由数学归纳法，\(P(n)\) 对所有 \(0\le n\le N-1\) 成立。
若
\[
  Q_{N-1}\ge\theta_{\mathrm{eng}},
\]
则有限区间累计证据达到工程替代阈值；按定义~\ref{def:tex58-tunnel-substitution}，
该证据支持
\[
  \Tunnel_G\succcurlyeq_{\mathrm{eng}}\Solve_{\PDE}^{\mathrm{descent}}.
\]
证毕。
\end{proof}

\subsection{定理：微分/偏微分是 L1 层的最小有意义差分}
\label{subsec:tex58-l1-minimal-meaningful-difference}

\begin{theorem}[\(L_1\) 最小有意义偏位定理]
\label{thm:tex58-l1-minimal-meaningful}
在给定 Law-Space 结构 \(\mathfrak L\) 下，若一个变化单元要同时满足
\[
  \Meaningful
  \wedge
  \Directional
  \wedge
  \Composable
  \wedge
  \Auditable
  \wedge
  \Integrable,
\]
则它的最小层级不是裸离散层 \(L_0\)，而是偏微/微分层 \(L_1\)。形式化地，
\[
  \boxed{
  \partial_L
  =
  \min_{\mathfrak L}
  \left\{
  \text{有意义、可组合、可积分、可审计的变化单元}
  \right\}.
  }
\]
\end{theorem}

\begin{Certificate}[证书：裸差分不可单独签发语义性质]
\label{cert:tex58-l1-minimal-meaningful}
证书登记谓词链：
\[
\begin{array}{rl}
  L_1: & \Delta_0u_i=u_{i+1}-u_i\in L_0,\\[2pt]
  L_2: & L_0\not\models
  \begin{gathered}[t]
    \Meaningful(\Delta_0u_i)\wedge\Directional(\Delta_0u_i)\\
    {}\wedge\Composable(\Delta_0u_i)\wedge\Auditable(\Delta_0u_i)\\
    {}\wedge\Integrable(\Delta_0u_i),
  \end{gathered}\\[8pt]
  L_3: & \sigma_{\mathfrak L}:L_0\to L_1,\\[2pt]
  L_4: & \sigma_{\mathfrak L}(\Delta_0u_i)
  \models
  \begin{gathered}[t]
    \Meaningful\wedge\Directional\wedge\Composable\\
    {}\wedge\Auditable\wedge\Integrable,
  \end{gathered}\\[8pt]
  L_5: & L_1\wedge L_2\wedge L_3\wedge L_4
  \Rightarrow
  \partial_L\text{ 是 }\mathfrak L\text{ 下的最小有意义偏位}.
\end{array}
\]
\end{Certificate}

\begin{proof}[证明（基于数学符号推演逻辑的谓词因果链）]
裸差分
\[
  \Delta_0u_i=u_{i+1}-u_i
\]
只说明两个记录之间存在差异。没有 \(\mathfrak L\) 结构时，它不能单独推出
\[
  \text{有效下降方向},
  \qquad
  \text{可积分路径},
  \qquad
  \text{边界条件保持},
  \qquad
  \text{结构闭合},
  \qquad
  \text{工程增益传导}.
\]
故 \(L_0\not\models\Meaningful(\Delta_0u_i)\) 及其相关可审计性质。

若要使该差异成为可用于求解和下降的变化量，必须施加语义解释映射
\[
  \sigma_{\mathfrak L}:L_0\to L_1.
\]
于是实际被使用的对象是
\[
  \sigma_{\mathfrak L}(\Delta_0u_i),
\]
而不是裸差分 \(\Delta_0u_i\) 本身。该对象携带方向、约束、可积性和审计接口，
因此位于 \(L_1\) 层。故可用于求解、下降、优化与兑现的差分，本质上是 \(L_1\)
语义结构解释后的差分。

由定义~\ref{def:tex58-lawspace-l1}，
\[
  \partial_Lu
  =
  \lim_{h\to0}^{\mathfrak L}
  \frac{u(x+h)-u(x)}{h}
\]
正是该层级中的最小有意义偏位。证毕。
\end{proof}

\subsection{定理：继续裸差分会坠入 L0 离散障碍}
\label{subsec:tex58-l0-obstruction}

\begin{theorem}[\(L_1\to L_0\) 离散障碍定理]
\label{thm:tex58-l0-obstruction}
若对 \(L_1\) 层的有意义偏位继续做裸差分，则该操作不会自动得到更精细的偏微结构。
形式化地，
\[
  \Delta_0(\partial_Lu)
  \not\models
  \partial_L^2u
\]
除非存在额外语义修复
\[
  \mathcal R_{\mathfrak L}:L_0\to L_1
  \quad\text{或}\quad
  \mathcal R_{\mathfrak L}:L_0\to L_{\ge1}.
\]
无修复时，存在离散障碍
\[
  \Obs_{L_0}\ne0.
\]
\end{theorem}

\begin{Certificate}[证书：二阶结构需要可比较性]
\label{cert:tex58-l0-obstruction}
证书登记谓词链：
\[
\begin{aligned}
  O_1 &: \partial_Lu(x)\in T_x^{\mathfrak L},\quad
         \partial_Lu(x+h)\in T_{x+h}^{\mathfrak L},\\
  O_2 &: T_x^{\mathfrak L}\text{ 与 }T_{x+h}^{\mathfrak L}\text{ 未给出联络/坐标/语义兼容时不可直接相减},\\
  O_3 &: O_2\Rightarrow \Delta_0(\partial_Lu)\not\models\partial_L^2u,\\
  O_4 &: \mathcal R_{\mathfrak L}(\Obs_{L_0})\to0
         \text{ 或 }\JacGap(\mathcal R_{\mathfrak L}(\Delta_0(\partial_Lu)))\downarrow
         \Rightarrow \text{二阶结构可恢复}.
\end{aligned}
\]
\end{Certificate}

\begin{proof}[证明（基于数学符号推演逻辑的谓词因果链）]
二阶偏微分 \(\partial_L^2u\) 不只是两个一阶偏位记录的裸相减。
要比较
\[
  \partial_Lu(x)
  \quad\text{与}\quad
  \partial_Lu(x+h),
\]
必须保证二者处在可比较结构中。几何语言中这需要联络；数值语言中这需要网格一致性；
语义语言中这需要上下文一致性；\GFramework{} 语言中这需要同伦修复与 JacGap 审计。

若没有这些结构，则
\[
  \partial_Lu(x+h)-\partial_Lu(x)
\]
只是两个局部偏位记录之间的差异，不能直接解释为 \(\partial_L^2u\)。
因此
\[
  \Delta_0(\partial_Lu)\not\models\partial_L^2u.
\]
此时记离散障碍为 \(\Obs_{L_0}\ne0\)。
若引入 \(\mathcal R_{\mathfrak L}\) 并使
\[
  \mathcal R_{\mathfrak L}(\Obs_{L_0})\to0
\]
或使修复后的 JacGap 下降，则该裸离散障碍可以被重新提升到 \(L_1\) 或更高层结构。
证毕。
\end{proof}

\subsection{引理：PDE 工程求解的核心是有效下降路径搜索}
\label{subsec:tex58-pde-descent-search-lemma}

\begin{lemma}[PDE 工程核心的下降路径搜索引理]
\label{lem:tex58-pde-core-descent-search}
在残差能量
\[
  E_{\PDE}(u)
  =
  \frac12\norm{\Pcal[u]}^2
  +
  \frac{\lambda}{2}\norm{\Bcal[u]}^2
\]
忠实的条件下，PDE 工程求解的核心可重写为
\[
  E_{\PDE}\downarrow,
  \qquad
  \Phi\uparrow,
  \qquad
  \Risk\text{ 受控}.
\]
也就是说，工程层的关键不是形式上写出 \(\Pcal[u]=0\)，而是找到使残差下降、目标兑现且风险受控的有效变分路径。
\end{lemma}

\begin{Certificate}[证书：残差能量等价 + 工程目标附加]
\label{cert:tex58-pde-core-descent-search}
证书登记谓词链：
\[
\begin{aligned}
  P_1 &: E_{\PDE}(u)=0
  \Longleftrightarrow
  \Pcal[u]=0\wedge\Bcal[u]=0
  \quad\text{（定理~\ref{thm:tex58-solution-energy-equivalence} 的同型实例）},\\
  P_2 &: \min_uE_{\PDE}(u)\Rightarrow \text{低残差求解目标},\\
  P_3 &: \Phi\uparrow\wedge\Risk\text{ 受控}\Rightarrow \text{工程可用性},\\
  P_4 &: P_1\wedge P_2\wedge P_3
  \Rightarrow \Solve_{\PDE}^{\mathrm{descent}}\text{ 是 PDE 工程求解核心模块}.
\end{aligned}
\]
\end{Certificate}

\begin{proof}[证明（基于数学符号推演逻辑的谓词因果链）]
由第~\ref{sec:tex58-pde-residual-descent-gain} 章的残差能量等价结论，
\[
  E_{\PDE}(u)=0
  \Longleftrightarrow
  \Pcal[u]=0\wedge\Bcal[u]=0.
\]
因此严格求解可写成
\[
  \min_uE_{\PDE}(u)
\]
并以 \(E_{\PDE}\downarrow\) 作为求解过程的下降目标。
在工程系统中，仅有低残差还不足以签发可用性；还需
\[
  \Phi\uparrow,
  \qquad
  \Risk\text{ 受控}.
\]
三者合取正是定义~\ref{def:tex58-pde-descent-search} 中的
\(\Solve_{\PDE}^{\mathrm{descent}}\)。证毕。
\end{proof}

\subsection{命题：微观隧穿是统计力学路径空间上的变分搜索算子}
\label{subsec:tex58-tunnel-path-action-proposition}

\begin{proposition}[统计力学路径作用量命题]
\label{prop:tex58-path-action}
令 \(\Gamma\) 为候选路径族
\[
  \Gamma=\{\gamma:[0,1]\to\mathcal S\}.
\]
定义路径作用量
\[
  \mathcal A(\gamma)
  =
  \int_0^1
  \left(
  aE_{\PDE}(\gamma(s))
  +
  b\JacGap(\gamma(s))
  +
  c\Risk(\gamma(s))
  -
  d\Phi(\gamma(s))
  \right)ds,
  \qquad a,b,c,d>0.
\]
定义统计力学路径权重
\[
  \mathbb P_\beta(\gamma)
  =
  \frac{\exp(-\beta\mathcal A(\gamma))}{Z_\beta},
  \qquad
  Z_\beta=\int_{\Gamma}\exp(-\beta\mathcal A(\gamma))\,d\mu(\gamma).
\]
则低作用量路径倾向于同时满足
\[
  E_{\PDE}\downarrow,
  \qquad
  \JacGap\downarrow,
  \qquad
  \Risk\downarrow,
  \qquad
  \Phi\uparrow.
\]
\end{proposition}

\begin{Certificate}[证书：正权重惩罚项 + 负号收益项]
\label{cert:tex58-path-action}
证书登记谓词链：
\[
\begin{aligned}
  A_1 &: a,b,c,d>0,\\
  A_2 &: \mathcal A(\gamma)
  =
  \int(aE_{\PDE}+b\JacGap+c\Risk-d\Phi)\,ds,\\
  A_3 &: A_1\wedge A_2
  \Rightarrow
  \mathcal A\downarrow\text{ 奖励 }E_{\PDE},\JacGap,\Risk\text{ 下降并奖励 }\Phi\text{ 上升},\\
  A_4 &: \mathbb P_\beta(\gamma)\propto e^{-\beta\mathcal A(\gamma)}
  \Rightarrow
  \beta\text{ 增大时低作用量路径权重增强}.
\end{aligned}
\]
\end{Certificate}

\begin{proof}[证明（基于数学符号推演逻辑的谓词因果链）]
由 \(a,b,c>0\)，路径作用量中
\[
  aE_{\PDE}+b\JacGap+c\Risk
\]
是惩罚项。降低 \(E_{\PDE}\)、降低 \(\JacGap\) 与降低 \(\Risk\) 都会降低 \(\mathcal A\)。
由 \(-d\Phi\) 且 \(d>0\)，提高 \(\Phi\) 也会降低 \(\mathcal A\)。
因此低作用量路径同时偏向
\[
  E_{\PDE}\downarrow,\quad
  \JacGap\downarrow,\quad
  \Risk\downarrow,\quad
  \Phi\uparrow.
\]
再由
\[
  \mathbb P_\beta(\gamma)
  =
  Z_\beta^{-1}e^{-\beta\mathcal A(\gamma)},
\]
在 \(\beta\) 增大时，较低 \(\mathcal A\) 的路径获得更高相对权重。
所以 \(\Tunnel_G\) 可解释为统计力学路径空间上的全局变分搜索算子。证毕。
\end{proof}

\subsection{定理：微观隧穿工程替代 PDE 下降路径搜索}
\label{subsec:tex58-tunnel-substitution-theorem}

\begin{theorem}[微观隧穿工程替代定理]
\label{thm:tex58-tunnel-engineering-substitution}
若在有限证据域 \(\mathcal D_{\mathrm{replayid}}\) 中满足
\[
  \Fid_{\JacGap\to\PDE}=\top,
  \qquad
  \Hit_{\TPH}^{\partial_L\to\Phi}=\top,
  \qquad
  \Eff_{\partial\JacGap}=\top,
  \qquad
  \Eff_{\mathrm{equity}}=\top,
\]
并且 \(\Fid_{\JacGap\to\PDE}=\top\) 表示
\[
  \JacGap\downarrow
  \Rightarrow
  E_{\PDE}\downarrow
  \quad\text{或至少}\quad
  \JacGap\downarrow
  \Rightarrow
  E_{\mathrm{struct}}\downarrow,
\]
则
\[
  \Tunnel_G
  \succcurlyeq_{\mathrm{eng}}
  \Solve_{\PDE}^{\mathrm{descent}}.
\]
等价地，
\[
  \mathcal D_{\mathrm{replayid}}
  \models
  \Subst_{\PDE}^{\mathrm{eng}}(\Tunnel_G)=\top.
\]
\end{theorem}

\begin{Certificate}[证书：JacGap 忠实代理 + TPH 传导 + 工程兑现]
\label{cert:tex58-tunnel-engineering-substitution}
证书登记谓词链：
\[
\begin{aligned}
  U_1 &: \Fid_{\JacGap\to\PDE}=\top
  \Rightarrow
  \JacGap\downarrow\Rightarrow E_{\PDE}\downarrow
  \text{ 或 }E_{\mathrm{struct}}\downarrow,\\
  U_2 &: \Eff_{\partial\JacGap}=\top
  \Rightarrow
  \sum_t[-\Delta_{\partial_L}\JacGap_t]_+
  >
  \sum_t[\Delta_{\partial_L}\JacGap_t]_+
  \quad\text{（推论~\ref{cor:tex58-partial-descent-realized}）},\\
  U_3 &: U_1\wedge U_2\Rightarrow E_{\PDE}\text{ 或结构残差在累计意义下降},\\
  U_4 &: \Hit_{\TPH}^{\partial_L\to\Phi}=\top
  \wedge\Eff_{\mathrm{equity}}=\top
  \Rightarrow \Phi\text{ 工程兑现并由 }\partial_L\text{ 传导支持},\\
  U_5 &: \text{命题~\ref{prop:tex58-path-action}}\Rightarrow
  \Tunnel_G\text{ 搜索低作用量路径},\\
  U_6 &: U_3\wedge U_4\wedge U_5
  \Rightarrow
  \Tunnel_G\succcurlyeq_{\mathrm{eng}}\Solve_{\PDE}^{\mathrm{descent}}.
\end{aligned}
\]
\end{Certificate}

\begin{proof}[证明（基于数学符号推演逻辑的谓词因果链）]
PDE 下降路径搜索任务要求
\[
  E_{\PDE}\downarrow,
  \qquad
  \Phi\uparrow,
  \qquad
  \Risk\text{ 受控}.
\]
由 \(\Eff_{\partial\JacGap}=\top\)，有限回放区间内 JacGap 累计下降成立。
由 \(\Fid_{\JacGap\to\PDE}=\top\)，JacGap 累计下降可忠实代理为 PDE 残差或结构残差下降。
因此 \(\Tunnel_G\) 生成的下降证据覆盖了 \(\Solve_{\PDE}^{\mathrm{descent}}\) 的残差下降要求。

由 \(\Hit_{\TPH}^{\partial_L\to\Phi}=\top\)，偏微/微分下降经 \TPHLHTS{} 命中传导到工程目标 \(\Phi\)。
由 \(\Eff_{\mathrm{equity}}=\top\)，
\[
  \Phi_N-\Phi_0>0,
\]
整体工程兑现成立。由第~\ref{sec:tex58-tph-lhts-finite-replay-impact} 章的主定理，还可登记
\[
  \PositiveImpact(\partial_L)>\NegativeImpact(\partial_L).
\]
这给出目标兑现与风险损耗被累计正影响支配的工程证书。

由命题~\ref{prop:tex58-path-action}，\(\Tunnel_G\) 不是裸差分过程，而是在路径空间中搜索低作用量路径；
低作用量路径同时压低 \(E_{\PDE}\)、\(\JacGap\)、\(\Risk\)，并提高 \(\Phi\)。
因此它完成的是 PDE 工程求解中有效下降路径搜索模块，而不是替代 PDE 理论整体。
故
\[
  \Tunnel_G\succcurlyeq_{\mathrm{eng}}\Solve_{\PDE}^{\mathrm{descent}},
\]
即
\[
  \mathcal D_{\mathrm{replayid}}\models\Subst_{\PDE}^{\mathrm{eng}}(\Tunnel_G)=\top.
\]
证毕。
\end{proof}

\subsection{推论：可用差分属于 L1 语义解释层}
\label{subsec:tex58-usable-difference-corollary}

\begin{corollary}[可用差分的 \(L_1\) 解释推论]
\label{cor:tex58-usable-difference-l1}
可用于 PDE 求解、同伦命中、残差下降与工程兑现的差分不是无附加语义的 \(L_0\) 裸差分，
而是
\[
  \Delta_{\mathrm{usable}}
  =
  \sigma_{\mathfrak L}(\Delta_0)
  \in L_1.
\]
因此
\[
  \Delta_0\in L_0,
  \qquad
  \Delta_{\mathrm{usable}}\in L_1.
\]
\end{corollary}

\begin{Certificate}[证书：裸差分经 Law-Space 解释后才可用]
\label{cert:tex58-usable-difference-l1}
证书登记谓词链：
\[
\begin{aligned}
  R_1 &: \Delta_0\in L_0,\\
  R_2 &: L_0\not\models \Usable(\Delta_0),\\
  R_3 &: \sigma_{\mathfrak L}:L_0\to L_1,\\
  R_4 &: \Delta_{\mathrm{usable}}=\sigma_{\mathfrak L}(\Delta_0),\\
  R_5 &: R_1\wedge R_2\wedge R_3\wedge R_4
  \Rightarrow
  \Delta_{\mathrm{usable}}\in L_1.
\end{aligned}
\]
\end{Certificate}

\begin{proof}[证明（基于数学符号推演逻辑的谓词因果链）]
由定义~\ref{def:tex58-lawspace-l0}，\(\Delta_0\) 只是裸离散差异。
由定理~\ref{thm:tex58-l1-minimal-meaningful}，裸差分本身不自动携带方向性、可积性、闭合性或审计性。
若要用于 PDE 求解、同伦命中或工程兑现，必须施加
\[
  \sigma_{\mathfrak L}:L_0\to L_1.
\]
因此实际可用对象为
\[
  \Delta_{\mathrm{usable}}
  =
  \sigma_{\mathfrak L}(\Delta_0)
  \in L_1.
\]
证毕。
\end{proof}

\subsection{备注：替代范围、非替代范围与最终口径}
\label{subsec:tex58-tunnel-substitution-remarks}

\begin{remark}[这不是普通离散差分替代 PDE]
普通有限差分路径可概括为
\[
  \Pcal[u]=0
  \quad\leadsto\quad
  \Pcal_h[u_h]=0,
\]
其核心证明链是
\[
  \text{一致性}+\text{稳定性}\Rightarrow\text{收敛性}.
\]
\GFramework{} 的微观隧穿/同伦命中路径则是
\[
  \Pcal[u]=0
  \quad\leadsto\quad
  E_{\PDE}+\JacGap+\Risk-\Phi
  \quad\leadsto\quad
  \min_{\gamma\in\Gamma}\mathcal A(\gamma).
\]
其核心证明链是
\[
  \text{路径族覆盖}
  +
  \text{JacGap 下降}
  +
  \text{TPH 命中}
  +
  \text{工程兑现}
  \Rightarrow
  \text{有效下降路径成立}.
\]
所以二者不同：普通差分是把 PDE 离散化；微观隧穿是把 PDE 求解任务变分化、
路径化、统计力学化、同伦修复化与证书化。
\end{remark}

\begin{remark}[有限证据域边界]
本章证明的是有限证据域内的工程替代性强支持命题。它证明
\[
  \Tunnel_G
\]
可以工程性替代 PDE 求解中的下降路径搜索任务；它不证明 \(\Tunnel_G\) 已经替代 PDE 的全部数学理论。
因此，本章不声称所有 PDE 的解析解都可由微观隧穿直接给出，不声称所有未来运行都成立，
不声称无限候选集上的绝对全局最优已经由有限运行实现，也不声称 PDE 的存在性、唯一性和正则性理论可以取消。
\end{remark}

\begin{remark}[最终严格结论]
最终结论可压缩为：
\[
  \boxed{
  \begin{gathered}
  \text{微观隧穿/同伦命中机制，在 \GFramework{} 中构成一种}\\
  \text{基于统计力学路径族、同伦修复与 JacGap 证书的全局变分搜索算子。}
  \end{gathered}
  }
\]
\[
  \boxed{
  \begin{gathered}
  \text{在有限回放证据域内，若 JacGap 下降、TPH-LHTS 命中传导、}\\
  \text{整体工程兑现与 JacGap 对 PDE 残差的忠实代理同时成立，}\\
  \text{则该机制可以有效替代 PDE 求解中的核心下降路径搜索任务。}
  \end{gathered}
  }
\]
其替代依据不是 \(L_0\) 裸差分比微分更基础，而是：
\[
  \boxed{
  \partial_L\text{ 是 }L_1\text{ 层有意义的最小偏位；继续坠入 }L_0
  \text{ 会产生离散障碍，}
  }
\]
\[
  \boxed{
  \text{必须通过同伦修复、语义解释与 JacGap 审计重新提升为可解释、可积分、可审计的结构。}
  }
\]
\end{remark}

%%%%%%%%%%%%%%%%%%%%%%%%%%%%%%%%%%%%%%%%%%%%%%%%%%%%%%%%%%%%
\clearpage

\section{形式逻辑：PDE 的 L-积分测度化替代与解释等价}
\label{sec:tex58-l-measure-pde-substitution}
%%%%%%%%%%%%%%%%%%%%%%%%%%%%%%%%%%%%%%%%%%%%%%%%%%%%%%%%%%%%

上一章证明的是：微观隧穿/同伦命中可以在有限证据域内工程性替代 PDE 求解中的有效下降路径搜索任务。
本章把该结论提升到更强但仍然条件化的形式：\PDE{} 的局部微分表达可以在 Law-Space 中重写为
\(L\)-积分/测度表达；\PDE{} 的函数求解可以重写为 \(L\)-路径测度的低作用量命中；
二者在同一解释函子下给出相同工程语义。

弱形式、测试函数与残差能量的标准背景可参考
\cite{Evans2010PDE,BrennerScott2008FEM,Brezis2011FunctionalAnalysis}；
Radon--Nikodym 密度、测度变化与“几乎处处”等价的背景可参考
\cite{Bogachev2007MeasureTheory}；路径测度、作用量和测度空间梯度流背景可参考
\cite{AmbrosioGigliSavare2008GradientFlows,FeynmanHibbs1965PathIntegrals,FreidlinWentzell2012RandomPerturbations}。
仓内语义规范场、PDE 命中范式、\LHTS{} 与 \TPHLHTS{} 证书链分别参考
\cite{RepoTex37SemanticGauge,RepoTex55PDEParadigm,RepoTex56LHTS,RepoTex57TPHLHTS}。
本章仍只引用 \(NoteBook\) 下 \(\_pub\) 稿件和 \(docs\) 发布稿，不引用 \(NoteBook/\) 内非 \(\_pub\) 目录稿件。

本章最终登记的命题为
\[
  \boxed{
  \mathcal D_{\mathrm{replayid}}
  \models
  \Subst_{\PDE}^{\Lmeas}(\Tunnel_G)=\top
  }
\]
其中
\[
  \Subst_{\PDE}^{\Lmeas}(\Tunnel_G)=\top
  \quad\Longleftrightarrow\quad
  \ExprSubst_{\PDE}^{\Lmeas}
  \wedge
  \SolveSubst_{\PDE}^{\Lmeas}
  \wedge
  \InterpEquiv_{\PDE}^{\Lmeas}.
\]
因此，该替代不是 \(L_0\) 裸差分替代 PDE，而是
\[
  \boxed{
  \begin{gathered}
  \text{用 }L\text{-积分/测度结构替代 PDE 的局部微分表达，}\\
  \text{用 }L\text{-路径测度低作用量命中替代 PDE 的求解对象，}\\
  \text{并在同一 Law-Space 解释函子下保持严格解释等价。}
  \end{gathered}
  }
\]

\subsection{对象域与三类替代谓词}
\label{subsec:tex58-l-measure-domain-definitions}

\begin{definition}[Law-Space 测度结构]
\label{def:tex58-lawspace-measure-structure}
定义 Law-Space 测度结构为
\[
  \mathfrak L
  \DefEq
  (X,\Sigma_L,\mu_L,\nabla_L,\mathcal H_L,\Pi_L,\mathcal C_L),
\]
其中 \(X\) 是状态空间，\(\Sigma_L\) 是 \(L\)-可测结构，\(\mu_L\) 是 \(L\)-测度，
\(\nabla_L\) 是 \(L\)-联络或语义联络，\(\mathcal H_L\) 是同伦修复结构，
\(\Pi_L\) 是从路径/测度解到工程状态解的投影，\(\mathcal C_L\) 是解释函子。
\end{definition}

\begin{definition}[偏微分作为测度变化的局部密度]
\label{def:tex58-partial-as-density}
设 \(D_Lu\) 是 \(u\) 在 \(\mathfrak L\) 中诱导的 \(L\)-变化测度。若
\[
  D_Lu\ll \mu_L,
\]
则 Radon--Nikodym 型密度存在，并定义
\[
  \partial_Lu
  \DefEq
  \frac{dD_Lu}{d\mu_L}.
\]
于是对任意 \(A\in\Sigma_L\)，
\[
  D_Lu(A)
  =
  \int_A \partial_Lu\,d\mu_L.
\]
这说明 \(\partial_Lu\) 是 \(D_Lu\) 相对于 \(\mu_L\) 的局部密度；
微分/偏微分是 \(L\)-积分/测度结构的局部表达，而不是比 \(L\)-测度更原始的对象。
\end{definition}

\begin{definition}[表达替代、求解替代与解释等价]
\label{def:tex58-three-l-measure-substitutions}
定义三类替代谓词如下：
\[
  \ExprSubst_{\PDE}^{\Lmeas}=\top
\]
表示 PDE 的局部微分表达可由 \(L\)-积分/测度表达等价替代；
\[
  \SolveSubst_{\PDE}^{\Lmeas}=\top
\]
表示 PDE 的函数求解可由 \(L\)-残差能量或 \(L\)-路径测度低作用量命中替代；
\[
  \InterpEquiv_{\PDE}^{\Lmeas}=\top
\]
表示 PDE 表达与 \(L\)-积分/测度表达在同一解释函子 \(\mathcal C_L\) 下严格等价。
总替代谓词定义为
\[
  \Subst_{\PDE}^{\Lmeas}(\Tunnel_G)=\top
  \quad\Longleftrightarrow\quad
  \ExprSubst_{\PDE}^{\Lmeas}
  \wedge
  \SolveSubst_{\PDE}^{\Lmeas}
  \wedge
  \InterpEquiv_{\PDE}^{\Lmeas}.
\]
\end{definition}

\subsection{公理降级为定理：解释谓词有限闭合}
\label{subsec:tex58-l-meaning-induction}

\begin{theorem}[公理降级：解释谓词有限闭合的归纳定理]
\label{thm:tex58-l-meaning-induction}
设有限解释谓词族为
\[
  \mathcal P_L^{\mathrm{sem}}
  =
  \{p_0,p_1,\dots,p_m\},
\]
其中各 \(p_j\) 分别登记残差、边界、约束、JacGap、风险、工程兑现与审计证书等解释项。
若对每个 \(j\) 都有
\[
  p_j\!\left(\mathcal C_L^{\Lmeas}(\nu)\right)
  =
  p_j\!\left(\mathcal C_L^{\PDE}(\Pi_L\nu)\right),
\]
则
\[
  \bigwedge_{j=0}^{m}
  p_j\!\left(\mathcal C_L^{\Lmeas}(\nu)\right)
  =
  \bigwedge_{j=0}^{m}
  p_j\!\left(\mathcal C_L^{\PDE}(\Pi_L\nu)\right).
\]
\end{theorem}

\begin{Certificate}[数学归纳法证书：解释谓词逐项闭合]
\label{cert:tex58-l-meaning-induction}
定义
\[
  H_n(\nu)
  \DefEq
  \bigwedge_{j=0}^{n}
  \left[
  p_j\!\left(\mathcal C_L^{\Lmeas}(\nu)\right)
  =
  p_j\!\left(\mathcal C_L^{\PDE}(\Pi_L\nu)\right)
  \right].
\]
证书链登记为
\[
\begin{aligned}
  M_1 &: H_0(\nu)
  \Longleftrightarrow
  p_0(\mathcal C_L^{\Lmeas}(\nu))
  =
  p_0(\mathcal C_L^{\PDE}(\Pi_L\nu)),\\
  M_2 &: H_n(\nu)\wedge
  \left[
  p_{n+1}(\mathcal C_L^{\Lmeas}(\nu))
  =
  p_{n+1}(\mathcal C_L^{\PDE}(\Pi_L\nu))
  \right]
  \Rightarrow
  H_{n+1}(\nu),\\
  M_3 &: M_1\wedge M_2
  \Rightarrow
  \forall n\le m,\ H_n(\nu),\\
  M_4 &: H_m(\nu)
  \Rightarrow
  \InterpEquiv_{\PDE}^{\Lmeas}=\top
  \quad\text{在有限解释谓词族上成立}.
\end{aligned}
\]
\end{Certificate}

\begin{proof}[证明（基于数学符号推演逻辑的谓词因果链）]
归纳基 \(n=0\) 时，由假设直接得到
\[
  p_0(\mathcal C_L^{\Lmeas}(\nu))
  =
  p_0(\mathcal C_L^{\PDE}(\Pi_L\nu)),
\]
故 \(H_0(\nu)\) 成立。

归纳假设：设 \(H_n(\nu)\) 成立，即
\[
  \bigwedge_{j=0}^{n}
  \left[
  p_j(\mathcal C_L^{\Lmeas}(\nu))
  =
  p_j(\mathcal C_L^{\PDE}(\Pi_L\nu))
  \right]
\]
成立。

归纳推进：由逐项交换条件，对 \(p_{n+1}\) 有
\[
  p_{n+1}(\mathcal C_L^{\Lmeas}(\nu))
  =
  p_{n+1}(\mathcal C_L^{\PDE}(\Pi_L\nu)).
\]
因此
\[
\begin{aligned}
  H_{n+1}(\nu)
  &=
  H_n(\nu)
  \wedge
  \left[
  p_{n+1}(\mathcal C_L^{\Lmeas}(\nu))
  =
  p_{n+1}(\mathcal C_L^{\PDE}(\Pi_L\nu))
  \right]\\
  &=\top.
\end{aligned}
\]
由数学归纳法，
\[
  \forall n\le m,\quad H_n(\nu)=\top.
\]
取 \(n=m\)，有限解释谓词族上的解释严格等价成立。证毕。
\end{proof}

\subsection{定理：PDE 表达可被 L-积分测度表达替代}
\label{subsec:tex58-expression-substitution-theorem}

\begin{definition}[L-积分弱表达]
\label{def:tex58-l-weak-expression}
设 PDE 系统为
\[
  \Pcal[u]=0,
  \qquad
  \Bcal[u]=0.
\]
令 \(\mathcal T_L\) 为 \(L\)-测试函数族，\(\mathcal T_{\partial L}\) 为边界测试函数族。
定义 \(L\)-积分弱表达为
\[
  \mathcal W_L(u;\varphi,\psi)
  \DefEq
  \int_X
  \ip{\Pcal[u]}{\varphi}\,d\mu_L
  +
  \int_{\partial X}
  \ip{\Bcal[u]}{\psi}\,d\sigma_L,
\]
其中 \(\varphi\in\mathcal T_L\)，\(\psi\in\mathcal T_{\partial L}\)。
测试族满足分离性是指
\[
  \left(
  \forall \varphi\in\mathcal T_L,\ 
  \int_X\ip{f}{\varphi}\,d\mu_L=0
  \right)
  \Rightarrow
  f=0\quad L\text{-a.e.},
\]
且边界测试族对 \(\Bcal[u]\) 具有同样分离性。
\end{definition}

\begin{theorem}[PDE 表达的 L-积分替代定理]
\label{thm:tex58-expression-substitution}
在定义~\ref{def:tex58-l-weak-expression} 的分离性条件下，
\[
  \Pcal[u]=0,\quad \Bcal[u]=0
  \quad\Longleftrightarrow\quad
  \forall(\varphi,\psi)\in\mathcal T_L\times\mathcal T_{\partial L},\
  \mathcal W_L(u;\varphi,\psi)=0.
\]
因此
\[
  \ExprSubst_{\PDE}^{\Lmeas}=\top.
\]
\end{theorem}

\begin{Certificate}[证书：局部残差到积分弱表达的分离闭合]
\label{cert:tex58-expression-substitution}
证书链登记为
\[
\begin{aligned}
  E_1 &: \Pcal[u]=0\wedge \Bcal[u]=0,\\
  E_2 &: E_1\Rightarrow
  \forall(\varphi,\psi),\
  \mathcal W_L(u;\varphi,\psi)=0,\\
  E_3 &: \forall(\varphi,\psi),\
  \mathcal W_L(u;\varphi,\psi)=0,\\
  E_4 &: E_3\wedge \mathrm{Sep}(\mathcal T_L,\mathcal T_{\partial L})
  \Rightarrow
  \Pcal[u]=0\ L\text{-a.e.}\ \wedge\ \Bcal[u]=0,\\
  E_5 &: E_2\wedge E_4
  \Rightarrow
  \ExprSubst_{\PDE}^{\Lmeas}=\top.
\end{aligned}
\]
\end{Certificate}

\begin{proof}[证明（基于数学符号推演逻辑的谓词因果链）]
若
\[
  \Pcal[u]=0,
  \qquad
  \Bcal[u]=0,
\]
则对任意 \((\varphi,\psi)\) 有
\[
  \int_X\ip{\Pcal[u]}{\varphi}\,d\mu_L=0,
  \qquad
  \int_{\partial X}\ip{\Bcal[u]}{\psi}\,d\sigma_L=0.
\]
故
\[
  \mathcal W_L(u;\varphi,\psi)=0.
\]
反过来，若对所有测试对都有
\[
  \mathcal W_L(u;\varphi,\psi)=0,
\]
则分别取边界测试为零和内部测试为零，可得
\[
  \forall\varphi,\quad
  \int_X\ip{\Pcal[u]}{\varphi}\,d\mu_L=0,
\]
以及
\[
  \forall\psi,\quad
  \int_{\partial X}\ip{\Bcal[u]}{\psi}\,d\sigma_L=0.
\]
由测试族分离性，
\[
  \Pcal[u]=0\quad L\text{-a.e.},
  \qquad
  \Bcal[u]=0.
\]
所以 PDE 局部微分表达与 \(L\)-积分/测度表达等价。证毕。
\end{proof}

\subsection{定理：PDE 求解可被 L-路径测度求解替代}
\label{subsec:tex58-solve-substitution-theorem}

\begin{definition}[L-积分残差能量与 L-路径作用量]
\label{def:tex58-l-residual-energy-action}
定义 \(L\)-积分残差能量为
\[
  E_L(u)
  \DefEq
  \sup_{\|(\varphi,\psi)\|_{\mathcal T_L}\le 1}
  \left|
  \mathcal W_L(u;\varphi,\psi)
  \right|^2.
\]
定义 PDE 解集为
\[
  \Sol_{\PDE}
  \DefEq
  \{u:\Pcal[u]=0,\ \Bcal[u]=0\}.
\]
令 \(\Gamma_L\) 为 \(L\)-路径空间，\(\gamma:[0,1]\to X\)。定义 \(L\)-路径作用量
\[
  \mathcal A_L(\gamma)
  \DefEq
  \int_0^1
  \left(
  aE_L(\gamma(s))
  +
  b\JacGap(\gamma(s))
  +
  c\Risk(\gamma(s))
  -
  d\Phi(\gamma(s))
  \right)ds,
  \qquad
  a,b,c,d>0.
\]
定义路径测度
\[
  \nu_\beta(d\gamma)
  \DefEq
  \frac{e^{-\beta\mathcal A_L(\gamma)}}{Z_\beta}\,dM_L(\gamma),
  \qquad
  Z_\beta
  =
  \int_{\Gamma_L}e^{-\beta\mathcal A_L(\gamma)}\,dM_L(\gamma).
\]
\end{definition}

\begin{theorem}[PDE 求解的 L-路径测度替代定理]
\label{thm:tex58-solve-substitution}
在定理~\ref{thm:tex58-expression-substitution} 的条件下，
\[
  E_L(u)=0
  \quad\Longleftrightarrow\quad
  u\in\Sol_{\PDE}.
\]
进一步，若
\[
  \gamma^\ast\in\Argmin_{\gamma\in\Gamma_L}\mathcal A_L(\gamma),
  \qquad
  E_L(\Pi_L\gamma^\ast)=0,
\]
则
\[
  \Pi_L\gamma^\ast\in\Sol_{\PDE}.
\]
有限工程阈值版本为：若 \(E_L(\Pi_L\gamma^\ast)\le\epsilon\)，则 \(\Pi_L\gamma^\ast\)
是 \(\epsilon\)-低残差工程解。
因此
\[
  \SolveSubst_{\PDE}^{\Lmeas}=\top.
\]
\end{theorem}

\begin{Certificate}[证书：残差能量零点到路径测度低作用量命中]
\label{cert:tex58-solve-substitution}
证书链登记为
\[
\begin{aligned}
  S_1 &: E_L(u)=0
  \Longleftrightarrow
  \sup_{\|(\varphi,\psi)\|\le1}|\mathcal W_L(u;\varphi,\psi)|^2=0,\\
  S_2 &: S_1\Rightarrow
  \forall(\varphi,\psi),\ \mathcal W_L(u;\varphi,\psi)=0,\\
  S_3 &: S_2\wedge \ExprSubst_{\PDE}^{\Lmeas}
  \Rightarrow u\in\Sol_{\PDE},\\
  S_4 &: \gamma^\ast\in\Argmin\mathcal A_L
  \wedge E_L(\Pi_L\gamma^\ast)=0
  \Rightarrow
  \Pi_L\gamma^\ast\in\Sol_{\PDE},\\
  S_5 &: S_1\wedge S_2\wedge S_3\wedge S_4
  \Rightarrow
  \SolveSubst_{\PDE}^{\Lmeas}=\top.
\end{aligned}
\]
\end{Certificate}

\begin{proof}[证明（基于数学符号推演逻辑的谓词因果链）]
由 \(E_L\) 的定义，
\[
  E_L(u)=0
\]
等价于所有单位测试对上的弱残差都为零：
\[
  \forall(\varphi,\psi),\quad
  \mathcal W_L(u;\varphi,\psi)=0.
\]
由定理~\ref{thm:tex58-expression-substitution}，
\[
  \forall(\varphi,\psi),\
  \mathcal W_L(u;\varphi,\psi)=0
  \Longleftrightarrow
  \Pcal[u]=0,\ \Bcal[u]=0.
\]
故
\[
  E_L(u)=0
  \Longleftrightarrow
  u\in\Sol_{\PDE}.
\]

路径测度形式中，低温极限 \(\beta\to\infty\) 使 \(\nu_\beta\) 的质量集中到低作用量路径集合。
若命中路径 \(\gamma^\ast\) 满足
\[
  E_L(\Pi_L\gamma^\ast)=0,
\]
则由前半部分立即得到
\[
  \Pi_L\gamma^\ast\in\Sol_{\PDE}.
\]
若只满足 \(E_L(\Pi_L\gamma^\ast)\le\epsilon\)，则得到 \(\epsilon\)-低残差工程解。证毕。
\end{proof}

\subsection{引理：解释函子交换给出解释严格等价}
\label{subsec:tex58-interpretation-equivalence-lemma}

\begin{lemma}[解释函子交换引理]
\label{lem:tex58-interpretation-equivalence}
设
\[
  \Pi_L:\Sol_{\Lmeas}\to\Sol_{\PDE}
\]
为 \(L\)-测度解到 PDE 解或低残差代表元的投影。若
\[
  \mathcal C_L^{\Lmeas}(\nu)
  =
  \mathcal C_L^{\PDE}(\Pi_L\nu),
\]
则
\[
  \InterpEquiv_{\PDE}^{\Lmeas}=\top.
\]
\end{lemma}

\begin{Certificate}[证书：解释函子交换图]
\label{cert:tex58-interpretation-equivalence}
交换图证书为
\[
\begin{array}{ccc}
  \Sol_{\Lmeas} & \xrightarrow{\Pi_L} & \Sol_{\PDE}\\
  \mathcal C_L^{\Lmeas}\downarrow & & \downarrow \mathcal C_L^{\PDE}\\
  \Meaning_L & = & \Meaning_L
\end{array}
\]
等价地，
\[
\begin{aligned}
  I_1 &: \nu\in\Sol_{\Lmeas},\\
  I_2 &: \Pi_L\nu\in\Sol_{\PDE},\\
  I_3 &: \mathcal C_L^{\Lmeas}(\nu)
  =
  \mathcal C_L^{\PDE}(\Pi_L\nu),\\
  I_4 &: I_1\wedge I_2\wedge I_3
  \Rightarrow
  \InterpEquiv_{\PDE}^{\Lmeas}=\top.
\end{aligned}
\]
\end{Certificate}

\begin{proof}[证明（基于数学符号推演逻辑的谓词因果链）]
解释严格等价不是要求 PDE 符号表达与 \(L\)-测度表达外形相同，而是要求它们在同一解释函子下给出相同意义对象。
由假设，
\[
  \mathcal C_L^{\Lmeas}(\nu)
  =
  \mathcal C_L^{\PDE}(\Pi_L\nu).
\]
故任意解释谓词 \(p_j\) 满足
\[
  p_j(\mathcal C_L^{\Lmeas}(\nu))
  =
  p_j(\mathcal C_L^{\PDE}(\Pi_L\nu)).
\]
再由定理~\ref{thm:tex58-l-meaning-induction}，有限解释谓词族上的残差、边界、约束、风险、兑现和审计含义全部一致。
因此
\[
  \InterpEquiv_{\PDE}^{\Lmeas}=\top.
\]
证毕。
\end{proof}

\subsection{命题：微观隧穿是 L-测度路径空间上的全局变分算子}
\label{subsec:tex58-tunnel-l-measure-proposition}

\begin{proposition}[L-测度微观隧穿命题]
\label{prop:tex58-tunnel-l-measure-operator}
\GFramework{} 的微观隧穿/同伦命中算子可定义为
\[
  \Tunnel_G:
  \Gamma_L
  \to
  \Argmin_{\gamma\in\Gamma_L}\mathcal A_L(\gamma).
\]
在正权重 \(a,b,c,d>0\) 下，低作用量命中满足工程偏序
\[
  \mathcal A_L(\gamma)\downarrow
  \Rightarrow
  E_L\downarrow
  \wedge
  \JacGap\downarrow
  \wedge
  \Risk\downarrow
  \wedge
  \Phi\uparrow.
\]
\end{proposition}

\begin{Certificate}[证书：低作用量到低残差高兑现]
\label{cert:tex58-tunnel-l-measure-operator}
证书链登记为
\[
\begin{aligned}
  T_1 &: \mathcal A_L(\gamma)
  =
  \int_0^1(aE_L+b\JacGap+c\Risk-d\Phi)\,ds,\\
  T_2 &: a,b,c,d>0,\\
  T_3 &: \mathcal A_L(\gamma)\downarrow
  \Rightarrow
  aE_L+b\JacGap+c\Risk-d\Phi\downarrow,\\
  T_4 &: T_2\wedge T_3
  \Rightarrow
  E_L\downarrow\wedge \JacGap\downarrow
  \wedge \Risk\downarrow\wedge \Phi\uparrow,\\
  T_5 &: T_4\Rightarrow
  \Tunnel_G\text{ 是 }L\text{-测度路径空间上的全局变分算子}.
\end{aligned}
\]
\end{Certificate}

\begin{proof}[证明（基于数学符号推演逻辑的谓词因果链）]
由定义~\ref{def:tex58-l-residual-energy-action}，
\[
  \mathcal A_L(\gamma)
  =
  \int_0^1(aE_L+b\JacGap+c\Risk-d\Phi)\,ds.
\]
当 \(a,b,c,d>0\) 时，降低 \(\mathcal A_L\) 的方向优先选择低 \(E_L\)、低 \(\JacGap\)、低 \(\Risk\)
和高 \(\Phi\) 的路径。于是
\[
  \Tunnel_G:
  \Gamma_L\to\Argmin\mathcal A_L
\]
不是把 PDE 网格化后逐点差分，而是在 \(L\)-测度路径空间中命中低作用量、低残差、
低 JacGap、低风险和高兑现的路径支撑。证毕。
\end{proof}

\subsection{定理：PDE 的 L-积分测度替代总定理}
\label{subsec:tex58-total-l-measure-substitution}

\begin{theorem}[PDE 的 L-积分/测度替代总定理]
\label{thm:tex58-total-l-measure-substitution}
若
\[
  \ExprSubst_{\PDE}^{\Lmeas}=\top,
  \qquad
  \SolveSubst_{\PDE}^{\Lmeas}=\top,
  \qquad
  \InterpEquiv_{\PDE}^{\Lmeas}=\top,
\]
则
\[
  \Subst_{\PDE}^{\Lmeas}=\top.
\]
即 PDE 的局部微分表达、函数求解对象与解释语义可被 \(L\)-积分/测度结构整体替代，并保持解释严格等价。
\end{theorem}

\begin{Certificate}[证书：表达替代、求解替代、解释等价三元合取]
\label{cert:tex58-total-l-measure-substitution}
证书链登记为
\[
\begin{aligned}
  U_1 &: \ExprSubst_{\PDE}^{\Lmeas}=\top,\\
  U_2 &: \SolveSubst_{\PDE}^{\Lmeas}=\top,\\
  U_3 &: \InterpEquiv_{\PDE}^{\Lmeas}=\top,\\
  U_4 &: U_1\wedge U_2\wedge U_3
  \Longleftrightarrow
  \Subst_{\PDE}^{\Lmeas}=\top,\\
  U_5 &: \Subst_{\PDE}^{\Lmeas}=\top
  \Rightarrow
  \begin{gathered}[t]
  \text{表达替代成立，求解替代成立，}\\
  \text{解释严格等价成立}.
  \end{gathered}
\end{aligned}
\]
\end{Certificate}

\begin{proof}[证明（基于数学符号推演逻辑的谓词因果链）]
由表达替代，
\[
  \Pcal[u]=0,\ \Bcal[u]=0
  \Longleftrightarrow
  \forall(\varphi,\psi),\ \mathcal W_L(u;\varphi,\psi)=0.
\]
因此 PDE 的局部微分表达可由 \(L\)-积分弱表达替代。

由求解替代，
\[
  u\in\Sol_{\PDE}
  \Longleftrightarrow
  E_L(u)=0,
\]
并且在路径测度形式下，低作用量命中经 \(\Pi_L\) 投影得到 PDE 解或 \(\epsilon\)-低残差工程解。
因此 PDE 的求解对象可由 \(L\)-路径测度低作用量命中替代。

由解释严格等价，
\[
  \mathcal C_L^{\Lmeas}(\nu)
  =
  \mathcal C_L^{\PDE}(\Pi_L\nu).
\]
因此表达对象不同、求解对象不同，但在 Law-Space 解释函子下意义相同。
三者合取正是定义~\ref{def:tex58-three-l-measure-substitutions} 中的
\(\Subst_{\PDE}^{\Lmeas}=\top\)。证毕。
\end{proof}

\subsection{推论：有限证据域中的微观隧穿总替代证书}
\label{subsec:tex58-finite-domain-l-measure-corollary}

\begin{corollary}[有限证据域 L-测度替代推论]
\label{cor:tex58-finite-domain-l-measure-substitution}
在 \(\mathfrak L\) 的分离性、投影相容性与解释函子交换条件下，若有限回放证据域满足
\[
  \Eff_{\partial\JacGap}
  \wedge
  \Hit_{\TPH}^{\partial_L\to\Phi}
  \wedge
  \Eff_{\mathrm{equity}},
\]
则
\[
  \mathcal D_{\mathrm{replayid}}
  \models
  \Subst_{\PDE}^{\Lmeas}(\Tunnel_G)=\top.
\]
\end{corollary}

\begin{Certificate}[证书：有限证据到 L-测度总替代]
\label{cert:tex58-finite-domain-l-measure-substitution}
证书链登记为
\[
\begin{aligned}
  V_1 &: \Eff_{\partial\JacGap}=\top,\\
  V_2 &: \Hit_{\TPH}^{\partial_L\to\Phi}=\top,\\
  V_3 &: \Eff_{\mathrm{equity}}=\top,\\
  V_4 &: V_1\wedge V_2\wedge V_3
  \Rightarrow
  \sum_t I^+_{\partial_L}(t)>\sum_t I^-_{\partial_L}(t),\\
  V_5 &: \ExprSubst_{\PDE}^{\Lmeas}\wedge
  \SolveSubst_{\PDE}^{\Lmeas}\wedge
  \InterpEquiv_{\PDE}^{\Lmeas},\\
  V_6 &: V_4\wedge V_5
  \Rightarrow
  \mathcal D_{\mathrm{replayid}}
  \models
  \Subst_{\PDE}^{\Lmeas}(\Tunnel_G)=\top.
\end{aligned}
\]
\end{Certificate}

\begin{proof}[证明（基于数学符号推演逻辑的谓词因果链）]
由第~\ref{sec:tex58-tph-lhts-finite-replay-impact} 节的有限回放证据定理，
\[
  \Eff_{\partial\JacGap}
  \wedge
  \Hit_{\TPH}^{\partial_L\to\Phi}
  \wedge
  \Eff_{\mathrm{equity}}
  \Rightarrow
  \sum_t I^+_{\partial_L}(t)>\sum_t I^-_{\partial_L}(t).
\]
这说明在有限证据域内，偏微/微分下降、\TPHLHTS{} 命中传导和整体工程兑现共同成立。

由定理~\ref{thm:tex58-expression-substitution}、定理~\ref{thm:tex58-solve-substitution}
和引理~\ref{lem:tex58-interpretation-equivalence}，在分离性、投影相容性与解释函子交换条件下，
\[
  \ExprSubst_{\PDE}^{\Lmeas}
  \wedge
  \SolveSubst_{\PDE}^{\Lmeas}
  \wedge
  \InterpEquiv_{\PDE}^{\Lmeas}
\]
成立。
由定理~\ref{thm:tex58-total-l-measure-substitution}，
\[
  \Subst_{\PDE}^{\Lmeas}=\top.
\]
再由命题~\ref{prop:tex58-tunnel-l-measure-operator}，\(\Tunnel_G\) 正是 \(L\)-测度路径空间上的低作用量命中算子，
故
\[
  \mathcal D_{\mathrm{replayid}}
  \models
  \Subst_{\PDE}^{\Lmeas}(\Tunnel_G)=\top.
\]
证毕。
\end{proof}

\subsection{备注：L-等价类边界与最终口径}
\label{subsec:tex58-l-measure-remarks}

\begin{remark}[L-等价类边界]
表达替代与求解替代成立在 \(L\)-等价类意义下。可定义
\[
  u\sim_L v
  \quad\Longleftrightarrow\quad
  u=v\quad \mu_L\text{-a.e.}
\]
或更一般地
\[
  u\sim_L v
  \quad\Longleftrightarrow\quad
  \mathcal C_L(u)=\mathcal C_L(v).
\]
因此，解释严格等价不是符号外形相同，而是残差、边界、约束、工程兑现、风险审计与 JacGap 证书这些解释谓词完全一致。
\end{remark}

\begin{remark}[不是裸差分替代 PDE]
本章结论不把 \(L_0\) 裸差分提升为第一性对象。相反，\(\partial_Lu=dD_Lu/d\mu_L\) 表明：
\[
  \boxed{
  \text{偏微分是 }L\text{-测度变化的局部密度。}
  }
\]
所以 \(L\)-积分/测度替代不是降级为粗糙差分，而是从局部密度表达提升为全局测度表达。
可用差分仍需经由 \(\sigma_{\mathfrak L}\) 提升到 \(L_1\) 语义层。
\end{remark}

\begin{remark}[最终严格结论]
最终结论为
\[
  \boxed{
  \begin{gathered}
  \text{PDE 的局部微分表达可以被 }L\text{-积分/测度表达替代；}\\
  \text{PDE 的函数求解可以被 }L\text{-路径测度/低作用量命中替代；}\\
  \text{二者在 Law-Space 解释函子下严格等价。}
  \end{gathered}
  }
\]
在有限证据域 \(\mathcal D_{\mathrm{replayid}}\) 中，该结论登记为
\[
  \boxed{
  \mathcal D_{\mathrm{replayid}}
  \models
  \left(
  \Eff_{\partial\JacGap}
  \wedge
  \Hit_{\TPH}^{\partial_L\to\Phi}
  \wedge
  \Eff_{\mathrm{equity}}
  \right)
  \Rightarrow
  \Subst_{\PDE}^{\Lmeas}(\Tunnel_G)=\top.
  }
\]
该式证明的是有限证据域内的 \(L\)-积分/测度化工程替代证书，不取消 PDE 的存在性、
唯一性、正则性与一般弱解理论，也不声称所有未来运行或所有无限候选集全局最优自动成立。
\end{remark}

%%%%%%%%%%%%%%%%%%%%%%%%%%%%%%%%%%%%%%%%%%%%%%%%%%%%%%%%%%%%
\clearpage

\section{形式逻辑：显式 PDE 边际控制到 L-测度隧穿范式的迁移}
\label{sec:tex58-explicit-to-lmeasure-tunnel-paradigm}
%%%%%%%%%%%%%%%%%%%%%%%%%%%%%%%%%%%%%%%%%%%%%%%%%%%%%%%%%%%%

\begin{remark}[章节目标与引用口径]
本节把旧范式
\[
  \mathsf{A}
  =
  \bigl(
  \PDE_{\mathrm{exp}},
  \Marg_{\PDE},
  \Ctrl_{\partial}^{\mathrm{exp}},
  \Phi
  \bigr)
\]
形式化迁移为新范式
\[
  \mathsf{B}
  =
  \bigl(
  \PDE_{\mathrm{imp}}^L,
  \Hom_{\mathrm{twist}},
  \Tunnel_G^{\TPHLHTS},
  \Desc_{\partial}^{L},
  \Phi
  \bigr).
\]
核心结论是
\[
  \boxed{
  \mathsf{A}
  \equiv_{\mathrm{solve+interp}}
  \mathsf{B},
  \qquad
  \mathsf{B}
  \succcurlyeq_{\mathrm{expr+method}}
  \mathsf{A}.
  }
\]
若存在非显式同伦路径产生更低作用量，则第二个关系升级为严格优势
\[
  \mathsf{B}
  \succ_{\mathrm{expr+method}}
  \mathsf{A}.
\]

本节沿用前文关于 PDE 低残差集合、\(L\)-测度表达、路径作用量、\TPHLHTS{} 命中证书与解释函子交换性的设定
\cite{RepoTex55PDEParadigm,RepoTex53StrategyCluster,RepoTex56LHTS,RepoTex57TPHLHTS}，
并继承 \(docs\) 发布稿中关于 JacGap、结构障碍与工程兑现目标的背景
\cite{RepoAppMathVol1,RepoEconomicsVol}。
外部数学基础仍采用 PDE 弱形式、泛函分析、测度论、梯度流、数值优化、路径积分与随机扰动理论的标准参考
\cite{Evans2010PDE,Brezis2011FunctionalAnalysis,Bogachev2007MeasureTheory,
AmbrosioGigliSavare2008GradientFlows,NocedalWright2006NumericalOptimization,
FeynmanHibbs1965PathIntegrals,FreidlinWentzell2012RandomPerturbations}。
\end{remark}

\subsection{对象域：旧范式、新范式与迁移谓词}
\label{subsec:tex58-paradigm-objects}

\begin{definition}[显式 PDE 范式]
\label{def:tex58-explicit-pde-paradigm}
令
\[
  \PDE_{\mathrm{exp}}
  :
  \Pcal[u]=0,
  \qquad
  \Bcal[u]=0
\]
为显式偏微系统。定义显式残差能量
\[
  E_{\mathrm{exp}}(u)
  \DefEq
  \frac12\norm{\Pcal[u]}_{\Hcal_P}^{2}
  +
  \frac{\lambda}{2}\norm{\Bcal[u]}_{\Hcal_B}^{2},
  \qquad
  \lambda>0.
\]
定义基于显式 PDE 的边际函数
\[
  \Marg_{\PDE}(u;\delta u)
  \DefEq
  D\Phi(u)[\delta u]
  -
  \rho\,D\Risk(u)[\delta u]
  -
  \eta\,DE_{\mathrm{exp}}(u)[\delta u],
  \qquad
  \rho,\eta>0.
\]
显式偏微控制算子记为
\[
  \Ctrl_{\partial}^{\mathrm{exp}}(u)
  \in
  \argmax_{\delta u\in\Ucal}
  \Marg_{\PDE}(u;\delta u).
\]
旧范式定义为
\[
  \mathsf{A}
  \DefEq
  \bigl(
  \PDE_{\mathrm{exp}},
  \Marg_{\PDE},
  \Ctrl_{\partial}^{\mathrm{exp}},
  \Phi
  \bigr).
\]
\end{definition}

\begin{definition}[L-测度隐式 PDE 范式]
\label{def:tex58-lmeasure-implicit-paradigm}
令 Law-Space 结构为
\[
  \mathfrak L
  =
  (X,\Sigma_L,\mu_L,\mathcal T_L,\mathcal C_L),
\]
其中 \(\mu_L\) 是 \(L\)-测度，\(\mathcal T_L\) 是测试函数族，\(\mathcal C_L\) 是解释函子。
对测试对 \((\varphi,\psi)\in\mathcal T_L\times\mathcal T_{\partial L}\)，定义 \(L\)-积分弱表达
\[
  \mathcal W_L(u;\varphi,\psi)
  \DefEq
  \int_X
  \ip{\Pcal[u]}{\varphi}\,d\mu_L
  +
  \int_{\partial X}
  \ip{\Bcal[u]}{\psi}\,d\sigma_L.
\]
定义 \(L\)-隐式残差
\[
  E_L(u)
  \DefEq
  \sup_{\substack{
    \norm{\varphi}_{\mathcal T_L}\le1\\
    \norm{\psi}_{\mathcal T_{\partial L}}\le1
  }}
  \left|
  \mathcal W_L(u;\varphi,\psi)
  \right|^2.
\]
于是
\[
  \PDE_{\mathrm{imp}}^L
  :
  E_L(u)=0.
\]
\end{definition}

\begin{definition}[同伦扭曲、路径作用量与 TPH-LHTS 微观隧穿命中]
\label{def:tex58-twist-action-hit}
定义 \(L\)-路径空间
\[
  \Gamma_L
  \DefEq
  \{\gamma:[0,1]\to\Ucal_L\}.
\]
同伦扭曲族定义为
\[
  \Hom_{\theta}:\Gamma_L\to\Gamma_L,
  \qquad
  \gamma_{\theta}=\Hom_{\theta}(\gamma_0).
\]
路径作用量定义为
\[
  \mathcal A_L(\gamma)
  \DefEq
  \int_0^1
  \left(
  aE_L(\gamma(s))
  +
  b\JacGap(\gamma(s))
  +
  c\Risk(\gamma(s))
  -
  d\Phi(\gamma(s))
  \right)\,ds,
  \qquad
  a,b,c,d>0.
\]
\TPHLHTS{} 微观隧穿命中谓词定义为
\[
\begin{aligned}
  \Hit_{\TPHLHTS}(\gamma)=\top
  \Longleftrightarrow
  \exists s^\ast\in[0,1]\ \text{使得}\quad
  &E_L(\gamma(s^\ast))\le\varepsilon_E,\\
  &\JacGap(\gamma(s^\ast))\le\varepsilon_J,\\
  &\Risk(\gamma(s^\ast))\le\varepsilon_R,\\
  &\Phi(\gamma(s^\ast))-\Phi(\gamma(0))>0.
\end{aligned}
\]
微观隧穿算子定义为
\[
  \Tunnel_G^{\TPHLHTS}
  :
  \Gamma_L
  \to
  \Argmin_{\gamma\in\Gamma_L}\mathcal A_L(\gamma).
\]
新范式定义为
\[
  \mathsf{B}
  \DefEq
  \bigl(
  \PDE_{\mathrm{imp}}^L,
  \Hom_{\mathrm{twist}},
  \Tunnel_G^{\TPHLHTS},
  \Desc_{\partial}^{L},
  \Phi
  \bigr).
\]
\end{definition}

\begin{definition}[范式迁移谓词]
\label{def:tex58-paradigm-migration-predicates}
定义表达等价谓词
\[
  \Eq_{\mathrm{expr}}(\mathsf A,\mathsf B)
  \DefEq
  \forall u\in\Ucal_L,\quad
  \bigl[
  E_{\mathrm{exp}}(u)=0
  \Longleftrightarrow
  E_L(u)=0
  \bigr].
\]
定义求解等价谓词
\[
  \Eq_{\mathrm{solve}}(\mathsf A,\mathsf B)
  \DefEq
  \Solve(\PDE_{\mathrm{exp}})
  \equiv
  \Solve(\Tunnel_G^{\TPHLHTS}).
\]
定义解释等价谓词
\[
  \Eq_{\mathrm{interp}}(\mathsf A,\mathsf B)
  \DefEq
  \forall \gamma\in\Gamma_L,\quad
  \mathcal C_L^{\mathrm{exp}}(\Pi_L\gamma)
  =
  \mathcal C_L^{\mathrm{imp}}(\gamma).
\]
定义表达--方法优势谓词
\[
  \Opt_{\mathrm{expr+method}}(\mathsf B,\mathsf A)
  \DefEq
  \Dom(\PDE_{\mathrm{exp}})
  \subseteq
  \Dom(\PDE_{\mathrm{imp}}^L)
  \ \wedge\
  \mathcal A_L(\gamma^\ast)\le
  \mathcal A_L(\gamma_{\mathrm{exp}}),
\]
其中
\[
  \gamma^\ast\in\Argmin_{\gamma\in\Gamma_L}\mathcal A_L(\gamma),
  \qquad
  \gamma_{\mathrm{exp}}\in\Gamma_L
\]
是显式边际控制路径在 \(L\)-路径空间中的嵌入。
\end{definition}

\subsection{公理降级为定理：四分量范式迁移的归纳闭合}
\label{subsec:tex58-paradigm-induction}

\begin{theorem}[公理降级：四分量范式迁移归纳定理]
\label{thm:tex58-paradigm-component-induction}
令四个迁移分量为
\[
\begin{aligned}
  q_0 &: \PDE_{\mathrm{exp}}\mapsto\PDE_{\mathrm{imp}}^L,\\
  q_1 &: \Marg_{\PDE}\mapsto\mathcal A_L,\\
  q_2 &: \Ctrl_{\partial}^{\mathrm{exp}}\mapsto\Tunnel_G^{\TPHLHTS},\\
  q_3 &: \Phi\mapsto\Phi.
\end{aligned}
\]
若每个 \(q_i\) 均保持其对应的有效性谓词
\[
  Q_i=\top,
  \qquad
  i=0,1,2,3,
\]
其中
\[
\begin{aligned}
  Q_0&\DefEq \Eq_{\mathrm{expr}}(\mathsf A,\mathsf B),\\
  Q_1&\DefEq
  \left[
  \Marg_{\PDE}\text{ 的局部方向选择可嵌入 }\mathcal A_L\text{ 的路径代价}
  \right],\\
  Q_2&\DefEq
  \left[
  \Ctrl_{\partial}^{\mathrm{exp}}\text{ 的候选路径属于 }\Gamma_L
  \text{ 且被 }\Tunnel_G^{\TPHLHTS}\text{ 比较}
  \right],\\
  Q_3&\DefEq
  \left[
  \Phi\text{ 在两种范式中由同一解释函子审计}
  \right],
\end{aligned}
\]
则累计迁移谓词
\[
  M_3
  \DefEq
  \bigwedge_{i=0}^{3}Q_i
\]
成立，并推出
\[
  \mathsf A
  \preccurlyeq_{\mathrm{embed}}
  \mathsf B.
\]
\end{theorem}

\begin{Certificate}[数学归纳法证书：四分量迁移闭合]
\label{cert:tex58-paradigm-component-induction}
定义
\[
  M_n
  \DefEq
  \bigwedge_{i=0}^{n}Q_i,
  \qquad
  n=0,1,2,3.
\]
归纳证书登记为
\[
\begin{aligned}
  I_0 &: M_0=Q_0,\\
  I_1 &: \forall n\in\{0,1,2\},\quad M_n\wedge Q_{n+1}\Rightarrow M_{n+1},\\
  I_2 &: M_3=Q_0\wedge Q_1\wedge Q_2\wedge Q_3,\\
  I_3 &: M_3\Rightarrow
  \mathsf A\preccurlyeq_{\mathrm{embed}}\mathsf B.
\end{aligned}
\]
\end{Certificate}

\begin{proof}[证明（数学归纳法证书；基于数学符号推演逻辑的谓词因果链）]
基步：
\[
  M_0=\bigwedge_{i=0}^{0}Q_i=Q_0.
\]
这表示显式 PDE 表达已经迁移为 \(L\)-测度隐式表达，且解集意义下的表达等价被保留。

归纳步：设 \(M_n=\bigwedge_{i=0}^{n}Q_i\) 成立，并给定 \(Q_{n+1}\)。
则
\[
  M_n\wedge Q_{n+1}
  =
  \left(\bigwedge_{i=0}^{n}Q_i\right)\wedge Q_{n+1}
  =
  \bigwedge_{i=0}^{n+1}Q_i
  =
  M_{n+1}.
\]
对 \(n=0,1,2\) 依次应用归纳步，得到
\[
  M_3=Q_0\wedge Q_1\wedge Q_2\wedge Q_3.
\]
由 \(Q_0\) 得表达等价；由 \(Q_1\) 得边际函数方向选择可写入路径作用量；
由 \(Q_2\) 得显式偏微控制路径被纳入 \(\Gamma_L\) 并受 \(\Tunnel_G^{\TPHLHTS}\) 比较；
由 \(Q_3\) 得工程兑现目标 \(\Phi\) 在两种范式中保持同一审计含义。
故旧范式的每个有效分量均可在新范式中找到保真嵌入，
\[
  \mathsf A\preccurlyeq_{\mathrm{embed}}\mathsf B.
\]
证毕。
\end{proof}

\subsection{定理：显式 PDE 与 L-测度隐式 PDE 的表达等价}
\label{subsec:tex58-expression-equivalence}

\begin{theorem}[表达等价定理]
\label{thm:tex58-explicit-implicit-expression-equivalence}
若测试函数族满足残差分离性
\[
  \Sep_L^{\PDE}=\top,
\]
即
\[
  \forall u\in\Ucal_L,\quad
  \left[
  \forall(\varphi,\psi)\in\mathcal T_L\times\mathcal T_{\partial L},\
  \mathcal W_L(u;\varphi,\psi)=0
  \right]
  \Rightarrow
  \left[
  \Pcal[u]=0\wedge\Bcal[u]=0
  \right],
\]
则
\[
  \PDE_{\mathrm{exp}}
  \equiv
  \PDE_{\mathrm{imp}}^L,
  \qquad
  \text{即}
  \qquad
  \Pcal[u]=0\wedge\Bcal[u]=0
  \Longleftrightarrow
  E_L(u)=0.
\]
\end{theorem}

\begin{Certificate}[证书：零残差到零弱表达，分离性反推强残差]
\label{cert:tex58-explicit-implicit-expression-equivalence}
证书链为
\[
\begin{aligned}
  E_1 &: \Pcal[u]=0\wedge\Bcal[u]=0,\\
  E_2 &: E_1\Rightarrow
  \forall(\varphi,\psi),\ \mathcal W_L(u;\varphi,\psi)=0,\\
  E_3 &: E_2\Rightarrow E_L(u)=0,\\
  E_4 &: E_L(u)=0\Rightarrow
  \forall(\varphi,\psi),\ \mathcal W_L(u;\varphi,\psi)=0,\\
  E_5 &: E_4\wedge\Sep_L^{\PDE}
  \Rightarrow \Pcal[u]=0\wedge\Bcal[u]=0.
\end{aligned}
\]
\end{Certificate}

\begin{proof}[证明（基于数学符号推演逻辑的谓词因果链）]
若
\[
  \Pcal[u]=0,
  \qquad
  \Bcal[u]=0,
\]
则对任意测试对 \((\varphi,\psi)\) 有
\[
  \mathcal W_L(u;\varphi,\psi)
  =
  \int_X\ip{0}{\varphi}\,d\mu_L
  +
  \int_{\partial X}\ip{0}{\psi}\,d\sigma_L
  =
  0.
\]
因此
\[
  E_L(u)
  =
  \sup_{\norm{\varphi}\le1,\norm{\psi}\le1}
  |\mathcal W_L(u;\varphi,\psi)|^2
  =
  0.
\]

反向设 \(E_L(u)=0\)。由 \(E_L\) 的非负性和上确界定义可得
\[
  \forall(\varphi,\psi)\in\mathcal T_L\times\mathcal T_{\partial L},
  \qquad
  \mathcal W_L(u;\varphi,\psi)=0.
\]
再由残差分离性 \(\Sep_L^{\PDE}\)，推出
\[
  \Pcal[u]=0,
  \qquad
  \Bcal[u]=0.
\]
故显式 PDE 与 \(L\)-测度隐式 PDE 在解集合意义下表达等价。证毕。
\end{proof}

\subsection{引理：投影忠实性与解释函子交换}
\label{subsec:tex58-projection-commutation-lemma}

\begin{lemma}[投影--解释交换引理]
\label{lem:tex58-projection-interpretation-commutation}
设投影
\[
  \Pi_L:\Gamma_L\to\Ucal_L
\]
满足投影忠实性
\[
  \Faith(\Pi_L)=\top
  \DefEq
  \forall\gamma\in\Gamma_L,\quad
  E_L(\Pi_L\gamma)=0
  \Longleftrightarrow
  \Pi_L\gamma\in\Sol_0.
\]
又设解释函子满足交换性
\[
  \Comm(\mathcal C_L,\Pi_L)=\top
  \DefEq
  \forall\gamma\in\Gamma_L,\quad
  \mathcal C_L^{\mathrm{exp}}(\Pi_L\gamma)
  =
  \mathcal C_L^{\mathrm{imp}}(\gamma).
\]
则任意被 \(L\)-路径命中的零残差投影解，与其显式 PDE 解在工程解释上严格等价。
\end{lemma}

\begin{Certificate}[证书：零残差投影 + 函子交换]
\label{cert:tex58-projection-interpretation-commutation}
证书链为
\[
\begin{aligned}
  C_1 &: \gamma^\ast\in\Gamma_L,\\
  C_2 &: E_L(\Pi_L\gamma^\ast)=0,\\
  C_3 &: C_2\wedge\Faith(\Pi_L)
  \Rightarrow \Pi_L\gamma^\ast\in\Sol_0,\\
  C_4 &: u^\ast=\Pi_L\gamma^\ast,\\
  C_5 &: C_4\wedge\Comm(\mathcal C_L,\Pi_L)
  \Rightarrow
  \mathcal C_L^{\mathrm{exp}}(u^\ast)
  =
  \mathcal C_L^{\mathrm{imp}}(\gamma^\ast).
\end{aligned}
\]
\end{Certificate}

\begin{proof}[证明（基于数学符号推演逻辑的谓词因果链）]
取 \(\gamma^\ast\in\Gamma_L\)，并令
\[
  u^\ast=\Pi_L\gamma^\ast.
\]
若
\[
  E_L(u^\ast)=E_L(\Pi_L\gamma^\ast)=0,
\]
则由投影忠实性得到
\[
  u^\ast\in\Sol_0.
\]
由解释函子交换性，
\[
  \mathcal C_L^{\mathrm{exp}}(\Pi_L\gamma^\ast)
  =
  \mathcal C_L^{\mathrm{imp}}(\gamma^\ast).
\]
代入 \(u^\ast=\Pi_L\gamma^\ast\)，得到
\[
  \mathcal C_L^{\mathrm{exp}}(u^\ast)
  =
  \mathcal C_L^{\mathrm{imp}}(\gamma^\ast).
\]
因此显式 PDE 解与隐式路径测度解在工程解释上严格等价。证毕。
\end{proof}

\subsection{定理：求解与解释等价}
\label{subsec:tex58-solve-interpretation-equivalence}

\begin{theorem}[求解--解释等价定理]
\label{thm:tex58-solve-interpretation-equivalence}
若以下条件成立：
\[
  \Sep_L^{\PDE}=\top,
  \qquad
  \Faith(\Pi_L)=\top,
  \qquad
  \Cover(\Gamma_L,\Sol_0)=\top,
  \qquad
  \Comm(\mathcal C_L,\Pi_L)=\top,
\]
其中路径族覆盖性为
\[
  \Cover(\Gamma_L,\Sol_0)=\top
  \DefEq
  \forall u^\ast\in\Sol_0,\ \exists\gamma^\ast\in\Gamma_L,\
  \Pi_L\gamma^\ast=u^\ast,
\]
则
\[
  \mathsf A
  \equiv_{\mathrm{solve+interp}}
  \mathsf B.
\]
\end{theorem}

\begin{Certificate}[证书：表达等价 + 覆盖 + 忠实投影 + 解释交换]
\label{cert:tex58-solve-interpretation-equivalence}
证书链为
\[
\begin{aligned}
  S_1 &: \Sep_L^{\PDE}
  \Rightarrow
  \PDE_{\mathrm{exp}}\equiv\PDE_{\mathrm{imp}}^L,\\
  S_2 &: u^\ast\in\Sol_0
  \wedge
  \Cover(\Gamma_L,\Sol_0)
  \Rightarrow
  \exists\gamma^\ast\in\Gamma_L,\ \Pi_L\gamma^\ast=u^\ast,\\
  S_3 &: E_L(\Pi_L\gamma^\ast)=0
  \wedge
  \Faith(\Pi_L)
  \Rightarrow
  \Pi_L\gamma^\ast\in\Sol_0,\\
  S_4 &: \Comm(\mathcal C_L,\Pi_L)
  \Rightarrow
  \mathcal C_L^{\mathrm{exp}}(\Pi_L\gamma^\ast)
  =
  \mathcal C_L^{\mathrm{imp}}(\gamma^\ast),\\
  S_5 &: S_1\wedge S_2\wedge S_3\wedge S_4
  \Rightarrow
  \mathsf A\equiv_{\mathrm{solve+interp}}\mathsf B.
\end{aligned}
\]
\end{Certificate}

\begin{proof}[证明（基于数学符号推演逻辑的谓词因果链）]
由定理~\ref{thm:tex58-explicit-implicit-expression-equivalence}，
\[
  \Sep_L^{\PDE}=\top
  \Rightarrow
  \left[
  \PDE_{\mathrm{exp}}
  \equiv
  \PDE_{\mathrm{imp}}^L
  \right].
\]
因此显式解集与 \(L\)-隐式零残差解集一致。

先从旧范式到新范式。若
\[
  u^\ast\in\Sol_0,
\]
由覆盖性存在 \(\gamma^\ast\in\Gamma_L\) 使
\[
  \Pi_L\gamma^\ast=u^\ast.
\]
由表达等价，
\[
  E_L(u^\ast)=0,
\]
故
\[
  E_L(\Pi_L\gamma^\ast)=0.
\]
于是旧范式解可由新范式路径投影表示。

再从新范式到旧范式。若 \(\gamma^\ast\in\Gamma_L\) 命中零残差投影
\[
  E_L(\Pi_L\gamma^\ast)=0,
\]
由投影忠实性得到
\[
  \Pi_L\gamma^\ast\in\Sol_0.
\]
令 \(u^\ast=\Pi_L\gamma^\ast\)，则 \(u^\ast\) 是显式 PDE 解。

最后由引理~\ref{lem:tex58-projection-interpretation-commutation}，
\[
  \mathcal C_L^{\mathrm{exp}}(u^\ast)
  =
  \mathcal C_L^{\mathrm{imp}}(\gamma^\ast).
\]
因此二者不仅求解对象等价，而且工程解释等价，得到
\[
  \mathsf A
  \equiv_{\mathrm{solve+interp}}
  \mathsf B.
\]
证毕。
\end{proof}

\subsection{命题：表达域扩张与路径方法优化}
\label{subsec:tex58-expression-method-optimization}

\begin{proposition}[表达--方法优化命题]
\label{prop:tex58-expression-method-optimization}
若显式 PDE 范式要求经典偏导数以逐点意义存在，而 \(L\)-隐式表达只要求
\(\mathcal W_L(u;\varphi,\psi)\) 在测试函数与 \(L\)-测度下有意义，则
\[
  \Dom(\PDE_{\mathrm{exp}})
  \subseteq
  \Dom(\PDE_{\mathrm{imp}}^L).
\]
若存在
\[
  u\in
  \Dom(\PDE_{\mathrm{imp}}^L)
  \setminus
  \Dom(\PDE_{\mathrm{exp}}),
\]
则该包含为严格包含。

若显式边际控制路径可嵌入 \(L\)-路径空间，
\[
  \gamma_{\mathrm{exp}}\in\Gamma_L,
\]
且
\[
  \gamma^\ast\in\Argmin_{\gamma\in\Gamma_L}\mathcal A_L(\gamma),
\]
则
\[
  \mathcal A_L(\gamma^\ast)
  \le
  \mathcal A_L(\gamma_{\mathrm{exp}}).
\]
若存在非显式同伦路径 \(\widetilde\gamma\in\Gamma_L\) 使得
\[
  \mathcal A_L(\widetilde\gamma)
  <
  \mathcal A_L(\gamma_{\mathrm{exp}}),
\]
则
\[
  \mathcal A_L(\gamma^\ast)
  <
  \mathcal A_L(\gamma_{\mathrm{exp}}).
\]
\end{proposition}

\begin{Certificate}[证书：弱表达域包含 + 作用量最小性]
\label{cert:tex58-expression-method-optimization}
证书链为
\[
\begin{aligned}
  O_1 &: u\in\Dom(\PDE_{\mathrm{exp}})
  \Rightarrow
  \Pcal[u],\Bcal[u]\text{ 逐点可解释},\\
  O_2 &: O_1\Rightarrow
  \mathcal W_L(u;\varphi,\psi)\text{ 对所有测试对有意义},\\
  O_3 &: O_2\Rightarrow u\in\Dom(\PDE_{\mathrm{imp}}^L),\\
  O_4 &: \gamma^\ast\in\Argmin_{\gamma\in\Gamma_L}\mathcal A_L(\gamma)
  \Rightarrow
  \forall\gamma\in\Gamma_L,\ \mathcal A_L(\gamma^\ast)\le\mathcal A_L(\gamma),\\
  O_5 &: \gamma_{\mathrm{exp}}\in\Gamma_L
  \Rightarrow
  \mathcal A_L(\gamma^\ast)\le\mathcal A_L(\gamma_{\mathrm{exp}}),\\
  O_6 &: \exists\widetilde\gamma,\
  \mathcal A_L(\widetilde\gamma)<\mathcal A_L(\gamma_{\mathrm{exp}})
  \Rightarrow
  \mathcal A_L(\gamma^\ast)<\mathcal A_L(\gamma_{\mathrm{exp}}).
\end{aligned}
\]
\end{Certificate}

\begin{proof}[证明（基于数学符号推演逻辑的谓词因果链）]
若 \(u\in\Dom(\PDE_{\mathrm{exp}})\)，则 \(\Pcal[u]\) 与 \(\Bcal[u]\) 在显式 PDE 所需的经典意义下可定义。
因此对任意测试对 \((\varphi,\psi)\)，积分表达
\[
  \mathcal W_L(u;\varphi,\psi)
\]
也可在 \(L\)-测度下解释，得到
\[
  u\in\Dom(\PDE_{\mathrm{imp}}^L).
\]
故
\[
  \Dom(\PDE_{\mathrm{exp}})
  \subseteq
  \Dom(\PDE_{\mathrm{imp}}^L).
\]
若存在属于右侧但不属于左侧的弱正则、测度型或路径型对象，则该包含严格。

另一方面，由
\[
  \gamma^\ast\in\Argmin_{\gamma\in\Gamma_L}\mathcal A_L(\gamma)
\]
可得
\[
  \forall\gamma\in\Gamma_L,\qquad
  \mathcal A_L(\gamma^\ast)\le\mathcal A_L(\gamma).
\]
取 \(\gamma=\gamma_{\mathrm{exp}}\)，得到
\[
  \mathcal A_L(\gamma^\ast)
  \le
  \mathcal A_L(\gamma_{\mathrm{exp}}).
\]
若存在 \(\widetilde\gamma\in\Gamma_L\) 满足
\[
  \mathcal A_L(\widetilde\gamma)
  <
  \mathcal A_L(\gamma_{\mathrm{exp}}),
\]
则由最小性
\[
  \mathcal A_L(\gamma^\ast)
  \le
  \mathcal A_L(\widetilde\gamma)
  <
  \mathcal A_L(\gamma_{\mathrm{exp}}).
\]
这给出严格方法优化。证毕。
\end{proof}

\subsection{定理：TPH-LHTS 命中下的偏微下降工程兑现}
\label{subsec:tex58-tph-lhts-descent-realization}

\begin{theorem}[\TPHLHTS{} 命中兑现定理]
\label{thm:tex58-tph-lhts-descent-realization}
若存在路径 \(\gamma^\ast\in\Gamma_L\) 使
\[
  \Hit_{\TPHLHTS}(\gamma^\ast)=\top,
\]
且偏微 JacGap 有效性与权益兑现有效性成立：
\[
  \Eff_{\partial\JacGap}=\top,
  \qquad
  \Eff_{\mathrm{equity}}=\top,
\]
则新范式给出偏微下降工程兑现：
\[
  \Desc_{\partial}^{L}(\gamma^\ast)
  \Rightarrow
  \Phi\uparrow
  \wedge
  E_L\downarrow
  \wedge
  \JacGap\downarrow
  \wedge
  \Risk\downarrow.
\]
\end{theorem}

\begin{Certificate}[证书：命中阈值到工程兑现]
\label{cert:tex58-tph-lhts-descent-realization}
证书链为
\[
\begin{aligned}
  R_1 &: \Hit_{\TPHLHTS}(\gamma^\ast)=\top,\\
  R_2 &: R_1\Rightarrow
  \exists s^\ast,\
  E_L(\gamma^\ast(s^\ast))\le\varepsilon_E
  \wedge
  \JacGap(\gamma^\ast(s^\ast))\le\varepsilon_J\\
  &\hspace{42mm}\wedge\
  \Risk(\gamma^\ast(s^\ast))\le\varepsilon_R
  \wedge
  \Phi(\gamma^\ast(s^\ast))>\Phi(\gamma^\ast(0)),\\
  R_3 &: R_2\wedge\Eff_{\partial\JacGap}
  \Rightarrow
  E_L\downarrow\wedge\JacGap\downarrow,\\
  R_4 &: R_2\wedge\Eff_{\mathrm{equity}}
  \Rightarrow
  \Phi\uparrow\wedge\Risk\downarrow,\\
  R_5 &: R_3\wedge R_4
  \Rightarrow
  \Desc_{\partial}^{L}(\gamma^\ast)\Rightarrow\Phi\text{ 工程兑现}.
\end{aligned}
\]
\end{Certificate}

\begin{proof}[证明（基于数学符号推演逻辑的谓词因果链）]
由命中谓词定义，存在 \(s^\ast\in[0,1]\) 使
\[
  E_L(\gamma^\ast(s^\ast))\le\varepsilon_E,
  \qquad
  \JacGap(\gamma^\ast(s^\ast))\le\varepsilon_J,
\]
\[
  \Risk(\gamma^\ast(s^\ast))\le\varepsilon_R,
  \qquad
  \Phi(\gamma^\ast(s^\ast))-\Phi(\gamma^\ast(0))>0.
\]
由 \(\Eff_{\partial\JacGap}\)，低 \(E_L\) 与低 \(\JacGap\) 被登记为偏微下降和结构障碍下降；
由 \(\Eff_{\mathrm{equity}}\)，\(\Phi\) 的正增量与风险阈值共同登记为工程兑现有效。
四个谓词合取后得到
\[
  \Phi\uparrow,
  \qquad
  E_L\downarrow,
  \qquad
  \JacGap\downarrow,
  \qquad
  \Risk\downarrow.
\]
这正是新范式中偏微下降工程兑现的形式含义。证毕。
\end{proof}

\subsection{总定理：求解解释等价、表达方法优化}
\label{subsec:tex58-paradigm-main-theorem}

\begin{theorem}[范式迁移总定理]
\label{thm:tex58-paradigm-migration-main}
若
\[
  \Sep_L^{\PDE}
  \wedge
  \Faith(\Pi_L)
  \wedge
  \Cover(\Gamma_L,\Sol_0)
  \wedge
  \Comm(\mathcal C_L,\Pi_L)
  \wedge
  \Hit_{\TPHLHTS}
  \wedge
  \Eff_{\partial\JacGap}
  \wedge
  \Eff_{\mathrm{equity}}
\]
成立，并且显式边际控制路径可嵌入 \(L\)-路径空间，则
\[
  \boxed{
  \mathsf A
  \equiv_{\mathrm{solve+interp}}
  \mathsf B
  }
\]
且
\[
  \boxed{
  \mathsf B
  \succcurlyeq_{\mathrm{expr+method}}
  \mathsf A.
  }
\]
若进一步存在非显式同伦路径 \(\widetilde\gamma\) 使
\[
  \mathcal A_L(\widetilde\gamma)
  <
  \mathcal A_L(\gamma_{\mathrm{exp}}),
\]
则
\[
  \boxed{
  \mathsf B
  \succ_{\mathrm{expr+method}}
  \mathsf A.
  }
\]
\end{theorem}

\begin{Certificate}[证书：总谓词合成]
\label{cert:tex58-paradigm-migration-main}
证书链为
\[
\begin{aligned}
  T_1 &: \Sep_L^{\PDE}\Rightarrow
  \Eq_{\mathrm{expr}}(\mathsf A,\mathsf B),\\
  T_2 &: \Faith(\Pi_L)\wedge\Cover(\Gamma_L,\Sol_0)
  \Rightarrow
  \Eq_{\mathrm{solve}}(\mathsf A,\mathsf B),\\
  T_3 &: \Comm(\mathcal C_L,\Pi_L)
  \Rightarrow
  \Eq_{\mathrm{interp}}(\mathsf A,\mathsf B),\\
  T_4 &: \Hit_{\TPHLHTS}\wedge\Eff_{\partial\JacGap}\wedge\Eff_{\mathrm{equity}}
  \Rightarrow
  \Desc_{\partial}^{L}\Rightarrow\Phi\uparrow,\\
  T_5 &: \Dom(\PDE_{\mathrm{exp}})
  \subseteq
  \Dom(\PDE_{\mathrm{imp}}^L),\\
  T_6 &: \gamma^\ast\in\Argmin_{\gamma\in\Gamma_L}\mathcal A_L
  \wedge
  \gamma_{\mathrm{exp}}\in\Gamma_L
  \Rightarrow
  \mathcal A_L(\gamma^\ast)\le\mathcal A_L(\gamma_{\mathrm{exp}}),\\
  T_7 &: T_1\wedge T_2\wedge T_3
  \Rightarrow
  \mathsf A\equiv_{\mathrm{solve+interp}}\mathsf B,\\
  T_8 &: T_4\wedge T_5\wedge T_6
  \Rightarrow
  \mathsf B\succcurlyeq_{\mathrm{expr+method}}\mathsf A.
\end{aligned}
\]
\end{Certificate}

\begin{proof}[证明（基于数学符号推演逻辑的谓词因果链）]
由定理~\ref{thm:tex58-explicit-implicit-expression-equivalence}，
\[
  \Sep_L^{\PDE}
  \Rightarrow
  \PDE_{\mathrm{exp}}\equiv\PDE_{\mathrm{imp}}^L.
\]
由定理~\ref{thm:tex58-solve-interpretation-equivalence}，
\[
  \Faith(\Pi_L)
  \wedge
  \Cover(\Gamma_L,\Sol_0)
  \wedge
  \Comm(\mathcal C_L,\Pi_L)
  \Rightarrow
  \mathsf A\equiv_{\mathrm{solve+interp}}\mathsf B.
\]
由定理~\ref{thm:tex58-tph-lhts-descent-realization}，
\[
  \Hit_{\TPHLHTS}
  \wedge
  \Eff_{\partial\JacGap}
  \wedge
  \Eff_{\mathrm{equity}}
  \Rightarrow
  \Desc_{\partial}^{L}\Rightarrow\Phi\uparrow.
\]
由命题~\ref{prop:tex58-expression-method-optimization}，
\[
  \Dom(\PDE_{\mathrm{exp}})
  \subseteq
  \Dom(\PDE_{\mathrm{imp}}^L)
\]
且
\[
  \mathcal A_L(\gamma^\ast)
  \le
  \mathcal A_L(\gamma_{\mathrm{exp}}).
\]
因此新范式在表达域上不窄于旧范式，在路径方法上不劣于旧范式，得到
\[
  \mathsf B
  \succcurlyeq_{\mathrm{expr+method}}
  \mathsf A.
\]
若存在 \(\widetilde\gamma\) 满足
\[
  \mathcal A_L(\widetilde\gamma)
  <
  \mathcal A_L(\gamma_{\mathrm{exp}}),
\]
则由作用量最小性
\[
  \mathcal A_L(\gamma^\ast)
  \le
  \mathcal A_L(\widetilde\gamma)
  <
  \mathcal A_L(\gamma_{\mathrm{exp}}),
\]
故新范式在方法上严格优于旧范式，并与表达域扩张合成为
\[
  \mathsf B
  \succ_{\mathrm{expr+method}}
  \mathsf A.
\]
证毕。
\end{proof}

\subsection{推论：旧范式作为新范式的显式局部特例}
\label{subsec:tex58-old-paradigm-special-case}

\begin{corollary}[显式局部特例推论]
\label{cor:tex58-old-paradigm-special-case}
在约束
\[
  \Hom_{\mathrm{twist}}=\mathrm{id},
  \qquad
  \nu=\delta_{\gamma_{\mathrm{exp}}},
  \qquad
  \Gamma_L=\{\gamma_{\mathrm{exp}}\}
\]
下，新范式退化为旧范式的显式、光滑、局部、单路径版本：
\[
  \mathsf A
  =
  \mathsf B
  \bigm|_{
    \Hom=\mathrm{id},\
    \nu=\delta_{\gamma_{\mathrm{exp}}},\
    \Gamma_L=\{\gamma_{\mathrm{exp}}\}
  }.
\]
因此旧范式可嵌入新范式。
\end{corollary}

\begin{Certificate}[证书：同伦恒等化 + 路径测度退化 + 单路径搜索]
\label{cert:tex58-old-paradigm-special-case}
证书链为
\[
\begin{aligned}
  P_1 &: \Hom_{\mathrm{twist}}=\mathrm{id}
  \Rightarrow
  \text{无同伦扭曲},\\
  P_2 &: \nu=\delta_{\gamma_{\mathrm{exp}}}
  \Rightarrow
  \text{路径测度集中于显式控制路径},\\
  P_3 &: \Gamma_L=\{\gamma_{\mathrm{exp}}\}
  \Rightarrow
  \Argmin_{\gamma\in\Gamma_L}\mathcal A_L(\gamma)=\{\gamma_{\mathrm{exp}}\},\\
  P_4 &: P_1\wedge P_2\wedge P_3
  \Rightarrow
  \Tunnel_G^{\TPHLHTS}
  \text{ 退化为显式偏微控制路径选择},\\
  P_5 &: P_4\Rightarrow
  \mathsf A=
  \mathsf B\bigm|_{\Hom=\mathrm{id},\nu=\delta_{\gamma_{\mathrm{exp}}},\Gamma_L=\{\gamma_{\mathrm{exp}}\}}.
\end{aligned}
\]
\end{Certificate}

\begin{proof}[证明（基于数学符号推演逻辑的谓词因果链）]
若
\[
  \Hom_{\mathrm{twist}}=\mathrm{id},
\]
则路径不发生同伦扭曲。若
\[
  \nu=\delta_{\gamma_{\mathrm{exp}}},
\]
则路径测度全部集中在显式控制路径。若
\[
  \Gamma_L=\{\gamma_{\mathrm{exp}}\},
\]
则路径作用量最小化不再比较额外路径，而是只返回
\[
  \gamma_{\mathrm{exp}}.
\]
于是
\[
  \Tunnel_G^{\TPHLHTS}
  :
  \Gamma_L\to\Argmin_{\gamma\in\Gamma_L}\mathcal A_L(\gamma)
\]
退化为显式局部控制路径选择；\(L\)-测度隐式表达在光滑且分离条件成立时与显式 PDE 表达等价。
因此旧范式是新范式在恒等同伦、Dirac 路径测度与单路径候选族下的特例。证毕。
\end{proof}

\begin{corollary}[范式迁移压缩式]
\label{cor:tex58-paradigm-migration-compressed}
该迁移可压缩登记为
\[
  \boxed{
  \begin{aligned}
  &\left[
  \PDE_{\mathrm{exp}}
  +
  \Marg_{\PDE}
  +
  \Ctrl_{\partial}^{\mathrm{exp}}
  \Rightarrow
  \Phi
  \right]\\
  &\qquad\leadsto
  \left[
  \PDE_{\mathrm{imp}}^L
  +
  \Hom_{\mathrm{twist}}
  +
  \Tunnel_G^{\TPHLHTS}
  +
  \Desc_{\partial}^{L}
  \Rightarrow
  \Phi
  \right].
  \end{aligned}
  }
\]
并满足
\[
  \boxed{
  \begin{aligned}
  \PDE_{\mathrm{exp}}
  &\equiv_{\mathrm{solve+interp}}
  \PDE_{\mathrm{imp}}^L
  +
  \Tunnel_G^{\TPHLHTS}
  \end{aligned}
  }
\]
以及
\[
  \boxed{
  \begin{aligned}
  \PDE_{\mathrm{imp}}^L
  +
  \Tunnel_G^{\TPHLHTS}
  &\succcurlyeq_{\mathrm{expr+method}}
  \PDE_{\mathrm{exp}}\\
  &\quad+
  \Marg_{\PDE}
  +
  \Ctrl_{\partial}^{\mathrm{exp}}.
  \end{aligned}
  }
\]
\end{corollary}

\begin{Certificate}[证书：压缩式对应关系]
\label{cert:tex58-paradigm-migration-compressed}
压缩式证书为
\[
\begin{aligned}
  K_1 &: \PDE_{\mathrm{exp}}\leftrightarrow\PDE_{\mathrm{imp}}^L,\\
  K_2 &: \Marg_{\PDE}\leftrightarrow\mathcal A_L,\\
  K_3 &: \Ctrl_{\partial}^{\mathrm{exp}}\leftrightarrow\Tunnel_G^{\TPHLHTS},\\
  K_4 &: \Phi\leftrightarrow\Phi,\\
  K_5 &: K_1\wedge K_2\wedge K_3\wedge K_4
  \Rightarrow
  \text{求解解释等价与表达方法优化同时登记}.
\end{aligned}
\]
\end{Certificate}

\begin{proof}[证明（基于数学符号推演逻辑的谓词因果链）]
\(K_1\) 由表达等价定理给出；\(K_2\) 由边际方向代价嵌入路径作用量给出；
\(K_3\) 由显式控制路径嵌入 \(\Gamma_L\) 并由 \(\Tunnel_G^{\TPHLHTS}\) 统一比较给出；
\(K_4\) 表示工程兑现目标在迁移前后不改变。
四者合取后，旧范式的局部控制链转化为新范式的全局变分链。
再由范式迁移总定理，求解解释等价与表达方法优化同时成立。证毕。
\end{proof}

\subsection{备注：等价、优化与边界口径}
\label{subsec:tex58-paradigm-migration-remarks}

\begin{remark}[等价的精确定义]
本节中的“等价”指
\[
  \boxed{
  \text{PDE 解集合等价、Law-Space 工程解释等价、兑现谓词等价。}
  }
\]
它不要求两种范式具有相同的符号外观，也不要求采用相同的数值求解路径。
\end{remark}

\begin{remark}[优化的精确定义]
本节中的“优化”指
\[
  \boxed{
  \begin{gathered}
  \text{表达域不缩小，路径搜索空间不缩小，显式局部路径可嵌入，}\\
  \text{非显式同伦路径可参与比较。}
  \end{gathered}
  }
\]
因此新范式保留旧范式的可解释解集合，同时把求解路线从显式边际函数控制扩展为
\[
  L\text{-测度隐式残差}
  +
  \text{同伦扭曲}
  +
  \text{路径作用量}
  +
  \TPHLHTS{}\text{ 微观隧穿命中}
  +
  L\text{-偏微下降工程兑现}.
\]
\end{remark}

\begin{remark}[最终登记口径]
最终结论可写为
\[
  \boxed{
  \begin{gathered}
  \text{该转化把“PDE 显式表达 + 显式边际函数偏微控制”的局部控制范式，}\\
  \text{提升为“PDE \(L\)-测度隐式表达 + 同伦扭曲 + \TPHLHTS{} 微观隧穿命中”的全局变分范式；}\\
  \text{二者在 PDE 解集合与 Law-Space 工程解释上严格等价，}\\
  \text{后者在表达域、路径搜索能力、离散障碍修复能力与工程兑现证书化方面更强。}
  \end{gathered}
  }
\]
这就是 PDE 求解和解释等价的问题转化，以及 PDE 表达方式与求解方法的优化。
\end{remark}

%%%%%%%%%%%%%%%%%%%%%%%%%%%%%%%%%%%%%%%%%%%%%%%%%%%%%%%%%%%%
\clearpage

\section{形式逻辑：L-测度替代框架的结构增益与成熟度证书}
\label{sec:tex58-l-measure-gain-maturity}
%%%%%%%%%%%%%%%%%%%%%%%%%%%%%%%%%%%%%%%%%%%%%%%%%%%%%%%%%%%%

前四章已经给出从偏微残差下降、有限回放正影响证书、微观隧穿下降路径替代，到
\(L\)-积分/测度表达替代的形式系统链条。为了避免把该链条仅写成叙述性总结，本章把“增益”
本身形式化为可登记的元证明对象：它不是只证明某个算法能辅助 PDE 求解，而是把 PDE 局部微分范式
重写为
\[
  \boxed{
  \begin{gathered}
  \text{局部微分表达}
  \Rightarrow
  L\text{-积分/测度表达替代}\\
  \Rightarrow
  L\text{-路径测度/微观隧穿求解替代}
  \Rightarrow
  \text{Law-Space 解释严格等价}.
  \end{gathered}
  }
\]
这一元证明依赖前三类基础：PDE 弱形式与泛函分析基础
\cite{Evans2010PDE,BrennerScott2008FEM,Brezis2011FunctionalAnalysis}；
测度密度、几乎处处等价和路径测度基础
\cite{Bogachev2007MeasureTheory,AmbrosioGigliSavare2008GradientFlows}；
仓内 \(\GFramework{}\) 的语义规范、PDE 命中、\LHTS{} 与 \TPHLHTS{} 证书链
\cite{RepoTex37SemanticGauge,RepoTex55PDEParadigm,RepoTex56LHTS,RepoTex57TPHLHTS}。
本章继续只引用 \(NoteBook\) 下 \(\_pub\) 稿件及 \(docs\) 发布稿。

本章的目标不是引入新的 PDE 解概念，而是给前述结果登记成熟度：
\[
  \boxed{
  \Maturity_{\mathrm{tex58}}
  =
  \text{有限证据域强支持命题}
  +
  \text{一般替代定理框架}
  +
  \text{四类待补强全称条件}.
  }
\]

\subsection{对象域、增益谓词与成熟度条件}
\label{subsec:tex58-gain-maturity-definitions}

\begin{definition}[三层结构增益谓词]
\label{def:tex58-three-layer-gain}
定义三层结构增益谓词为
\[
\begin{aligned}
  \Gain_{\mathrm{expr}}
  &\DefEq
  \ExprSubst_{\PDE}^{\Lmeas},\\
  \Gain_{\mathrm{solve}}
  &\DefEq
  \SolveSubst_{\PDE}^{\Lmeas},\\
  \Gain_{\mathrm{interp}}
  &\DefEq
  \InterpEquiv_{\PDE}^{\Lmeas}.
\end{aligned}
\]
定义总结构增益为
\[
  \Gain_{\mathrm{struct}}
  \DefEq
  \Gain_{\mathrm{expr}}
  \wedge
  \Gain_{\mathrm{solve}}
  \wedge
  \Gain_{\mathrm{interp}}.
\]
由定义~\ref{def:tex58-three-l-measure-substitutions}，
\[
  \Gain_{\mathrm{struct}}
  \Longleftrightarrow
  \Subst_{\PDE}^{\Lmeas}(\Tunnel_G)=\top.
\]
\end{definition}

\begin{definition}[四类全称化补强条件]
\label{def:tex58-four-strengthening-conditions}
定义四类补强条件：
\[
\begin{aligned}
  \Sep_L &:=
  \Sep(\mathcal T_L,\mathcal T_{\partial L}),\\
  \Faith_L &:=
  \Faith(\Pi_L),\\
  \Comm_L &:=
  \Comm(\mathcal C_L,\Pi_L),\\
  \Cover_L &:=
  \Cover(\Gamma_L,\mathcal U_L).
\end{aligned}
\]
其中 \(\Sep_L\) 表示测试函数族对内部残差与边界残差具有分离性；
\(\Faith_L\) 表示投影忠实性，即 \(E_L(\nu)=0\Rightarrow E_{\PDE}(\Pi_L\nu)=0\)；
\(\Comm_L\) 表示解释函子交换性，即
\[
  \mathcal C_L^{\Lmeas}(\nu)
  =
  \mathcal C_L^{\PDE}(\Pi_L\nu);
\]
\(\Cover_L\) 表示路径族覆盖性，即 \(\Gamma_L\) 对候选解域 \(\mathcal U_L\) 具有足够覆盖，
使低作用量命中不是仅在过窄候选族内成立。
\end{definition}

\subsection{公理降级为定理：三层替代增益的归纳闭合}
\label{subsec:tex58-three-layer-gain-induction}

\begin{theorem}[公理降级：三层替代增益归纳定理]
\label{thm:tex58-three-layer-gain-induction}
设
\[
  g_0=\Gain_{\mathrm{expr}},
  \qquad
  g_1=\Gain_{\mathrm{solve}},
  \qquad
  g_2=\Gain_{\mathrm{interp}}.
\]
若 \(g_0,g_1,g_2\) 均成立，则
\[
  \bigwedge_{i=0}^{2}g_i
  =
  \Gain_{\mathrm{struct}}
  =
  \Subst_{\PDE}^{\Lmeas}(\Tunnel_G)=\top.
\]
\end{theorem}

\begin{Certificate}[数学归纳法证书：三层替代谓词逐层闭合]
\label{cert:tex58-three-layer-gain-induction}
定义累计替代谓词
\[
  G_n
  \DefEq
  \bigwedge_{i=0}^{n}g_i,
  \qquad
  n=0,1,2.
\]
证书链为
\[
\begin{aligned}
  A_1 &: G_0=g_0=\Gain_{\mathrm{expr}},\\
  A_2 &: G_n\wedge g_{n+1}\Rightarrow G_{n+1},\qquad n=0,1,\\
  A_3 &: G_2=g_0\wedge g_1\wedge g_2,\\
  A_4 &: G_2
  \Longleftrightarrow
  \ExprSubst_{\PDE}^{\Lmeas}
  \wedge
  \SolveSubst_{\PDE}^{\Lmeas}
  \wedge
  \InterpEquiv_{\PDE}^{\Lmeas},\\
  A_5 &: A_4
  \Longleftrightarrow
  \Subst_{\PDE}^{\Lmeas}(\Tunnel_G)=\top.
\end{aligned}
\]
\end{Certificate}

\begin{proof}[证明（基于数学符号推演逻辑的谓词因果链）]
归纳基 \(n=0\)：
\[
  G_0=g_0=\Gain_{\mathrm{expr}}.
\]
归纳假设：设 \(G_n=\bigwedge_{i=0}^{n}g_i\) 成立。
归纳推进：
\[
  G_{n+1}
  =
  \left(\bigwedge_{i=0}^{n}g_i\right)\wedge g_{n+1}.
\]
故 \(G_n\wedge g_{n+1}\Rightarrow G_{n+1}\)。对 \(n=0,1\) 连续应用该推进，得到
\[
  G_2
  =
  g_0\wedge g_1\wedge g_2.
\]
代回 \(g_i\) 的定义，
\[
  G_2
  =
  \ExprSubst_{\PDE}^{\Lmeas}
  \wedge
  \SolveSubst_{\PDE}^{\Lmeas}
  \wedge
  \InterpEquiv_{\PDE}^{\Lmeas}.
\]
由定义~\ref{def:tex58-three-l-measure-substitutions}，
\[
  G_2
  \Longleftrightarrow
  \Subst_{\PDE}^{\Lmeas}(\Tunnel_G)=\top.
\]
证毕。
\end{proof}

\subsection{定理：命题强度升级为表达求解解释三层替代}
\label{subsec:tex58-strength-upgrade-theorem}

\begin{theorem}[命题强度升级定理]
\label{thm:tex58-strength-upgrade}
若
\[
  \Subst_{\PDE}^{\Lmeas}(\Tunnel_G)=\top,
\]
则该结论严格强于“\(\Tunnel_G\) 只替代 PDE 下降路径搜索任务”的单层工程替代口径。形式化地，
\[
  \Subst_{\PDE}^{\Lmeas}(\Tunnel_G)=\top
  \Rightarrow
  \left[
  \Tunnel_G
  \succcurlyeq_{\mathrm{eng}}
  \Solve_{\PDE}^{\mathrm{descent}}
  \right],
\]
但反向蕴含不自动成立。
\end{theorem}

\begin{Certificate}[证书：三层替代蕴含单层下降路径替代]
\label{cert:tex58-strength-upgrade}
证书链为
\[
\begin{aligned}
  B_1 &: \Subst_{\PDE}^{\Lmeas}(\Tunnel_G)=\top,\\
  B_2 &: B_1\Rightarrow \SolveSubst_{\PDE}^{\Lmeas}=\top,\\
  B_3 &: \SolveSubst_{\PDE}^{\Lmeas}=\top
  \Rightarrow
  E_L\downarrow\ \text{或}\ E_L=0\ \text{可由低作用量命中给出},\\
  B_4 &: B_3\Rightarrow
  \Tunnel_G\succcurlyeq_{\mathrm{eng}}\Solve_{\PDE}^{\mathrm{descent}},\\
  B_5 &: \left[
  \Tunnel_G\succcurlyeq_{\mathrm{eng}}\Solve_{\PDE}^{\mathrm{descent}}
  \right]
  \not\Rightarrow
  \ExprSubst_{\PDE}^{\Lmeas}\wedge\InterpEquiv_{\PDE}^{\Lmeas}.
\end{aligned}
\]
\end{Certificate}

\begin{proof}[证明（基于数学符号推演逻辑的谓词因果链）]
由定理~\ref{thm:tex58-three-layer-gain-induction}，
\[
  \Subst_{\PDE}^{\Lmeas}(\Tunnel_G)=\top
\]
包含
\[
  \SolveSubst_{\PDE}^{\Lmeas}=\top.
\]
由定理~\ref{thm:tex58-solve-substitution}，求解替代意味着 PDE 求解可转化为 \(E_L=0\)
或 \(L\)-路径测度低作用量命中。因此低作用量命中包含 PDE 下降路径搜索任务，得到
\[
  \Tunnel_G\succcurlyeq_{\mathrm{eng}}\Solve_{\PDE}^{\mathrm{descent}}.
\]
然而，仅有下降路径搜索替代只能说明某类路径使残差下降；它不能单独推出 PDE 局部微分表达已被
\(L\)-积分/测度表达替代，也不能单独推出 Law-Space 解释函子交换。因此反向蕴含不自动成立。
证毕。
\end{proof}

\subsection{引理：替代基础不是 L0 裸差分而是 L-测度提升}
\label{subsec:tex58-not-l0-difference-lemma}

\begin{lemma}[非 \(L_0\) 裸差分替代引理]
\label{lem:tex58-not-l0-difference}
若
\[
  \partial_Lu=\frac{dD_Lu}{d\mu_L},
\]
且可用差分满足
\[
  \Delta_{\mathrm{usable}}=\sigma_{\mathfrak L}(\Delta_0)\in L_1,
\]
则 PDE 的 \(L\)-测度替代不是 \(L_0\) 裸差分替代，而是
\[
  \text{局部密度表达}
  \leadsto
  \text{全局 }L\text{-测度表达}
\]
的表达层级提升。
\end{lemma}

\begin{Certificate}[证书：裸差分障碍到测度提升]
\label{cert:tex58-not-l0-difference}
证书链为
\[
\begin{aligned}
  C_1 &: \Delta_0\in L_0,\\
  C_2 &: L_0\not\models
  \Directional\wedge\Integrable\wedge\Auditable,\\
  C_3 &: \Delta_{\mathrm{usable}}=\sigma_{\mathfrak L}(\Delta_0)\in L_1,\\
  C_4 &: \partial_Lu=dD_Lu/d\mu_L,\\
  C_5 &: C_3\wedge C_4
  \Rightarrow
  \partial_Lu\text{ 是 }L\text{-测度变化的局部密度},\\
  C_6 &: C_1\wedge C_2\wedge C_5
  \Rightarrow
  \text{替代基础不是 }L_0\text{ 裸差分}.
\end{aligned}
\]
\end{Certificate}

\begin{proof}[证明（基于数学符号推演逻辑的谓词因果链）]
由定理~\ref{thm:tex58-l1-minimal-meaningful}，\(L_0\) 裸差分不自动携带方向性、
可积性和审计性；由推论~\ref{subsec:tex58-usable-difference-corollary}，可用于求解的差分必须经
\(\sigma_{\mathfrak L}\) 提升为 \(L_1\) 语义对象。
又由定义~\ref{def:tex58-partial-as-density}，
\[
  \partial_Lu=\frac{dD_Lu}{d\mu_L}.
\]
因此偏微分是 \(L\)-变化测度相对于 \(L\)-测度的局部密度。
把 PDE 写成 \(L\)-积分/测度形式，是从局部密度表达提升到全局测度表达，而不是用裸差分替代 PDE。
证毕。
\end{proof}

\subsection[命题：微观隧穿统一接口算子]{命题：微观隧穿是 PDE、测度、JacGap 与工程兑现的统一接口算子}
\label{subsec:tex58-unified-interface-proposition}

\begin{proposition}[统一接口算子命题]
\label{prop:tex58-unified-interface-operator}
在 \(L\)-路径空间 \(\Gamma_L\) 上，
\[
  \Tunnel_G:\Gamma_L\to\Argmin_{\gamma\in\Gamma_L}\mathcal A_L(\gamma)
\]
同时连接四类目标：
\[
  E_L\downarrow,
  \qquad
  \JacGap\downarrow,
  \qquad
  \Risk\downarrow,
  \qquad
  \Phi\uparrow.
\]
因此 \(\Tunnel_G\) 是 PDE、测度、变分、JacGap、\TPHLHTS{} 和工程兑现之间的统一接口算子。
\end{proposition}

\begin{Certificate}[证书：作用量四分量到统一接口]
\label{cert:tex58-unified-interface-operator}
证书链为
\[
\begin{aligned}
  D_1 &: \mathcal A_L(\gamma)
  =
  \int_0^1(aE_L+b\JacGap+c\Risk-d\Phi)\,ds,\\
  D_2 &: a,b,c,d>0,\\
  D_3 &: \gamma^\ast\in\Argmin\mathcal A_L,\\
  D_4 &: D_1\wedge D_2\wedge D_3
  \Rightarrow
  E_L\downarrow\wedge\JacGap\downarrow\wedge\Risk\downarrow\wedge\Phi\uparrow,\\
  D_5 &: D_4\Rightarrow
  \Tunnel_G\text{ 统一 PDE 残差、结构修复、风险约束与工程兑现}.
\end{aligned}
\]
\end{Certificate}

\begin{proof}[证明（基于数学符号推演逻辑的谓词因果链）]
由定义~\ref{def:tex58-l-residual-energy-action}，\(\mathcal A_L\) 将 \(E_L\)、\(\JacGap\)、
\(\Risk\) 与 \(\Phi\) 放入同一作用量。由于 \(a,b,c,d>0\)，低作用量路径倾向于降低
\(E_L,\JacGap,\Risk\)，并提升 \(\Phi\)。因此
\[
  \gamma^\ast\in\Argmin\mathcal A_L
\]
给出的是低残差、低结构障碍、低风险、高兑现的共同命中，而不是单一数值步骤。
所以 \(\Tunnel_G\) 在形式上承担 PDE/测度/变分/JacGap/\TPHLHTS{}/工程兑现的接口角色。证毕。
\end{proof}

\subsection{定理：成熟度边界与四重补强条件}
\label{subsec:tex58-maturity-boundary-theorem}

\begin{theorem}[成熟度边界定理]
\label{thm:tex58-maturity-boundary}
在有限证据域 \(\mathcal D_{\mathrm{replayid}}\) 内，若
\[
  \Eff_{\partial\JacGap}
  \wedge
  \Hit_{\TPH}^{\partial_L\to\Phi}
  \wedge
  \Eff_{\mathrm{equity}}
\]
成立，则得到有限证据域强支持命题
\[
  \mathcal D_{\mathrm{replayid}}
  \models
  \Subst_{\PDE}^{\Lmeas}(\Tunnel_G)=\top.
\]
若进一步满足
\[
  \Sep_L\wedge\Faith_L\wedge\Comm_L\wedge\Cover_L,
\]
则该强支持命题可提升为 Law-Space 内的一般替代定理框架。
\end{theorem}

\begin{Certificate}[证书：有限证据到一般框架的四重门槛]
\label{cert:tex58-maturity-boundary}
证书链为
\[
\begin{aligned}
  F_1 &: \Eff_{\partial\JacGap}\wedge
  \Hit_{\TPH}^{\partial_L\to\Phi}\wedge
  \Eff_{\mathrm{equity}},\\
  F_2 &: F_1\Rightarrow
  \mathcal D_{\mathrm{replayid}}
  \models
  \sum_tI^+_{\partial_L}(t)>\sum_tI^-_{\partial_L}(t),\\
  F_3 &: \Sep_L
  \Rightarrow
  \ExprSubst_{\PDE}^{\Lmeas},\\
  F_4 &: \Faith_L
  \Rightarrow
  \SolveSubst_{\PDE}^{\Lmeas}\text{ 与 }E_{\PDE}\text{ 相容},\\
  F_5 &: \Comm_L
  \Rightarrow
  \InterpEquiv_{\PDE}^{\Lmeas},\\
  F_6 &: \Cover_L
  \Rightarrow
  \Argmin_{\Gamma_L}\mathcal A_L\text{ 可代表候选解域内的低作用量搜索},\\
  F_7 &: F_2\wedge F_3\wedge F_4\wedge F_5\wedge F_6
  \Rightarrow
  \Maturity_{\mathrm{tex58}}=\text{一般替代定理框架}.
\end{aligned}
\]
\end{Certificate}

\begin{proof}[证明（基于数学符号推演逻辑的谓词因果链）]
由推论~\ref{cor:tex58-finite-domain-l-measure-substitution}，有限证据域条件给出
\[
  \mathcal D_{\mathrm{replayid}}
  \models
  \Subst_{\PDE}^{\Lmeas}(\Tunnel_G)=\top.
\]
这是有限证据域强支持命题。

若 \(\Sep_L\) 成立，则弱表达能够反推 PDE 残差为零，表达替代具备全称反推条件。
若 \(\Faith_L\) 成立，则 \(L\)-测度解经 \(\Pi_L\) 投影后保持 PDE 残差忠实性。
若 \(\Comm_L\) 成立，则两种表达在 \(\mathcal C_L\) 下解释一致。
若 \(\Cover_L\) 成立，则路径空间命中不是仅对狭窄候选集有效，而是在指定候选解域中具有覆盖意义。
四者合取后，有限证据域强支持命题具备推广为 Law-Space 内一般替代定理框架所需的形式门槛。证毕。
\end{proof}

\subsection{推论：结构性高增益与成熟度登记}
\label{subsec:tex58-structural-high-gain-corollary}

\begin{corollary}[结构性高增益推论]
\label{cor:tex58-structural-high-gain}
在第~\ref{sec:tex58-l-measure-pde-substitution} 节的三层替代已经成立、且本章四重补强条件被显式登记的条件下，
本文形成的结论不是单点算法收益，而是结构性高增益：
\[
  \boxed{
  \PDE
  \leadsto
  \Lmeas
  \leadsto
  \Gamma_L
  \xrightarrow{\Tunnel_G}
  \Argmin\mathcal A_L
  \leadsto
  \Phi\uparrow,\ \Risk\downarrow.
  }
\]
\end{corollary}

\begin{Certificate}[证书：结构性高增益登记]
\label{cert:tex58-structural-high-gain}
证书链为
\[
\begin{aligned}
  H_1 &: \ExprSubst_{\PDE}^{\Lmeas}
  \Rightarrow
  \PDE\text{ 局部表达被 }L\text{-积分/测度表达替代},\\
  H_2 &: \SolveSubst_{\PDE}^{\Lmeas}
  \Rightarrow
  \PDE\text{ 求解被 }L\text{-路径测度低作用量命中替代},\\
  H_3 &: \InterpEquiv_{\PDE}^{\Lmeas}
  \Rightarrow
  \text{二者解释严格等价},\\
  H_4 &: H_1\wedge H_2\wedge H_3
  \Rightarrow
  \Gain_{\mathrm{struct}},\\
  H_5 &: \Sep_L\wedge\Faith_L\wedge\Comm_L\wedge\Cover_L
  \Rightarrow
  \text{全称化补强条件已登记},\\
  H_6 &: H_4\wedge H_5
  \Rightarrow
  \Maturity_{\mathrm{tex58}}
  =
  \text{结构性高增益}.
\end{aligned}
\]
\end{Certificate}

\begin{proof}[证明（基于数学符号推演逻辑的谓词因果链）]
\(\ExprSubst\) 给出表达层增益，\(\SolveSubst\) 给出求解层增益，\(\InterpEquiv\)
给出解释层增益。三者合取由定理~\ref{thm:tex58-three-layer-gain-induction} 闭合为
\[
  \Gain_{\mathrm{struct}}
  =
  \Subst_{\PDE}^{\Lmeas}(\Tunnel_G)=\top.
\]
四重补强条件则把有限证据域强支持命题与一般 Law-Space 替代定理框架之间的缺口显式列出。
因此成熟度登记不是“已无边界的全称定理”，而是“有限证据域内强支持命题已经成型，
一般理论层面的替代定理框架已经成型，全称版本的分离性、忠实性、交换性和覆盖性条件已明确”。证毕。
\end{proof}

\subsection{备注：成熟度评分、边界与最终口径}
\label{subsec:tex58-gain-maturity-remarks}

\begin{remark}[成熟度评分的工程口径]
若把本章作为研究推进状态登记，可给出工程成熟度注记
\[
  \mathrm{Score}_{\mathrm{maturity}}
  =
  9.2/10.
\]
该分值不是数学定理，而是对当前证明骨架的工程评价：理论结构已从“下降路径搜索替代”
提升到“表达替代、求解替代、解释严格等价”；有限证据域证书已经闭合；一般全称版本仍依赖
\(\Sep_L,\Faith_L,\Comm_L,\Cover_L\) 四类条件的进一步实例化。
\end{remark}

\begin{remark}[不可越界外推]
本章不证明所有 PDE 的解析解均可由 \(\Tunnel_G\) 直接给出；不证明所有未来运行均成立；
不证明无限候选集全局最优已经由有限证据域实现；不取消 PDE 的存在性、唯一性、正则性、弱解理论与数值收敛理论。
本章证明的是：
\[
  \boxed{
  \begin{gathered}
  \text{有限证据域内，}L\text{-积分/测度替代与微观隧穿命中形成强支持证书；}\\
  \text{一般 Law-Space 内，表达替代、求解替代、解释等价的证明框架已经成型。}
  \end{gathered}
  }
\]
\end{remark}

\begin{remark}[最终严格口径]
本文的结构性结论可登记为
\[
  \boxed{
  \begin{gathered}
  \text{PDE 局部微分范式可被提升为 }L\text{-积分/测度范式；}\\
  \text{PDE 求解可被提升为 }L\text{-路径测度低作用量命中；}\\
  \text{微观隧穿 } \Tunnel_G \text{ 是连接 PDE、测度、变分、JacGap、TPH 与工程兑现的统一算子；}\\
  \text{二者在 Law-Space 解释函子下保持严格解释等价。}
  \end{gathered}
  }
\]
这就是本文所需的“结构性高增益”结论。
\end{remark}

\clearpage

%%%%%%%%%%%%%%%%%%%%%%%%%%%%%%%%%%%%%%%%%%%%%%%%%%%%%%%%%%%%
\section[形式逻辑：PDE 到 L-测度隧穿的战略降维]{形式逻辑：传统 PDE 求解到 L-测度隧穿求解的战略工程降维优势}
\label{sec:tex58-strategic-engineering-dimred}
%%%%%%%%%%%%%%%%%%%%%%%%%%%%%%%%%%%%%%%%%%%%%%%%%%%%%%%%%%%%

本章把传统 PDE 求解任务
\[
  \Pcal[u]=0,
  \qquad
  \Bcal[u]=0
\]
与 \(L\)-测度同伦隧穿求解任务放在同一谓词系统中比较。传统求解的核心是寻找满足方程、
边界条件和函数空间约束的函数 \(u\)，即压低或归零残差
\[
  E_{\PDE}(u)\to0.
\]
\(L\)-测度同伦隧穿范式则采用链条
\[
  \PDE
  \rightarrow
  \Lmeas
  \rightarrow
  \Gamma_L
  \rightarrow
  \Argmin_{\gamma\in\Gamma_L}\mathcal A_L(\gamma)
  \rightarrow
  \Phi\uparrow,\ \Risk\downarrow,\ \Audit=\top.
\]
该链条继承前文关于 PDE 残差下降、\(L\)-测度弱表达、路径作用量、
\TPHLHTS{} 命中与结构成熟度证书的结论
\cite{RepoTex55PDEParadigm,RepoTex56LHTS,RepoTex57TPHLHTS}，
并使用经典 PDE 弱形式、测度论、梯度流、数值优化与路径积分背景
\cite{Evans2010PDE,BrennerScott2008FEM,Bogachev2007MeasureTheory,
AmbrosioGigliSavare2008GradientFlows,NocedalWright2006NumericalOptimization,
FeynmanHibbs1965PathIntegrals,FreidlinWentzell2012RandomPerturbations}。
本章仍只引用 \(NoteBook\) 下 \(\_pub\) 稿件和 \(docs\) 发布稿作为仓内来源。

\subsection{对象域与比较谓词}
\label{subsec:tex58-strategic-objects}

\begin{definition}[传统 PDE 求解任务]
\label{def:tex58-traditional-pde-task}
定义传统 PDE 满足谓词为
\[
  \mathsf{StdPDE}(u)
  \DefEq
  \bigl(\Pcal[u]=0\bigr)\wedge\bigl(\Bcal[u]=0\bigr).
\]
定义传统 PDE 求解任务为
\[
  \mathsf{Solve}_{\mathrm{std}}
  \DefEq
  \exists u\in\Vcal,\quad \mathsf{StdPDE}(u).
\]
在残差口径下，该任务等价于寻找
\[
  E_{\PDE}(u)
  \DefEq
  \frac12\norm{\Pcal[u]}_{\Hcal_P}^{2}
  +
  \frac{\lambda}{2}\norm{\Bcal[u]}_{\Hcal_B}^{2}
  \to0,
  \qquad \lambda>0.
\]
因此传统任务的核心判断是：给定物理约束、边界条件和函数空间后，某个候选函数是否满足
对应 PDE 约束。
\end{definition}

\begin{definition}[L-测度同伦隧穿工程求解任务]
\label{def:tex58-lmeasure-tunnel-engineering-task}
令 \(\Gamma_L\) 为 \(L\)-路径空间，令
\[
  \mathcal A_L(\gamma)
  \DefEq
  \int_0^1
  \left(
  aE_L(\gamma(s))
  +
  b\JacGap(\gamma(s))
  +
  c\Risk(\gamma(s))
  -
  d\Phi(\gamma(s))
  \right)\,ds,
  \qquad a,b,c,d>0.
\]
定义多目标命中谓词
\[
  \mathsf{MOHit}_L(\gamma^\ast)
  \DefEq
  E_L\downarrow
  \wedge
  \JacGap\downarrow
  \wedge
  \Risk\downarrow
  \wedge
  \Phi\uparrow.
\]
定义 \(L\)-测度同伦隧穿工程求解谓词
\[
  \mathsf{LTunnel}(\gamma^\ast)
  \DefEq
  \gamma^\ast\in\Argmin_{\gamma\in\Gamma_L}\mathcal A_L(\gamma)
  \wedge
  \mathsf{MOHit}_L(\gamma^\ast)
  \wedge
  \Audit\bigl(\mathsf{Cert}_L(\gamma^\ast)\bigr)=\top.
\]
该定义与第~\ref{def:tex58-l-residual-energy-action}、\ref{def:tex58-twist-action-hit} 节的
路径作用量和 \TPHLHTS{} 命中定义一致。
\end{definition}

\begin{definition}[战略工程降维谓词]
\label{def:tex58-strategic-dimred-predicate}
令分散工程任务族为
\[
  \mathsf{Split}
  \DefEq
  \mathsf{EqSolve}
  \times
  \mathsf{PathOpt}
  \times
  \mathsf{RiskCtrl}
  \times
  \mathsf{EngVerify}
  \times
  \mathsf{AuditTrace}.
\]
若存在同一低作用量路径与同一证书链
\[
  \bigl(\gamma^\ast,\mathsf{Cert}_L(\gamma^\ast)\bigr)
\]
使得
\[
  \bigl(\gamma^\ast,\mathsf{Cert}_L(\gamma^\ast)\bigr)
  \vdash
  \mathsf{Split},
\]
则定义战略工程降维谓词为
\[
  \mathsf{DimRed}_{\mathrm{eng}}^L=\top.
\]
这里的“降维”不是拓扑维数或线性代数维数命题，而是工程组织维度命题：原本分散处理的
方程求解、路径优化、风险控制、工程验证和审计归因被压缩进同一个作用量与证书链。
\end{definition}

\subsection[公理降级：多目标审计命中归纳]{公理降级为定理：多目标可审计命中的归纳闭合}
\label{subsec:tex58-multiobjective-audit-induction}

\begin{theorem}[公理降级：多目标可审计命中归纳定理]
\label{thm:tex58-multiobjective-audit-induction}
给定候选路径 \(\gamma^\ast\in\Gamma_L\)。定义五个基础命中谓词
\[
\begin{aligned}
  p_0(\gamma^\ast)&\DefEq E_L\downarrow,\\
  p_1(\gamma^\ast)&\DefEq \JacGap\downarrow,\\
  p_2(\gamma^\ast)&\DefEq \Risk\downarrow,\\
  p_3(\gamma^\ast)&\DefEq \Phi\uparrow,\\
  p_4(\gamma^\ast)&\DefEq \Audit\bigl(\mathsf{Cert}_L(\gamma^\ast)\bigr)=\top.
\end{aligned}
\]
若
\[
  \gamma^\ast\in\Argmin_{\gamma\in\Gamma_L}\mathcal A_L(\gamma)
  \qquad\text{且}\qquad
  \bigwedge_{i=0}^{4}p_i(\gamma^\ast)
\]
成立，则
\[
  \mathsf{LTunnel}(\gamma^\ast)=\top.
\]
\end{theorem}

\begin{Certificate}[数学归纳法证书：五类命中谓词逐项闭合]
\label{cert:tex58-multiobjective-audit-induction}
定义累计谓词
\[
  M_n(\gamma^\ast)
  \DefEq
  \bigwedge_{i=0}^{n}p_i(\gamma^\ast),
  \qquad n=0,1,2,3,4.
\]
归纳证书登记为
\[
\begin{aligned}
  I_1 &: M_0(\gamma^\ast)=p_0(\gamma^\ast),\\
  I_2 &: M_n(\gamma^\ast)\wedge p_{n+1}(\gamma^\ast)
         \Rightarrow M_{n+1}(\gamma^\ast),
         \qquad n=0,1,2,3,\\
  I_3 &: M_4(\gamma^\ast)
         \Longleftrightarrow
         \mathsf{MOHit}_L(\gamma^\ast)
         \wedge
         \Audit\bigl(\mathsf{Cert}_L(\gamma^\ast)\bigr)=\top,\\
  I_4 &: I_3\wedge
         \gamma^\ast\in\Argmin_{\gamma\in\Gamma_L}\mathcal A_L(\gamma)
         \Rightarrow
         \mathsf{LTunnel}(\gamma^\ast)=\top.
\end{aligned}
\]
\end{Certificate}

\begin{proof}[证明（数学归纳法证书；基于数学符号推演逻辑的谓词因果链）]
归纳基：
\[
  M_0(\gamma^\ast)
  =
  \bigwedge_{i=0}^{0}p_i(\gamma^\ast)
  =
  p_0(\gamma^\ast).
\]
故 \(I_1\) 成立。

归纳步：假设
\[
  M_n(\gamma^\ast)=\bigwedge_{i=0}^{n}p_i(\gamma^\ast)
\]
成立。若再给定 \(p_{n+1}(\gamma^\ast)\)，则
\[
  M_n(\gamma^\ast)\wedge p_{n+1}(\gamma^\ast)
  =
  \left(\bigwedge_{i=0}^{n}p_i(\gamma^\ast)\right)
  \wedge p_{n+1}(\gamma^\ast)
  =
  \bigwedge_{i=0}^{n+1}p_i(\gamma^\ast).
\]
因此
\[
  M_n(\gamma^\ast)\wedge p_{n+1}(\gamma^\ast)
  \Rightarrow
  M_{n+1}(\gamma^\ast).
\]
对 \(n=0,1,2,3\) 连续应用，得到
\[
  M_4(\gamma^\ast)
  =
  p_0\wedge p_1\wedge p_2\wedge p_3\wedge p_4.
\]
代回 \(p_i\) 的定义，
\[
  M_4(\gamma^\ast)
  \Longleftrightarrow
  E_L\downarrow
  \wedge
  \JacGap\downarrow
  \wedge
  \Risk\downarrow
  \wedge
  \Phi\uparrow
  \wedge
  \Audit\bigl(\mathsf{Cert}_L(\gamma^\ast)\bigr)=\top.
\]
由定义~\ref{def:tex58-lmeasure-tunnel-engineering-task}，
若再合取
\[
  \gamma^\ast\in\Argmin_{\gamma\in\Gamma_L}\mathcal A_L(\gamma),
\]
则得到
\[
  \mathsf{LTunnel}(\gamma^\ast)=\top.
\]
证毕。
\end{proof}

\subsection[定理：PDE 到 L-测度工程求解提升]{定理：传统 PDE 求解到 L-测度工程求解的任务提升}
\label{subsec:tex58-task-lift-theorem}

\begin{theorem}[任务提升定理]
\label{thm:tex58-task-lift}
设存在投影
\[
  \Pi_L:\Gamma_L\to\Vcal
\]
以及阈值忠实条件
\[
  \mathsf{Faith}_L^{\varepsilon}
  :
  E_L(\gamma)\le\varepsilon_E
  \Rightarrow
  E_{\PDE}(\Pi_L\gamma)\le\varepsilon_{\PDE}.
\]
若 \(\mathsf{LTunnel}(\gamma^\ast)=\top\) 且 \(E_L(\gamma^\ast)\le\varepsilon_E\)，则
\[
  \Pi_L\gamma^\ast\in\Sol_{\varepsilon_{\PDE}}
\]
并且同时成立
\[
  \JacGap\downarrow,
  \qquad
  \Risk\downarrow,
  \qquad
  \Phi\uparrow,
  \qquad
  \Audit\bigl(\mathsf{Cert}_L(\gamma^\ast)\bigr)=\top.
\]
因此 \(L\)-测度同伦隧穿求解把传统的
\[
  E_{\PDE}(u)\to0
\]
提升为
\[
  E_L\downarrow,\quad
  \JacGap\downarrow,\quad
  \Risk\downarrow,\quad
  \Phi\uparrow,\quad
  \Audit=\top.
\]
\end{theorem}

\begin{Certificate}[证书：残差满足到多目标工程满足]
\label{cert:tex58-task-lift}
证书链为
\[
\begin{aligned}
  T_1 &: \mathsf{LTunnel}(\gamma^\ast)=\top
         \Rightarrow
         \mathsf{MOHit}_L(\gamma^\ast)
         \wedge
         \Audit(\mathsf{Cert}_L(\gamma^\ast))=\top,\\
  T_2 &: \mathsf{MOHit}_L(\gamma^\ast)
         \Rightarrow
         E_L\downarrow\wedge\JacGap\downarrow\wedge\Risk\downarrow\wedge\Phi\uparrow,\\
  T_3 &: E_L(\gamma^\ast)\le\varepsilon_E
         \wedge
         \mathsf{Faith}_L^{\varepsilon}
         \Rightarrow
         E_{\PDE}(\Pi_L\gamma^\ast)\le\varepsilon_{\PDE},\\
  T_4 &: E_{\PDE}(\Pi_L\gamma^\ast)\le\varepsilon_{\PDE}
         \Rightarrow
         \Pi_L\gamma^\ast\in\Sol_{\varepsilon_{\PDE}},\\
  T_5 &: T_1\wedge T_2\wedge T_3\wedge T_4
         \Rightarrow
         \text{传统残差求解被提升为多目标工程求解}.
\end{aligned}
\]
\end{Certificate}

\begin{proof}[证明（基于数学符号推演逻辑的谓词因果链）]
由定义~\ref{def:tex58-lmeasure-tunnel-engineering-task}，
\[
  \mathsf{LTunnel}(\gamma^\ast)=\top
\]
推出
\[
  \mathsf{MOHit}_L(\gamma^\ast)
  \wedge
  \Audit(\mathsf{Cert}_L(\gamma^\ast))=\top.
\]
由 \(\mathsf{MOHit}_L\) 的定义，得到
\[
  E_L\downarrow,\quad
  \JacGap\downarrow,\quad
  \Risk\downarrow,\quad
  \Phi\uparrow.
\]
若 \(E_L(\gamma^\ast)\le\varepsilon_E\)，则由阈值忠实条件
\[
  \mathsf{Faith}_L^{\varepsilon}
\]
推出
\[
  E_{\PDE}(\Pi_L\gamma^\ast)\le\varepsilon_{\PDE}.
\]
由阈值解定义~\ref{def:tex58-solve-pde}，这等价于
\[
  \Pi_L\gamma^\ast\in\Sol_{\varepsilon_{\PDE}}.
\]
因此，传统 PDE 求解中的残差满足只是 \(L\)-测度同伦隧穿工程求解的一个投影分量；
后者还同时登记结构间隙、风险、目标兑现与审计证书。证毕。
\end{proof}

\subsection[引理：L-路径搜索空间扩展]{引理：L-路径空间严格扩展显式函数搜索空间}
\label{subsec:tex58-l-path-search-expansion}

\begin{lemma}[搜索空间严格扩展引理]
\label{lem:tex58-l-path-search-expansion}
设存在显式函数到路径空间的嵌入
\[
  \iota:\Vcal\hookrightarrow\Gamma_L,
  \qquad
  \iota(u)(s)=u.
\]
若存在同伦扭曲或弱测度路径
\[
  \gamma_{\mathrm{tw}}\in\Gamma_L
  \qquad\text{且}\qquad
  \gamma_{\mathrm{tw}}\notin\iota(\Vcal),
\]
则
\[
  \iota(\Vcal)\subsetneq\Gamma_L.
\]
因此 \(L\)-路径空间中的求解搜索严格扩展了传统显式函数空间或其离散网格中的单点搜索。
\end{lemma}

\begin{Certificate}[证书：嵌入、外点与严格包含]
\label{cert:tex58-l-path-search-expansion}
证书链为
\[
\begin{aligned}
  L_1 &: \iota:\Vcal\hookrightarrow\Gamma_L,\\
  L_2 &: L_1\Rightarrow \iota(\Vcal)\subseteq\Gamma_L,\\
  L_3 &: \exists\gamma_{\mathrm{tw}}\in\Gamma_L\setminus\iota(\Vcal),\\
  L_4 &: L_2\wedge L_3\Rightarrow \iota(\Vcal)\subsetneq\Gamma_L,\\
  L_5 &: L_4\Rightarrow
         \text{\(L\)-路径搜索严格扩展传统显式函数搜索}.
\end{aligned}
\]
\end{Certificate}

\begin{proof}[证明（基于数学符号推演逻辑的谓词因果链）]
由 \(\iota\) 是嵌入可知
\[
  \iota(\Vcal)\subseteq\Gamma_L.
\]
若存在
\[
  \gamma_{\mathrm{tw}}\in\Gamma_L
  \quad\text{且}\quad
  \gamma_{\mathrm{tw}}\notin\iota(\Vcal),
\]
则 \(\Gamma_L\) 中至少有一个元素不属于 \(\iota(\Vcal)\)。因此包含关系不是等号，而是
\[
  \iota(\Vcal)\subsetneq\Gamma_L.
\]
这说明传统显式函数空间中的候选函数可作为常值路径嵌入 \(L\)-路径空间；
但 \(L\)-路径空间还允许同伦扭曲、非显式路径、弱表达路径和低作用量命中路径。
结合定理~\ref{thm:tex58-expression-substitution} 与定理~\ref{thm:tex58-solve-substitution}，
该扩展不是裸差分扩展，而是 \(L\)-积分/测度表达与 \(L\)-路径测度求解层级上的扩展。证毕。
\end{proof}

\subsection[命题：作用量与证书链压缩工程环节]{命题：分散工程环节可压缩为同一作用量与证书链}
\label{subsec:tex58-engineering-compression-proposition}

\begin{proposition}[工程组织压缩命题]
\label{prop:tex58-engineering-compression}
若 \(\gamma^\ast\) 满足
\[
  \mathsf{LTunnel}(\gamma^\ast)=\top,
\]
且证书链 \(\mathsf{Cert}_L(\gamma^\ast)\) 同时登记
\[
  E_L,\quad
  \JacGap,\quad
  \Risk,\quad
  \Phi,\quad
  \Audit,
\]
则存在证书读取映射
\[
  \Theta(\gamma^\ast,\mathsf{Cert}_L)
  =
  \bigl(
  \mathsf{EqSolve},
  \mathsf{PathOpt},
  \mathsf{RiskCtrl},
  \mathsf{EngVerify},
  \mathsf{AuditTrace}
  \bigr),
\]
从而
\[
  \mathsf{DimRed}_{\mathrm{eng}}^L=\top.
\]
\end{proposition}

\begin{Certificate}[证书：同一低作用量路径读取五类工程结论]
\label{cert:tex58-engineering-compression}
证书链为
\[
\begin{aligned}
  C_1 &: \mathsf{LTunnel}(\gamma^\ast)=\top
         \Rightarrow
         \gamma^\ast\in\Argmin\mathcal A_L,\\
  C_2 &: \gamma^\ast\in\Argmin\mathcal A_L
         \Rightarrow
         \mathsf{PathOpt},\\
  C_3 &: E_L(\gamma^\ast)\le\varepsilon_E
         \wedge
         \mathsf{Faith}_L^{\varepsilon}
         \Rightarrow
         \mathsf{EqSolve},\\
  C_4 &: \Risk\downarrow\Rightarrow\mathsf{RiskCtrl},\\
  C_5 &: \Phi\uparrow\wedge\JacGap\downarrow
         \Rightarrow
         \mathsf{EngVerify},\\
  C_6 &: \Audit(\mathsf{Cert}_L(\gamma^\ast))=\top
         \Rightarrow
         \mathsf{AuditTrace},\\
  C_7 &: C_2\wedge C_3\wedge C_4\wedge C_5\wedge C_6
         \Rightarrow
         \bigl(\gamma^\ast,\mathsf{Cert}_L\bigr)\vdash\mathsf{Split}.
\end{aligned}
\]
\end{Certificate}

\begin{proof}[证明（基于数学符号推演逻辑的谓词因果链）]
由 \(\mathsf{LTunnel}(\gamma^\ast)=\top\)，首先得到
\[
  \gamma^\ast\in\Argmin_{\gamma\in\Gamma_L}\mathcal A_L(\gamma),
\]
所以路径优化结论 \(\mathsf{PathOpt}\) 可由同一 \(\gamma^\ast\) 读取。
若进一步有 \(E_L(\gamma^\ast)\le\varepsilon_E\) 与 \(\mathsf{Faith}_L^{\varepsilon}\)，
则由定理~\ref{thm:tex58-task-lift} 得到 PDE 阈值求解结论 \(\mathsf{EqSolve}\)。

同一命中谓词还给出
\[
  \Risk\downarrow,
  \qquad
  \Phi\uparrow,
  \qquad
  \JacGap\downarrow,
  \qquad
  \Audit(\mathsf{Cert}_L(\gamma^\ast))=\top.
\]
于是风险控制、工程验证和审计归因分别可从同一证书链读取。故
\[
  \Theta(\gamma^\ast,\mathsf{Cert}_L)
  =
  \mathsf{Split}
\]
成立，从而
\[
  \mathsf{DimRed}_{\mathrm{eng}}^L=\top.
\]
证毕。
\end{proof}

\subsection{推论：战略工程降维优势}
\label{subsec:tex58-strategic-dimred-corollary}

\begin{corollary}[战略工程降维优势推论]
\label{cor:tex58-strategic-engineering-dimred}
在定理~\ref{thm:tex58-multiobjective-audit-induction}、定理~\ref{thm:tex58-task-lift}、
引理~\ref{lem:tex58-l-path-search-expansion} 与命题~\ref{prop:tex58-engineering-compression} 的条件下，
\(L\)-测度同伦隧穿 PDE 求解相对于传统 PDE 求解形成如下严格优势：
\[
  \boxed{
  \begin{gathered}
  \text{传统：}\quad
  \exists u,\ \Pcal[u]=0,\ \Bcal[u]=0,\\
  \text{提升：}\quad
  \exists\gamma^\ast\in\Gamma_L,\quad
  \mathsf{LTunnel}(\gamma^\ast)=\top,\\
  E_L\downarrow,\quad
  \JacGap\downarrow,\quad
  \Risk\downarrow,\quad
  \Phi\uparrow,\quad
  \Audit(\mathsf{Cert}_L)=\top,\\
  \mathsf{DimRed}_{\mathrm{eng}}^L=\top.
  \end{gathered}
  }
\]
因此其革命性不在于“更快算一个 PDE”，而在于把“求函数解”升级为
“在复杂约束系统中寻找可审计、可兑现、可控风险的全局命中路径”。
\end{corollary}

\begin{Certificate}[证书：任务提升、搜索扩展、证书内置与组织压缩]
\label{cert:tex58-strategic-engineering-dimred}
证书链为
\[
\begin{aligned}
  S_1 &: \mathsf{Solve}_{\mathrm{std}}
         \Rightarrow
         E_{\PDE}(u)\to0,\\
  S_2 &: \mathsf{LTunnel}(\gamma^\ast)=\top
         \Rightarrow
         E_L\downarrow\wedge\JacGap\downarrow\wedge\Risk\downarrow\wedge\Phi\uparrow
         \wedge\Audit=\top,\\
  S_3 &: \iota(\Vcal)\subsetneq\Gamma_L
         \Rightarrow
         \text{搜索空间严格扩展},\\
  S_4 &: \bigl(\gamma^\ast,\mathsf{Cert}_L\bigr)\vdash\mathsf{Split}
         \Rightarrow
         \mathsf{DimRed}_{\mathrm{eng}}^L=\top,\\
  S_5 &: S_2\wedge S_3\wedge S_4
         \Rightarrow
         \text{可审计、可兑现、可控风险的全局命中路径}.
\end{aligned}
\]
\end{Certificate}

\begin{proof}[证明（基于数学符号推演逻辑的谓词因果链）]
传统 PDE 求解由定义~\ref{def:tex58-traditional-pde-task} 登记为
\[
  \exists u\in\Vcal,\quad E_{\PDE}(u)\to0.
\]
该任务只直接判断 PDE 残差与边界残差是否满足。

由定理~\ref{thm:tex58-multiobjective-audit-induction}，
\[
  \mathsf{LTunnel}(\gamma^\ast)=\top
\]
闭合为
\[
  \gamma^\ast\in\Argmin\mathcal A_L,
  \quad
  E_L\downarrow,
  \quad
  \JacGap\downarrow,
  \quad
  \Risk\downarrow,
  \quad
  \Phi\uparrow,
  \quad
  \Audit(\mathsf{Cert}_L)=\top.
\]
由定理~\ref{thm:tex58-task-lift}，其中 \(E_L\) 的低残差分量可投影回 PDE 阈值解；
由引理~\ref{lem:tex58-l-path-search-expansion}，\(\Gamma_L\) 严格包含传统显式函数搜索的嵌入像；
由命题~\ref{prop:tex58-engineering-compression}，方程求解、路径优化、风险控制、工程验证和审计归因
可由同一个低作用量路径和证书链读取。

于是
\[
  \mathsf{Solve}_{\mathrm{std}}
  \leadsto
  \mathsf{LTunnel}
  \leadsto
  \mathsf{DimRed}_{\mathrm{eng}}^L
\]
成立。该推演说明，新范式不是传统 PDE 求解器的单点加速版本，而是将求解对象、
搜索空间、目标函数、风险约束和审计证书同时统一到 \(L\)-测度路径作用量中的工程组织升级。
证毕。
\end{proof}

\subsection{备注：最终严格口径}
\label{subsec:tex58-strategic-dimred-remarks}

\begin{remark}[四个正确结论]
本章最终登记四个正确结论：
\[
  \boxed{
  \begin{gathered}
  \text{第一，传统 PDE 求解是 }E_{\PDE}(u)\to0\text{ 的函数解任务；}\\
  \text{第二，}L\text{-测度同伦隧穿是 }
  E_L,\JacGap,\Risk,\Phi,\Audit\text{ 的多目标工程任务；}\\
  \text{第三，若存在非显式同伦扭曲或弱测度路径，则 } \Gamma_L
  \text{ 严格扩展传统搜索空间；}\\
  \text{第四，同一作用量与证书链可压缩分散工程环节，形成战略工程降维优势。}
  \end{gathered}
  }
\]
\end{remark}

\begin{remark}[边界口径]
本章不声称任意 PDE 的解析解都由单次隧穿直接给出，也不取消 PDE 的存在性、唯一性、
正则性、弱解理论、数值稳定性和收敛性要求。严格结论是：在 \(L\)-测度表达、
路径覆盖、投影忠实、解释交换和审计证书成立的条件下，传统 PDE 的函数解任务可以被提升为
可审计、可兑现、可控风险的全局命中路径任务。
\end{remark}

\clearpage

\begin{thebibliography}{99}
\raggedright

\bibitem{RepoTex37SemanticGauge}
Gao Zheng.
\newblock \emph{G 框架正则语义规范场到自然语言语义逻辑占位规范场的高阶同伦扭曲双射、LLM 运行时降级与逻辑引擎（工作笔记：形式系统版）}.
\newblock 仓内稿件：
\code{NoteBook/notebook2/tex37_pub/}\allowbreak
\code{gframework_semantic_gauge_field_natural_language_logic_placeholder_runtime_formal_system.tex}, 2026.

\bibitem{RepoTex55PDEParadigm}
Gao Zheng.
\newblock \emph{基于高阶正交确定性收敛的统计力学同伦命中：PDE 求解范式替代（工作笔记：形式系统版）}.
\newblock 仓内稿件：
\code{NoteBook/notebook3/tex55_pub/gframework_homotopy_hitting_pde_paradigm_replacement_formal_system.tex}, 2026.

\bibitem{RepoTex53StrategyCluster}
Gao Zheng.
\newblock \emph{统计力学参数格点、同伦扭曲命中与 PDE 策略簇：从时域求解到频域簇命中的近似量子软件层算法（工作笔记：形式系统版）}.
\newblock 仓内稿件：
\code{NoteBook/notebook3/tex53_pub/}\allowbreak
\code{gframework_statistical_lattice_homotopy_tunnel_pde_strategy_cluster_formal_system.tex}, 2026.

\bibitem{RepoAppMathVol1}
Gao Zheng.
\newblock \emph{GaoZheng G-Framework and GaoZheng G-Algebra: Applied Mathematics Volume I}.
\newblock 仓内发布稿：
\code{docs/AppMathVol1/Pub_GFramework_AppMath_Vol1.tex}, 2026.

\bibitem{RepoTex56LHTS}
Gao Zheng.
\newblock \emph{LHTS：延迟信用分配的格点同伦扭曲评分与命中证书体系（工作笔记：形式系统版）}.
\newblock 仓内稿件：
\code{NoteBook/notebook3/tex56_pub/}\allowbreak
\code{gframework_lhts_delayed_credit_lattice_homotopy_twist_scoring_formal_system.tex}, 2026.

\bibitem{RepoTex57TPHLHTS}
Gao Zheng.
\newblock \emph{TPH-LHTS 轨迹函数参数同伦选择：概率路径同伦类、广义分形轨迹族与 D 结构递归决策路径（工作笔记：形式系统版）}.
\newblock 仓内稿件：
\code{NoteBook/notebook3/tex57_pub/}\allowbreak
\code{gframework_tph_lhts_trajectory_parameter_homotopy_probability_fractal_d_structure_formal_system.tex}, 2026.

\bibitem{RepoEconomicsVol}
Gao Zheng.
\newblock \emph{GaoZheng G-Framework and GaoZheng G-Algebra: Economics}.
\newblock 仓内发布稿：
\code{docs/Economics/Pub_GFramework_Economics.tex}, 2026.

\bibitem{Evans2010PDE}
Lawrence C. Evans.
\newblock \emph{Partial Differential Equations}.
\newblock American Mathematical Society, second edition, 2010.

\bibitem{BrennerScott2008FEM}
Susanne C. Brenner and L. Ridgway Scott.
\newblock \emph{The Mathematical Theory of Finite Element Methods}.
\newblock Springer, third edition, 2008.

\bibitem{Brezis2011FunctionalAnalysis}
Haim Brezis.
\newblock \emph{Functional Analysis, Sobolev Spaces and Partial Differential Equations}.
\newblock Springer, 2011.

\bibitem{Bogachev2007MeasureTheory}
Vladimir I. Bogachev.
\newblock \emph{Measure Theory}.
\newblock Springer, 2007.

\bibitem{NocedalWright2006NumericalOptimization}
Jorge Nocedal and Stephen J. Wright.
\newblock \emph{Numerical Optimization}.
\newblock Springer, second edition, 2006.

\bibitem{AmbrosioGigliSavare2008GradientFlows}
Luigi Ambrosio, Nicola Gigli, and Giuseppe Savar\'e.
\newblock \emph{Gradient Flows in Metric Spaces and in the Space of Probability Measures}.
\newblock Birkh\"auser, second edition, 2008.

\bibitem{Markowitz1952PortfolioSelection}
Harry Markowitz.
\newblock Portfolio selection.
\newblock \emph{The Journal of Finance}, 7(1):77--91, 1952.

\bibitem{RockafellarUryasev2000CVaR}
R. Tyrrell Rockafellar and Stanislav Uryasev.
\newblock Optimization of conditional value-at-risk.
\newblock \emph{The Journal of Risk}, 2(3):21--41, 2000.

\bibitem{LeVeque2007FiniteDifference}
Randall J. LeVeque.
\newblock \emph{Finite Difference Methods for Ordinary and Partial Differential Equations}.
\newblock Society for Industrial and Applied Mathematics, 2007.

\bibitem{Strikwerda2004FiniteDifference}
John C. Strikwerda.
\newblock \emph{Finite Difference Schemes and Partial Differential Equations}.
\newblock Society for Industrial and Applied Mathematics, second edition, 2004.

\bibitem{FeynmanHibbs1965PathIntegrals}
Richard P. Feynman and Albert R. Hibbs.
\newblock \emph{Quantum Mechanics and Path Integrals}.
\newblock McGraw--Hill, 1965.

\bibitem{FreidlinWentzell2012RandomPerturbations}
Mark I. Freidlin and Alexander D. Wentzell.
\newblock \emph{Random Perturbations of Dynamical Systems}.
\newblock Springer, third edition, 2012.

\end{thebibliography}

\end{CJK*}
\end{document}
